% Superposition des ondes mécaniques : interférences et ondes stationnaires — Physique Tle C — Cours signé OnBuch+
% Programme officiel MINESEC. Compiler avec : tectonic P15-superposition-ondes-mecaniques.tex
\newcommand{\DOCMATIERE}{Physique}
\newcommand{\DOCNIVEAU}{Tle C}
\newcommand{\DOCMODULE}{Module 4 : Onde, matière et transformations nucléaires}
\newcommand{\DOCLECON}{P15}
\newcommand{\DOCTITRE}{Superposition des ondes mécaniques : interférences et ondes stationnaires}
% OnBuch+ — préambule « cours signés » ENRICHI (compilé avec tectonic / XeLaTeX)
% Système visuel « Pop » d'OnBuch+ : fond crème (jamais blanc), bordures encre
% épaisses, ombres « dures » décalées, titres Archivo Black en capitales.
% Version MATHÉMATIQUES : graphes (pgfplots/tikz), boîtes pédagogiques
% (définition, propriété, méthode, attention, le savais-tu, exercice résolu,
% exercice, TP, bilan), page de garde et signature OnBuch+ ancrée en bas.
%
% Macros attendues dans main.tex AVANT \input{preamble} :
%   \DOCMATIERE  \DOCNIVEAU  \DOCMODULE  \DOCLECON (numéro)  \DOCTITRE
\documentclass[11pt]{article}

\usepackage{fontspec}
\usepackage[french]{babel}
\usepackage[a4paper,margin=2cm,top=3.4cm,bottom=2.8cm,headheight=1.4cm,footskip=1.3cm]{geometry}
\usepackage{amsmath,amssymb,amsfonts}
\usepackage{mathtools} % \xmapsto, \coloneqq... souvent utilisés par les modèles
\usepackage{cancel} % simplifications barrées (bilans de forces)
% \abs{z} (module, valeur absolue) et \norm{u} (norme), utilisés par les modèles
\providecommand{\abs}{}\let\abs\relax\DeclarePairedDelimiter{\abs}{\lvert}{\rvert}
\providecommand{\norm}{}\let\norm\relax\DeclarePairedDelimiter{\norm}{\lVert}{\rVert}
\usepackage[table]{xcolor} % [table] : fournit aussi \rowcolors (alternance de lignes)
\usepackage{pagecolor}
\usepackage{tikz}
\usetikzlibrary{arrows.meta,positioning,calc,shapes.geometric,shapes.misc,shapes.arrows,decorations.pathmorphing,decorations.markings,patterns,shadows,mindmap,trees,fit,backgrounds,matrix,intersections,angles,quotes}
\usepackage{pgfplots}
\usepackage[version=4]{mhchem} % équations nucléaires \ce{^{235}_{92}U}
\usepackage[european,siunitx]{circuitikz} % schémas électriques
\pgfplotsset{compat=1.18}
\usepgfplotslibrary{fillbetween,groupplots} % aires entre courbes (calcul intégral)
\usepackage{enumitem}
\usepackage{fancyhdr}
% [most] charge toutes les bibliothèques tcolorbox (listings, minted, raster,
% documentation...) et gonfle inutilement la mémoire TeX sur un document long
% (~300 boîtes) : on ne charge que ce qui est réellement utilisé.
\usepackage[skins,breakable]{tcolorbox}
\usepackage{graphicx}
\usepackage{colortbl}
\usepackage{booktabs}
\usepackage{tabularx}
\usepackage{diagbox} % case d'en-tête diagonale des tableaux à double entrée
% Colonnes extensibles centrée / gauche / droite (C, L, R), souvent utilisées par les modèles
\newcolumntype{C}{>{\centering\arraybackslash}X}
\newcolumntype{L}{>{\raggedright\arraybackslash}X}
\newcolumntype{R}{>{\raggedleft\arraybackslash}X}
\usepackage{array}
\usepackage{multicol}
\usepackage{siunitx}
% parse-numbers=false : accepte aussi \SI{2,5 \times 10^{-3}}{...}
\sisetup{locale=FR,output-decimal-marker={,},group-minimum-digits=4,parse-numbers=false}
\DeclareSIUnit\ohm{\ensuremath{\symup{\Omega}}} % U+2126 absent de la police
\DeclareSIUnit\division{div} % réglages d'oscilloscope (ms/div, V/div)
\usepackage{qrcode}
% \degree : commande fréquemment utilisée par les modèles pour un angle ou une
% température en dehors de \SI{}{} (non fournie par les paquets chargés).
\providecommand{\degree}{^{\circ}}
% \qed (fin de démonstration), utilisé par les modèles sans amsthm
\providecommand{\qed}{\hfill$\square$}
% \barre : double barre des valeurs interdites dans un tableau de variations (tkz-tab)
\providecommand{\barre}{\ensuremath{\|}}
% environnement proof (amsthm non chargé) utilisé par les modèles
\ifdefined\proof\else\newenvironment{proof}[1][Démonstration]{\par\noindent\textit{#1.}\ }{\qed\par}\fi
\usepackage{lastpage}
\usepackage{hyperref}
\hypersetup{hidelinks}

% ============================================================
% Palette « Pop » OnBuch+ — mêmes hex que la classe P (plus_ui.dart)
% ============================================================
\definecolor{poporange}{HTML}{F59321}
\definecolor{popdark}{HTML}{E07A0C}
\definecolor{popcream}{HTML}{FFF6E8}
\definecolor{popink}{HTML}{1C1714}
\definecolor{poppurple}{HTML}{7A5AE0}
\definecolor{popgreen}{HTML}{2E9E6A}
\definecolor{poppink}{HTML}{E0457B}
\definecolor{popblue}{HTML}{2F7DE1}
\definecolor{popgold}{HTML}{C98A12}
\definecolor{popmuted}{HTML}{8A7F76}
\definecolor{popline}{HTML}{EADFD2}
\definecolor{popgray}{HTML}{8A7F76}
\colorlet{popgrey}{popgray}
% Alias tolérants (le modèle nomme parfois les couleurs différemment)
\colorlet{popwhite}{white}
\colorlet{popblack}{popink}
\colorlet{popred}{poppink}
\colorlet{popyellow}{popgold}
% Teintes claires pour fonds de tableaux / zones
\colorlet{popblueL}{popblue!10!white}
\colorlet{poporangeL}{poporange!14!white}
\colorlet{popgreenL}{popgreen!12!white}
\colorlet{poppurpleL}{poppurple!12!white}
\colorlet{poppinkL}{poppink!12!white}
\colorlet{popgoldL}{popgold!14!white}

\pagecolor{popcream}

% ============================================================
% Typographie
% ============================================================
\defaultfontfeatures{Ligatures=TeX}
\setmainfont{PlusJakartaSans-Medium.ttf}[Path=../../../fonts/,BoldFont=PlusJakartaSans-Bold.ttf]
% Maths en sans-serif (Fira Math) avec chiffres et lettres droites de Plus
% Jakarta, pour que les formules se fondent dans le texte.
\usepackage[math-style=ISO,bold-style=ISO]{unicode-math}
\setmathfont{FiraMath-Regular.otf}[Path=../../../fonts/]
\setmathfont{PlusJakartaSans-Medium.ttf}[Path=../../../fonts/,range={up/num,up/{Latin,latin},\mathup}]
\usepackage{icomma} % virgule décimale sans espace en mode math
% Caractères mathématiques Unicode absents des polices de texte (Plus Jakarta,
% Archivo Black) : rendus par leur équivalent mathématique, en texte comme en
% formule — sinon ils s'affichent en carré vide (ex. « Arithmétique dans ℤ »).
\usepackage{newunicodechar}
\newunicodechar{ℝ}{\ensuremath{\mathbb{R}}}
\newunicodechar{ℤ}{\ensuremath{\mathbb{Z}}}
\newunicodechar{ℂ}{\ensuremath{\mathbb{C}}}
\newunicodechar{ℕ}{\ensuremath{\mathbb{N}}}
\newunicodechar{ℚ}{\ensuremath{\mathbb{Q}}}
\newunicodechar{𝔻}{\ensuremath{\mathbb{D}}}
\newunicodechar{α}{\ensuremath{\alpha}}
\newunicodechar{β}{\ensuremath{\beta}}
\newunicodechar{γ}{\ensuremath{\gamma}}
\newunicodechar{δ}{\ensuremath{\delta}}
\newunicodechar{ε}{\ensuremath{\varepsilon}}
\newunicodechar{θ}{\ensuremath{\theta}}
\newunicodechar{λ}{\ensuremath{\lambda}}
\newunicodechar{μ}{\ensuremath{\mu}}
\newunicodechar{σ}{\ensuremath{\sigma}}
\newunicodechar{τ}{\ensuremath{\tau}}
\newunicodechar{φ}{\ensuremath{\varphi}}
\newunicodechar{ψ}{\ensuremath{\psi}}
\newunicodechar{ω}{\ensuremath{\omega}}
\newunicodechar{Σ}{\ensuremath{\Sigma}}
% majuscules produites par \MakeUppercase dans les titres (α-aminés -> Α)
\newunicodechar{Α}{\ensuremath{\alpha}}
\newunicodechar{Β}{\ensuremath{\beta}}
\newunicodechar{Γ}{\ensuremath{\Gamma}}
\newunicodechar{Δ}{\ensuremath{\Delta}}
\newunicodechar{Ε}{\ensuremath{\varepsilon}}
\newunicodechar{Θ}{\ensuremath{\Theta}}
\newunicodechar{Λ}{\ensuremath{\Lambda}}
\newunicodechar{Μ}{\ensuremath{\mu}}
\newunicodechar{Τ}{\ensuremath{\tau}}
\newunicodechar{Φ}{\ensuremath{\Phi}}
\newunicodechar{Ψ}{\ensuremath{\Psi}}
\newunicodechar{Ω}{\ensuremath{\Omega}}
\newunicodechar{↦}{\ensuremath{\mapsto}}
\newunicodechar{∈}{\ensuremath{\in}}
\newunicodechar{∉}{\ensuremath{\notin}}
\newunicodechar{∧}{\ensuremath{\wedge}}
\newunicodechar{⇔}{\ensuremath{\Leftrightarrow}}
\newunicodechar{⟺}{\ensuremath{\Longleftrightarrow}}
\newunicodechar{⇒}{\ensuremath{\Rightarrow}}
\newunicodechar{⟹}{\ensuremath{\Longrightarrow}}
\newunicodechar{∘}{\ensuremath{\circ}}
\newunicodechar{∪}{\ensuremath{\cup}}
\newunicodechar{∩}{\ensuremath{\cap}}
\newunicodechar{≡}{\ensuremath{\equiv}}
\newunicodechar{⊂}{\ensuremath{\subset}}
\newunicodechar{⊥}{\ensuremath{\perp}}
\newunicodechar{∃}{\ensuremath{\exists}}
\newunicodechar{∀}{\ensuremath{\forall}}
\newunicodechar{∛}{\ensuremath{\sqrt[3]{\,}}}
\newunicodechar{∜}{\ensuremath{\sqrt[4]{\,}}}
\newunicodechar{⃗}{\ensuremath{{}^{\to}}}
\newunicodechar{ⁿ}{\ifmmode^{n}\else\textsuperscript{\ensuremath{n}}\fi}
\newunicodechar{ˣ}{\ifmmode^{x}\else\textsuperscript{\ensuremath{x}}\fi}
\newunicodechar{ᵇ}{\ifmmode^{b}\else\textsuperscript{\ensuremath{b}}\fi}
\newunicodechar{ʳ}{\ifmmode^{r}\else\textsuperscript{\ensuremath{r}}\fi}
\newunicodechar{ᵘ}{\ifmmode^{u}\else\textsuperscript{\ensuremath{u}}\fi}
\newunicodechar{ʸ}{\ifmmode^{y}\else\textsuperscript{\ensuremath{y}}\fi}
\newunicodechar{ᵃ}{\ifmmode^{a}\else\textsuperscript{\ensuremath{a}}\fi}
\newunicodechar{ᵉ}{\ifmmode^{e}\else\textsuperscript{\ensuremath{e}}\fi}
\newunicodechar{ᵏ}{\ifmmode^{k}\else\textsuperscript{\ensuremath{k}}\fi}
\newunicodechar{ᵖ}{\ifmmode^{p}\else\textsuperscript{\ensuremath{p}}\fi}
\newunicodechar{ᴺ}{\ifmmode^{N}\else\textsuperscript{\ensuremath{N}}\fi}
\newunicodechar{ᶜ}{\ifmmode^{c}\else\textsuperscript{\ensuremath{c}}\fi}
\newunicodechar{ʰ}{\ifmmode^{h}\else\textsuperscript{\ensuremath{h}}\fi}
\newunicodechar{ᵠ}{\ifmmode^{q}\else\textsuperscript{\ensuremath{q}}\fi}
\newunicodechar{ₙ}{\ifmmode_{n}\else\textsubscript{\ensuremath{n}}\fi}
\newunicodechar{ₐ}{\ifmmode_{a}\else\textsubscript{\ensuremath{a}}\fi}
\newunicodechar{ᵢ}{\ifmmode_{i}\else\textsubscript{\ensuremath{i}}\fi}
\newunicodechar{ᵣ}{\ifmmode_{r}\else\textsubscript{\ensuremath{r}}\fi}
\newunicodechar{ₖ}{\ifmmode_{k}\else\textsubscript{\ensuremath{k}}\fi}
\newunicodechar{ᵦ}{\ifmmode_{\beta}\else\textsubscript{\ensuremath{\beta}}\fi}
\newunicodechar{ₓ}{\ifmmode_{x}\else\textsubscript{\ensuremath{x}}\fi}
\newfontfamily\popdisplay{ArchivoBlack-Regular.ttf}[Path=../../../fonts/]
\newfontfamily\poplogo{SpaceGrotesk-Bold.ttf}[Path=../../../fonts/]
\newfontfamily\popsemi{PlusJakartaSans-SemiBold.ttf}[Path=../../../fonts/]
\newfontfamily\popxbold{PlusJakartaSans-ExtraBold.ttf}[Path=../../../fonts/]
\newfontfamily\popxboldls{PlusJakartaSans-ExtraBold.ttf}[Path=../../../fonts/,LetterSpace=3.0]

\setlist[enumerate]{leftmargin=1.7em,itemsep=5pt,topsep=4pt}
\setlist[itemize]{leftmargin=1.7em,itemsep=4pt,topsep=4pt,label={\color{poporange}\textbullet}}
\setlength{\parindent}{0pt}
\setlength{\parskip}{5pt}
\renewcommand{\arraystretch}{1.25}

% pgfplots : style Pop par défaut
\pgfplotsset{
  every axis/.append style={axis line style={popink,line width=1pt},tick style={popink},
    grid style={popline,line width=0.5pt},label style={font=\small},tick label style={font=\footnotesize}},
  popcurve/.style={poporange,line width=1.8pt,no markers,smooth},
}
% Même style disponible dans un \draw TikZ simple (hors pgfplots)
\tikzset{popcurve/.style={poporange,line width=1.8pt}}
\tikzset{popgrid/.style={popline,very thin}}
\colorlet{popcurve}{poporange} % le nom du style est parfois utilisé comme couleur

% ============================================================
% Pill / squiggle
% ============================================================
\newcommand{\poppill}[3]{%
  \tikz[baseline=-0.55ex]\node[draw=popink,line width=1.2pt,rounded corners=2.6pt,%
    fill=#2,inner xsep=6.5pt,inner ysep=3pt]{%
    \popxboldls\fontsize{7}{7}\selectfont\color{#3}\MakeUppercase{#1}};%
}
\newcommand{\popsquiggle}[1]{%
  \tikz[baseline=-0.4ex]{
    \draw[#1,line width=2.6pt,line cap=round]
      plot [smooth] coordinates {(0,0.09) (0.34,0.01) (0.68,0.21) (1.02,0.01) (1.36,0.10)};
  }%
}
% Mot-clé surligné (marqueur orange sous le texte)
\newcommand{\cle}[1]{{\popxbold\color{popdark}#1}}

% ============================================================
% En-tête / pied
% ============================================================
\pagestyle{fancy}
\fancyhf{}
\renewcommand{\headrulewidth}{0pt}
\renewcommand{\footrulewidth}{0pt}
\lhead{\raisebox{-0.15ex}{\poplogo\fontsize{13}{13}\selectfont\color{popink}OnBuch\color{poporange}+}}
\rhead{\poppill{\DOCMATIERE\ \textbullet\ \DOCNIVEAU\ \textbullet\ Leçon \DOCLECON}{white}{popink}}
\lfoot{\popxbold\fontsize{7.6}{7.6}\selectfont\color{popmuted}\MakeUppercase{Cours signé OnBuch+}\ \textbullet\ {\popsemi\DOCTITRE}}
\rfoot{\tikz[baseline=-0.6ex]\node[rounded corners=8.5pt,fill=popink,minimum height=17pt,minimum width=17pt,inner xsep=5pt,inner sep=0]{\popxbold\fontsize{8}{8}\selectfont\color{popcream}\thepage\,/\,\pageref*{LastPage}};}

% ============================================================
% Titres
% ============================================================
\newcommand{\chaptitre}[2]{%
  \begin{tcolorbox}[enhanced,colback=poporange,colframe=popink,boxrule=2.6pt,arc=11pt,
      drop shadow=popink,left=15pt,right=15pt,top=12pt,bottom=11pt,before skip=3pt,after skip=18pt]
    \poppill{#2}{popink}{popcream}\par\vspace{9pt}
    {\raggedright\hyphenpenalty=10000\exhyphenpenalty=10000\popdisplay\fontsize{22}{24}\selectfont\color{popcream}\MakeUppercase{#1}\par}
  \end{tcolorbox}
}
% Section numérotée : pastille orange avec le numéro + titre Archivo
\newcounter{popsec}
\newcommand{\coursec}[1]{%
  \stepcounter{popsec}\setcounter{popsub}{0}%
  \par\vspace{16pt}\noindent
  \begin{minipage}{\linewidth}
  \tikz[baseline=-0.7ex]\node[draw=popink,line width=1.6pt,fill=poporange,rounded corners=4pt,minimum size=22pt,inner sep=0]{\popdisplay\fontsize{12}{12}\selectfont\color{popcream}\thepopsec};%
  \hspace{8pt}{\popdisplay\fontsize{15}{17}\selectfont\color{popink}\MakeUppercase{#1}}\par
  \vspace{4pt}\noindent{\color{poporange}\rule{36pt}{3.6pt}}
  \end{minipage}\par\nopagebreak\vspace{8pt}
}
\newcounter{popsub}
\newcommand{\courssub}[1]{%
  \stepcounter{popsub}%
  \par\vspace{9pt}\noindent{\popxbold\fontsize{11.5}{13}\selectfont\color{popink}{\color{poporange}\thepopsec.\thepopsub}\hspace{6pt}#1}\par\nopagebreak\vspace{4pt}
}

% ============================================================
% Boîtes pédagogiques — même recette : fond blanc, bordure épaisse colorée,
% ombre dure encre, badge plein. Titre optionnel [#1].
% ============================================================
\tcbset{popbase/.style={breakable,enhanced,colback=white,boxrule=2.3pt,arc=9pt,
  shadow={3pt}{-3pt}{0pt}{popink},left=14pt,right=14pt,top=11pt,bottom=11pt,before skip=11pt,after skip=15pt,
  fontupper=\popsemi\fontsize{10.6}{14.5}\selectfont\color{popink},
  fontlower=\popsemi\fontsize{10.6}{14.5}\selectfont\color{popink}}}
% Coche dessinée (le glyphe ✓ n'existe pas dans les polices du document)
\newcommand{\popcheck}[1]{\tikz[baseline=-0.5ex]\draw[#1,line width=1.6pt,line cap=round,line join=round](-0.13,0.0)--(-0.04,-0.09)--(0.13,0.1);}
\providecommand{\coche}{\popcheck{popgreen}}
\providecommand{\resultat}[1]{\par\smallskip\noindent\fcolorbox{popgreen}{popgreenL}{\parbox{\dimexpr\linewidth-2\fboxsep-2\fboxrule\relax}{\textbf{Résultat.} #1}}\par\smallskip}
\newcommand{\popboxhead}[3]{\poppill{#1}{#2}{white}\ifx\relax#3\relax\else\hspace{6pt}{\popxbold\fontsize{10.6}{12}\selectfont\color{popink}#3}\fi\par\vspace{7pt}}

\newtcolorbox{aretenir}[1][]{popbase,colframe=popgreen,before upper={\popboxhead{À retenir}{popgreen}{#1}}}
\newtcolorbox{exemplebox}[1][]{popbase,colframe=poporange,before upper={\popboxhead{Exemple}{poporange}{#1}}}
\newtcolorbox{definition}[1][]{popbase,colframe=popblue,before upper={\popboxhead{Définition}{popblue}{#1}}}
\newtcolorbox{propriete}[1][]{popbase,colframe=poppurple,before upper={\popboxhead{Propriété}{poppurple}{#1}}}
\newtcolorbox{methode}[1][]{popbase,colframe=popgold,before upper={\popboxhead{Méthode}{popgold}{#1}}}
\newtcolorbox{attention}[1][]{popbase,colframe=poppink,before upper={\popboxhead{Attention}{poppink}{#1}}}
\newtcolorbox{experience}[1][]{popbase,colframe=popblue,colback=popblueL,before upper={\popboxhead{Expérience}{popblue}{#1}}}
% « Le savais-tu ? » : carte violette pleine (contexte, culture, Cameroun)
\newtcolorbox{savaistu}[1][]{popbase,colback=poppurple,colframe=popink,
  fontupper=\popsemi\fontsize{10.4}{14.2}\selectfont\color{white},
  before upper={\poppill{Le savais-tu\,?}{white}{poppurple}\ifx\relax#1\relax\else\hspace{6pt}{\popxbold\color{white}#1}\fi\par\vspace{7pt}}}
% Exercice résolu : énoncé en haut, \tcblower puis la solution sur fond crème
\newcounter{popexo}\newcounter{popexr}
\newtcolorbox{exoresolu}[1][]{popbase,colframe=poporange,colbacklower=popcream,
  segmentation style={popink,line width=1.2pt,dashed},
  before upper={\stepcounter{popexr}\popboxhead{Exercice résolu \thepopexr}{poporange}{#1}},
  before lower={\poppill{Solution}{popink}{popcream}\par\vspace{6pt}}}
% Exercice (non résolu en place) — la correction est dans la partie Corrigés
\newtcolorbox{exercice}[1][]{popbase,colframe=popink,
  before upper={\stepcounter{popexo}\popboxhead{Exercice \thepopexo}{popink}{#1}}}
% Corrigé
\newtcolorbox{corrige}[1][]{popbase,colframe=popgreen,colback=popgreenL,
  before upper={\popboxhead{Corrigé}{popgreen}{#1}}}
% Objectifs / prérequis / bilan
\newtcolorbox{objectifs}{popbase,colframe=popink,colback=poporangeL,
  before upper={\popboxhead{Objectifs de la leçon}{popink}{}}}
\newtcolorbox{prerequis}{popbase,colframe=popink,colback=popblueL,
  before upper={\popboxhead{Prérequis}{popblue}{}}}
\newtcolorbox{bilan}{popbase,colframe=popink,colback=white,boxrule=2.8pt,
  before upper={\popboxhead{Fiche bilan}{poporange}{L'essentiel à connaître pour l'examen}}}
% Figure encadrée avec légende
\newtcolorbox{popfigure}[1][]{enhanced,breakable=false,colback=white,colframe=popline,boxrule=1.4pt,arc=8pt,
  left=10pt,right=10pt,top=10pt,bottom=8pt,before skip=10pt,after skip=12pt,halign=center,
  lower separated=false,halign lower=center,
  fontlower=\popsemi\fontsize{9}{11.5}\selectfont\color{popmuted},
  #1}
\newcommand{\legende}[1]{\par\vspace{6pt}{\popsemi\fontsize{9}{11.5}\selectfont\color{popmuted}#1}}

% ============================================================
% Signature OnBuch+
% ============================================================
\newcommand{\poplogoinline}[1]{{\poplogo\fontsize{#1}{#1}\selectfont\color{popink}OnBuch\color{poporange}+}}
\newcommand{\signatureonbuch}{%
  \begin{tcolorbox}[enhanced,colback=popink,colframe=popink,boxrule=2.2pt,arc=10pt,
      drop shadow=poporange,left=14pt,right=14pt,top=10pt,bottom=10pt,before skip=0pt,after skip=0pt]
    \tikz[baseline=-0.7ex]\node[circle,fill=popgreen,draw=popcream,line width=1.3pt,minimum size=22pt,inner sep=0]{\popcheck{white}};%
    \hspace{9pt}\begin{minipage}[c]{0.62\linewidth}
      {\popxbold\fontsize{12}{13}\selectfont\color{popcream}Signé\ \textbullet\ {\poplogo OnBuch\color{poporange}+}}\par\vspace{2pt}
      {\raggedright\popsemi\fontsize{8}{10.5}\selectfont\color{popline}\MakeUppercase{Équipe pédagogique OnBuch+}\\\MakeUppercase{Conforme au programme officiel MINESEC}\par}
    \end{minipage}\hfill
    {\poplogo\fontsize{22}{22}\selectfont\color{popcream}OnBuch\color{poporange}+}
  \end{tcolorbox}%
}

% ============================================================
% Page de garde
% #1 titre · #2 matière · #3 niveau · #4 module · #5 n° leçon · #6 durée ·
% #7 description / objectifs (texte ou itemize)
% ============================================================
\newcommand{\pagegarde}[7]{%
  \thispagestyle{empty}
  \newgeometry{a4paper,margin=1.7cm,top=1.6cm,bottom=1.4cm}%
  \begin{tikzpicture}[remember picture,overlay]
    \fill[poporange!18] ([xshift=-3.2cm,yshift=-3.6cm]current page.north east) circle (4.6cm);
    \fill[poppurple!14] ([xshift=2.4cm,yshift=7.6cm]current page.south west) circle (3.8cm);
  \end{tikzpicture}%
  {\poplogo\fontsize{30}{30}\selectfont\color{popink}OnBuch\color{poporange}+}\hfill
  \raisebox{0.6ex}{\poppill{Cours signé \textbullet\ Premium}{popink}{popcream}}\par
  \vspace{1pt}\popsquiggle{poppurple}\par
  \vspace{8pt}
  \begin{tcolorbox}[enhanced,colback=poporange,colframe=popink,boxrule=2.8pt,arc=13pt,
      drop shadow=popink,left=18pt,right=18pt,top=12pt,bottom=13pt,after skip=10pt]
    \poppill{#2 \textbullet\ #3}{popink}{popcream}\hspace{5pt}\poppill{#4}{white}{popink}\par\vspace{8pt}
    {\popdisplay\fontsize{14}{14}\selectfont\color{popink}LEÇON #5}\par\vspace{4pt}
    {\raggedright\hyphenpenalty=10000\exhyphenpenalty=10000\popdisplay\fontsize{25}{27}\selectfont\color{popcream}\MakeUppercase{#1}\par}
    \vspace{8pt}
    {\popxbold\fontsize{9.5}{9.5}\selectfont\color{popink}\MakeUppercase{Durée conseillée : #6}}
  \end{tcolorbox}
  \begin{tcolorbox}[enhanced,colback=white,colframe=popink,boxrule=2.2pt,arc=10pt,
      drop shadow=popink,left=16pt,right=16pt,top=10pt,bottom=10pt,after skip=8pt]
    \poppill{Dans cette leçon}{poppurple}{white}\par\vspace{5pt}
    \popsemi\fontsize{9.3}{12.6}\selectfont\color{popink}#7
  \end{tcolorbox}
  \noindent\begin{minipage}[t]{0.487\linewidth}
    \begin{tcolorbox}[enhanced,colback=popgreen,colframe=popink,boxrule=2.2pt,arc=10pt,
        drop shadow=popink,left=11pt,right=11pt,top=10pt,bottom=10pt,height=3.1cm,valign=center]
      \tcbox[colback=white,colframe=popink,boxrule=1.3pt,arc=5pt,boxsep=0pt,nobeforeafter,
        left=3pt,right=3pt,top=3pt,bottom=3pt,valign=center]{\qrcode[height=1.9cm]{https://play.google.com/store/apps/details?id=com.luvvix.onbuch&hl=fr}}%
      \hspace{8pt}\begin{minipage}[c]{0.5\linewidth}
        {\popdisplay\fontsize{11}{12.5}\selectfont\color{white}\MakeUppercase{Télécharge OnBuch}}\par\vspace{4pt}
        {\raggedright\popsemi\fontsize{8.4}{11}\selectfont\color{white}Gratuit \textbullet\ Google Play\\Tuteur IA, annales, concours, résultats\par}
      \end{minipage}
    \end{tcolorbox}
  \end{minipage}\hfill
  \begin{minipage}[t]{0.487\linewidth}
    \begin{tcolorbox}[enhanced,colback=poppurple,colframe=popink,boxrule=2.2pt,arc=10pt,
        drop shadow=popink,left=11pt,right=11pt,top=10pt,bottom=10pt,height=3.1cm,valign=center]
      \tikz[baseline=-0.7ex]\node[circle,fill=white,draw=popink,line width=1.3pt,minimum size=1.6cm,inner sep=0]{\popdisplay\fontsize{24}{24}\selectfont\color{poppurple}+};%
      \hspace{8pt}\begin{minipage}[c]{0.6\linewidth}
        {\popdisplay\fontsize{11}{12.5}\selectfont\color{white}\MakeUppercase{Passe à OnBuch+}}\par\vspace{4pt}
        {\raggedright\popsemi\fontsize{8.4}{11}\selectfont\color{white}Tous les cours signés, Tuteur IA illimité, fiches PDF — onglet OnBuch+ de l'app\par}
      \end{minipage}
    \end{tcolorbox}
  \end{minipage}
  \vspace{8pt}
  \popcontenu
  \vfill
  \signatureonbuch
  \restoregeometry
  \newpage
}

% Rangée « Ce cours contient » (page de garde)
\newcommand{\poptile}[3]{% #1 couleur #2 gros texte #3 légende
  \begin{tcolorbox}[enhanced,colback=#1,colframe=popink,boxrule=2pt,arc=9pt,shadow={2.5pt}{-2.5pt}{0pt}{popink},
      left=6pt,right=6pt,top=9pt,bottom=9pt,halign=center,valign=center,height=1.75cm,width=\linewidth,nobeforeafter]
    {\popdisplay\fontsize{12.5}{13}\selectfont\color{white}\MakeUppercase{#2}}\par\vspace{3pt}
    {\centering\popsemi\fontsize{7.8}{9.5}\selectfont\color{white}#3\par}
  \end{tcolorbox}}
\newcommand{\popcontenu}{%
  \par\noindent{\popxboldls\fontsize{7.5}{7.5}\selectfont\color{popmuted}CE COURS SIGNÉ CONTIENT}\par\vspace{7pt}
  \noindent\begin{minipage}[t]{0.235\linewidth}\poptile{popblue}{Cours}{complet, illustré, pas à pas}\end{minipage}\hfill
  \begin{minipage}[t]{0.235\linewidth}\poptile{poporange}{Méthodes}{et exercices résolus}\end{minipage}\hfill
  \begin{minipage}[t]{0.235\linewidth}\poptile{poppink}{Exercices}{type examen + corrigés}\end{minipage}\hfill
  \begin{minipage}[t]{0.235\linewidth}\poptile{popgreen}{Fiche bilan}{l'essentiel à retenir}\end{minipage}\par}

% Sommaire
\newtcolorbox{sommaire}{enhanced,colback=white,colframe=popink,boxrule=2.2pt,arc=10pt,
  drop shadow=popink,left=18pt,right=18pt,top=13pt,bottom=13pt,before skip=6pt,after skip=20pt,
  before upper={\poppill{Sommaire}{poppurple}{white}\par\vspace{9pt}\popsemi\fontsize{10.6}{14.5}\selectfont\color{popink}}}

% Signature de fin de cours — poussée tout en bas de la dernière page
\newcommand{\findecours}{%
  \par\vfill\nopagebreak
  \begin{center}\popsquiggle{poporange}\end{center}\vspace{4pt}
  \signatureonbuch
}
\AtBeginDocument{\providecommand{\llbracket}{\lBrack}\providecommand{\rrbracket}{\rBrack}} % intervalles d'entiers

%% FIGFIX-BEGIN v4 — correctifs globaux des figures (docs/onbuch-generation/tools/figfix)
\usepackage{adjustbox}\usepackage{etoolbox}
% toute figure TikZ/pgfplots plus large que la ligne est réduite à la largeur disponible
\BeforeBeginEnvironment{tikzpicture}{\begin{adjustbox}{max width=\linewidth}}
\AfterEndEnvironment{tikzpicture}{\end{adjustbox}}
% légendes pgfplots lisibles par-dessus les courbes
\pgfplotsset{every axis/.append style={legend style={fill=white,fill opacity=0.92,text opacity=1,draw=black!15,inner sep=3pt}}}
% étiquettes d'arêtes (syntaxe edge["..."]) avec fond blanc
\tikzset{every edge quotes/.append style={fill=white,inner sep=1.2pt}}
%% FIGFIX-END
\begin{document}

\pagegarde{Superposition des ondes mécaniques : interférences et ondes stationnaires}{Physique}{Tle C}{Module 4 : Onde, matière et transformations nucléaires}{P15}{6 h}{Cette leçon étudie la superposition des ondes mécaniques : les interférences à la surface de l'eau entre deux sources cohérentes et les ondes stationnaires sur une corde vibrante (expérience de Melde). Elle se conclut par l'étude de systèmes technologiques à ondes : échographie, télémétrie, sonar et radar.
\vspace{4pt}
\begin{multicols}{2}\small
\begin{itemize}
\item À la fin de cette leçon, je sais énoncer et appliquer le principe de superposition des ondes mécaniques
\item À la fin de cette leçon, je sais décrire et schématiser l'expérience d'interférences à la surface de l'eau avec deux sources cohérentes
\item À la fin de cette leçon, je sais déterminer la différence de marche en un point et caractériser les points de vibration maximale et les points de repos
\item À la fin de cette leçon, je sais calculer l'interfrange et exploiter une figure d'interférences observée en éclairage normal ou stroboscopique
\item À la fin de cette leçon, je sais décrire l'expérience de Melde et interpréter la formation d'une onde stationnaire
\item À la fin de cette leçon, je sais exploiter les conditions aux limites (nœuds, ventres) pour déterminer les modes propres d'une corde
\end{itemize}\end{multicols}}

\begin{sommaire}
\begin{multicols}{2}
\begin{enumerate}
\item Principe de superposition et sources cohérentes
\item Interférences à la surface de l'eau
\item Interfrange et exploitation de la figure d'interférences
\item Ondes stationnaires : expérience de Melde
\item Modes propres d'une corde et instruments de musique
\item Systèmes technologiques à ondes
\item Activité d'intégration
\item Méthodes et exercices résolus
\item Exercices
\item Corrigés des exercices
\item Fiche bilan
\end{enumerate}
\end{multicols}
\end{sommaire}

\begin{objectifs}
\begin{itemize}
\item À la fin de cette leçon, je sais énoncer et appliquer le principe de superposition des ondes mécaniques
\item À la fin de cette leçon, je sais décrire et schématiser l'expérience d'interférences à la surface de l'eau avec deux sources cohérentes
\item À la fin de cette leçon, je sais déterminer la différence de marche en un point et caractériser les points de vibration maximale et les points de repos
\item À la fin de cette leçon, je sais calculer l'interfrange et exploiter une figure d'interférences observée en éclairage normal ou stroboscopique
\item À la fin de cette leçon, je sais décrire l'expérience de Melde et interpréter la formation d'une onde stationnaire
\item À la fin de cette leçon, je sais exploiter les conditions aux limites (nœuds, ventres) pour déterminer les modes propres d'une corde
\item À la fin de cette leçon, je sais relier les modes propres d'une corde au fonctionnement des instruments de musique à cordes
\item À la fin de cette leçon, je sais expliquer le principe de fonctionnement de l'échographie, de la télémétrie, du sonar et du radar
\end{itemize}
\end{objectifs}

\begin{prerequis}
\begin{itemize}
\item Ondes mécaniques progressives : définition, célérité, longueur d'onde, période (classe de Première)
\item Relation fondamentale λ = cT = c/f
\item Retard d'une onde en un point et déphasage entre deux points
\item Célérité d'une onde sur une corde tendue : c = √(F/μ) (admissible, rappel)
\item Notions de géométrie : distance entre deux points, hyperbole comme ensemble de points
\end{itemize}
\end{prerequis}

\newpage

\coursec{Principe de superposition et sources cohérentes}

Au bord du fleuve Wouri, deux pirogues qui tanguent côte à côte créent des vagues qui se croisent : à certains endroits l'eau danse violemment, à d'autres elle reste étrangement calme. De même, quand le joueur de mvet pince une corde tendue entre deux points fixes, seuls certains sons précis sortent de l'instrument. Pourquoi deux ondes qui se rencontrent peuvent-elles se renforcer ou s'annuler ? Pourquoi une corde vibrante ne produit-elle que certaines fréquences ? Cette leçon répond à ces questions et explique le fonctionnement de l'échographie pratiquée à l'hôpital ou du sonar des pêcheurs.

Toute cette leçon repose sur une idée simple et puissante : lorsque plusieurs ondes traversent le même milieu, leurs effets \emph{s'ajoutent}. C'est le \cle{principe de superposition}, que nous allons énoncer rigoureusement après quelques rappels.

\courssub{Rappels : l'onde mécanique progressive}

\begin{definition}[Onde mécanique progressive]
Une \cle{onde mécanique progressive} est le phénomène de propagation d'une perturbation dans un milieu matériel (corde, eau, air), \textbf{sans transport de matière} mais avec \textbf{transport d'énergie}. La perturbation se propage à une vitesse appelée \cle{célérité} $c$ (en $\mathrm{m\cdot s^{-1}}$), qui dépend des propriétés du milieu (tension et masse linéique pour une corde, rigidité et inertie en général).
\end{definition}

Pour une corde de masse linéique $\mu$ tendue par une force de norme $T$, la célérité des ondes transversales est :
\[
c = \sqrt{\dfrac{T}{\mu}}
\]
Vérifions l'homogénéité : $[T] = \mathrm{kg\cdot m\cdot s^{-2}}$ et $[\mu] = \mathrm{kg\cdot m^{-1}}$, donc $[T/\mu] = \mathrm{m^2\cdot s^{-2}}$, et la racine carrée donne bien des $\mathrm{m\cdot s^{-1}}$. La formule est dimensionnellement correcte.

\begin{aretenir}[Grandeurs caractéristiques d'une onde sinusoïdale]
Une onde progressive sinusoïdale de fréquence $f$ (en Hz), de période $T = 1/f$ (en s), possède une \cle{longueur d'onde} :
\[
\lambda = c\,T = \dfrac{c}{f}
\]
$\lambda$ est la plus petite distance séparant deux points du milieu qui vibrent \emph{en phase}. L'\cle{élongation} $y(M,t)$ d'un point $M$ est son déplacement par rapport à sa position d'équilibre, mesuré perpendiculairement à la direction de propagation (onde transversale) ou le long de celle-ci (onde longitudinale).
\end{aretenir}

\begin{attention}[Sans transport de matière !]
Quand la vague passe sous la pirogue sur le Wouri, la pirogue monte et descend, mais elle n'est \textbf{pas emportée} par la vague : l'eau oscille sur place. C'est la \emph{perturbation} qui voyage, pas le milieu. Confondre vitesse de l'onde $c$ et vitesse des points du milieu est l'erreur la plus classique du chapitre.
\end{attention}

\courssub{Observation : croisement de deux impulsions sur une corde}

\begin{experience}[Croisement de deux impulsions]
\textbf{Matériel :} une longue corde (ou un ressort souple) tendue horizontalement entre deux opérateurs.

\textbf{Protocole :} chaque opérateur donne une secousse brève vers le haut à son extrémité. Deux impulsions naissent et se propagent l'une vers l'autre.

\textbf{Observations :}
\begin{enumerate}
\item Avant le croisement, chaque impulsion se propage sans déformation notable.
\item Au moment du croisement, la corde prend une forme dont la hauteur, en chaque point, est la \emph{somme} des hauteurs des deux impulsions : si les deux impulsions sont vers le haut, on observe un renforcement (bosse deux fois plus haute) ; si l'une est vers le haut et l'autre vers le bas, elles peuvent s'annuler momentanément.
\item Après le croisement, chaque impulsion reprend sa forme initiale et poursuit sa route \emph{comme si de rien n'était}.
\end{enumerate}
\end{experience}

\begin{popfigure}
\begin{center}
\begin{tikzpicture}[scale=1.0, >=Stealth]
% --- Avant ---
\draw[popmuted, dashed] (0,0) -- (5.2,0);
\draw[popblue, very thick] (0.4,0) .. controls (0.8,0) and (1.0,0.9) .. (1.3,0.9) .. controls (1.6,0.9) and (1.8,0) .. (2.2,0);
\draw[poporange, very thick] (3.0,0) .. controls (3.4,0) and (3.6,0.9) .. (3.9,0.9) .. controls (4.2,0.9) and (4.4,0) .. (4.8,0);
\draw[->, popblue] (1.3,1.15) -- (2.1,1.15);
\draw[->, poporange] (3.9,1.15) -- (3.1,1.15);
\node[popdark] at (2.6,-0.55) {\small \textbf{a)} Avant le croisement};
% --- Pendant ---
\begin{scope}[yshift=-2.6cm]
\draw[popmuted, dashed] (0,0) -- (5.2,0);
\draw[poppurple, very thick] (1.6,0) .. controls (2.0,0) and (2.2,1.8) .. (2.6,1.8) .. controls (3.0,1.8) and (3.2,0) .. (3.6,0);
\node[popdark] at (2.6,-0.55) {\small \textbf{b)} Pendant : superposition (renforcement)};
\node[poppurple] at (2.6,2.1) {\small $y = y_1 + y_2$};
\end{scope}
% --- Après ---
\begin{scope}[yshift=-5.2cm]
\draw[popmuted, dashed] (0,0) -- (5.2,0);
\draw[poporange, very thick] (0.4,0) .. controls (0.8,0) and (1.0,0.9) .. (1.3,0.9) .. controls (1.6,0.9) and (1.8,0) .. (2.2,0);
\draw[popblue, very thick] (3.0,0) .. controls (3.4,0) and (3.6,0.9) .. (3.9,0.9) .. controls (4.2,0.9) and (4.4,0) .. (4.8,0);
\draw[->, poporange] (1.3,1.15) -- (0.5,1.15);
\draw[->, popblue] (3.9,1.15) -- (4.7,1.15);
\node[popdark] at (2.6,-0.55) {\small \textbf{c)} Après : chaque impulsion a poursuivi sa route inchangée};
\end{scope}
\end{tikzpicture}
\end{center}
\legende{Croisement de deux impulsions sur une corde tendue. Au point de rencontre, les élongations s'ajoutent algébriquement ; après le croisement, chaque impulsion est intacte : les ondes sont indépendantes.}
\end{popfigure}

\courssub{Énoncé du principe de superposition}

L'expérience précédente se généralise à toutes les ondes mécaniques de faible amplitude.

\begin{definition}[Principe de superposition]
Lorsque deux ondes se propagent dans un même milieu et se croisent en un point $M$, l'élongation résultante du point $M$ est la \textbf{somme algébrique} des élongations que chaque onde produirait seule en ce point :
\[
y(M,t) = y_1(M,t) + y_2(M,t)
\]
Après le croisement, chaque onde poursuit sa propagation \textbf{sans être modifiée} : c'est l'\cle{indépendance des ondes}.
\end{definition}

\begin{attention}[Somme \emph{algébrique}]
Les élongations sont des grandeurs \emph{signées} : une élongation vers le haut ($y_1 > 0$) et une élongation vers le bas ($y_2 < 0$) peuvent se compenser. Si $y_1 = +3\ \mathrm{mm}$ et $y_2 = -3\ \mathrm{mm}$ en un point à un instant donné, alors $y = 0$ : la corde passe momentanément par sa position d'équilibre en ce point, sans que les ondes aient disparu !
\end{attention}

\begin{savaistu}[Une loi des petites perturbations]
Le principe de superposition n'est exact que pour des perturbations de \textbf{faible amplitude}, c'est-à-dire lorsque le milieu répond de façon \emph{linéaire} (la déformation est proportionnelle à la cause). Une grosse vague déferlante sur la plage de Kribi, un tsunami, ou une onde de choc ne respectent plus ce principe : les équations deviennent non linéaires et les ondes se déforment en se propageant. En Terminale, toutes les situations étudiées sont dans le régime linéaire.
\end{savaistu}

\begin{exemplebox}[Application numérique du principe de superposition]
Deux impulsions se croisent en un point $M$ d'une corde. À l'instant $t_0$, la première seule donnerait $y_1(M,t_0) = +4{,}0\ \mathrm{mm}$ et la seconde $y_2(M,t_0) = +2{,}5\ \mathrm{mm}$. L'élongation résultante est :
\[
y(M,t_0) = 4{,}0 + 2{,}5 = 6{,}5\ \mathrm{mm}
\]
Si la seconde impulsion était inversée ($y_2 = -2{,}5\ \mathrm{mm}$), on aurait $y = 4{,}0 - 2{,}5 = 1{,}5\ \mathrm{mm}$ : renforcement partiel vers le haut.
\end{exemplebox}

\courssub{Sources cohérentes et sources synchrones}

Passons maintenant des impulsions brèves aux \emph{ondes sinusoïdales entretenues}, émises par deux sources $S_1$ et $S_2$ vibrant en permanence — par exemple deux pointes qui frappent régulièrement la surface de l'eau dans une cuve à ondes.

\begin{popfigure}
\begin{center}
\begin{tikzpicture}[scale=0.85]
% Source S1
\foreach \r in {0.5,1.0,1.5,2.0,2.5,3.0,3.5}
  \draw[popblue!60] (1.5,0) circle (\r);
% Source S2
\foreach \r in {0.5,1.0,1.5,2.0,2.5,3.0,3.5}
  \draw[poporange!60] (5.5,0) circle (\r);
\fill[popblue] (1.5,0) circle (0.09);
\fill[poporange] (5.5,0) circle (0.09);
\node[below, popblue] at (1.5,-0.12) {$S_1$};
\node[below, poporange] at (5.5,-0.12) {$S_2$};
\draw[<->, >=Stealth, popdark] (1.5,3.9) -- (5.5,3.9) node[midway, above] {$d = S_1S_2$};
\fill[poppurple] (3.5,1.9) circle (0.09);
\node[above right, poppurple] at (3.5,1.9) {$M$};
\draw[poppurple, dashed] (1.5,0) -- (3.5,1.9) node[midway, sloped, above] {\scriptsize $d_1$};
\draw[poppurple, dashed] (5.5,0) -- (3.5,1.9) node[midway, sloped, above] {\scriptsize $d_2$};
\end{tikzpicture}
\end{center}
\legende{Deux sources ponctuelles $S_1$ et $S_2$ émettent des ondes circulaires de même fréquence à la surface de l'eau. En tout point $M$, les deux ondes se superposent ; l'état vibratoire de $M$ dépend des distances $d_1 = S_1M$ et $d_2 = S_2M$.}
\end{popfigure}

Imaginons d'abord deux sources qui vibrent \emph{n'importe comment} : l'une à $50\ \mathrm{Hz}$, l'autre à $60\ \mathrm{Hz}$, ou deux sources dont la phase varie au hasard. En un point $M$, le déphasage entre les deux ondes changerait sans cesse : tantôt les ondes s'ajouteraient, tantôt elles se retrancheraient, et l'œil (ou tout capteur) ne verrait qu'une agitation moyenne uniforme. Aucune structure stable n'apparaît.

\begin{definition}[Sources cohérentes, sources synchrones]
Deux sources $S_1$ et $S_2$ sont dites \cle{cohérentes} si :
\begin{itemize}
\item elles ont la \textbf{même fréquence} $f$ ;
\item leur \textbf{différence de phase reste constante} au cours du temps.
\end{itemize}
Si de plus elles vibrent \textbf{en phase} (même élongation au même instant : $y_{S_1}(t) = y_{S_2}(t)$), on dit qu'elles sont \cle{synchrones}. Les sources synchrones sont un cas particulier de sources cohérentes (différence de phase nulle).
\end{definition}

\begin{propriete}[Condition nécessaire d'interférences]
Seules des sources \textbf{cohérentes} produisent une \cle{figure d'interférences stable} : des régions où l'amplitude de vibration est maximale en permanence et des régions où elle est minimale (voire nulle) en permanence. En pratique, on réalise deux sources cohérentes en alimentant deux pointes par \emph{le même} vibreur (fourche à deux branches dans la cuve à ondes) : elles sont alors automatiquement synchrones.
\end{propriete}

\begin{attention}[Erreur fréquente au Bac]
Deux sources de \textbf{fréquences différentes} ne produisent \textbf{jamais} de figure d'interférences stable, même si elles sont très proches. Dire « il y a interférences dès que deux ondes se croisent » est faux : il y a bien superposition (elle a toujours lieu), mais le phénomène d'\emph{interférences} — des zones fixes de renforcement et d'annulation — exige la cohérence. De même, deux ampoules ordinaires indépendantes ne sont pas cohérentes entre elles ; en mécanique, on fabrique la cohérence avec un vibreur unique.
\end{attention}

\begin{exoresolu}[Superposition de deux signaux sinusoïdaux en un point]
En un point $M$ de la surface de l'eau arrivent deux ondes issues de sources synchrones de fréquence $f = 25\ \mathrm{Hz}$. À un instant $t$ donné, les élongations en $M$ sont :
\[
y_1(M,t) = 3{,}0\cos\!\left(50\pi t\right)\ \mathrm{mm}
\qquad
y_2(M,t) = 3{,}0\cos\!\left(50\pi t - \pi\right)\ \mathrm{mm}
\]
\begin{enumerate}
\item Les sources sont-elles cohérentes ? Justifier.
\item Déterminer l'élongation résultante $y(M,t)$ et commenter.
\end{enumerate}
\tcblower
\textbf{Solution.}
\begin{enumerate}
\item Les deux ondes ont la même fréquence $f = 25\ \mathrm{Hz}$ (pulsation $\omega = 2\pi f = 50\pi\ \mathrm{rad\cdot s^{-1}}$) et leur différence de phase en $M$ vaut $\Delta\varphi = \pi$, constante au cours du temps : les sources sont bien \textbf{cohérentes}.
\item Appliquons le principe de superposition :
\begin{align*}
y(M,t) &= y_1(M,t) + y_2(M,t)\\
&= 3{,}0\cos(50\pi t) + 3{,}0\cos(50\pi t - \pi)\\
&= 3{,}0\cos(50\pi t) - 3{,}0\cos(50\pi t) = 0
\end{align*}
car $\cos(\theta - \pi) = -\cos\theta$. Le point $M$ reste \textbf{immobile à tout instant} : les deux ondes y arrivent en \emph{opposition de phase} et s'annulent. C'est un point de vibration minimale (amplitude nulle) : l'eau y reste « étrangement calme », comme sur le Wouri. Nous verrons dans la section suivante que cela se produit lorsque la différence de marche $d_2 - d_1$ est un multiple impair de $\lambda/2$.
\end{enumerate}
\end{exoresolu}

\begin{aretenir}[L'essentiel de la section]
\begin{itemize}
\item \textbf{Principe de superposition :} en tout point $M$, $y(M,t) = y_1(M,t) + y_2(M,t)$ (somme algébrique) ; après le croisement, les ondes repartent inchangées.
\item Ce principe n'est valable que pour les \textbf{petites perturbations} (milieu linéaire).
\item Deux sources sont \textbf{cohérentes} si elles ont la même fréquence et une différence de phase constante ; elles sont \textbf{synchrones} si elles vibrent en phase.
\item La cohérence est la \textbf{condition nécessaire} pour obtenir une figure d'interférences stable.
\end{itemize}
\end{aretenir}

\begin{exercice}[Pour s'entraîner]
Deux pointes $S_1$ et $S_2$ sont fixées aux deux branches d'une même fourche solidaire d'un vibreur de fréquence $f = 40\ \mathrm{Hz}$ et affleurent la surface de l'eau d'une cuve à ondes.
\begin{enumerate}
\item Justifier que $S_1$ et $S_2$ sont des sources synchrones.
\item En un point $M$, les élongations sont $y_1 = 2{,}0\cos(80\pi t)$ mm et $y_2 = 2{,}0\cos(80\pi t)$ mm. Calculer l'amplitude résultante en $M$ et qualifier ce point.
\item Même question si $y_2 = 2{,}0\cos(80\pi t + \pi/2)$ mm : l'amplitude résultante vaut-elle $4{,}0$ mm ? (On rappelle que l'amplitude de la somme de deux signaux de même pulsation et de déphasage $\varphi$ est $A = \sqrt{A_1^2 + A_2^2 + 2A_1A_2\cos\varphi}$.)
\end{enumerate}
\end{exercice}

\begin{corrige}[Exercice : Pour s'entraîner]
\begin{enumerate}
\item Les deux pointes sont solidaires du \emph{même} vibreur : elles ont donc la même fréquence $f = 40\ \mathrm{Hz}$ et vibrent en phase à tout instant. Ce sont des sources \textbf{synchrones}, donc cohérentes.
\item Les deux signaux sont en phase ($\varphi = 0$) : $y = 4{,}0\cos(80\pi t)$ mm. L'amplitude résultante est $A = 4{,}0\ \mathrm{mm}$ : c'est un point de vibration \textbf{maximale}.
\item Avec $\varphi = \pi/2$, $\cos\varphi = 0$, donc $A = \sqrt{2{,}0^2 + 2{,}0^2} = \sqrt{8{,}0} \approx 2{,}8\ \mathrm{mm} \neq 4{,}0$ mm. Piège classique : on ne peut additionner directement les amplitudes que si les signaux sont \emph{en phase}.
\end{enumerate}
\end{corrige}

\coursec{Interférences à la surface de l'eau}

Dans la section précédente, nous avons établi le \cle{principe de superposition} : lorsque deux ondes se croisent en un point, l'élongation résultante est la somme algébrique des élongations de chaque onde. Nous avons aussi dégagé la notion de \cle{sources cohérentes} : deux sources de même fréquence dont le déphasage initial reste constant au cours du temps. Une question se pose alors naturellement : que voit-on \emph{concrètement} lorsque deux sources cohérentes émettent simultanément dans le même milieu ? La réponse est spectaculaire : il apparaît une \emph{figure d'interférences}, c'est-à-dire une répartition stable dans l'espace de zones où l'eau vibre fortement (amplitude maximale) et de zones où elle reste presque immobile (amplitude minimale). C'est cette figure que nous allons décrire, observer et expliquer quantitativement dans cette section.

\courssub{Description de l'expérience : la cuve à ondes}

\begin{experience}[Interférences à la surface de l'eau]
\textbf{Matériel :} une cuve à ondes (bac peu profond à fond transparent rempli d'eau), un vibreur muni de \textbf{deux pointes} $S_1$ et $S_2$ solidaires du même support, un stroboscope, un écran placé sous la cuve recevant l'image de la surface de l'eau grâce à un éclairage et un miroir (ou un vidéo-projecteur).

\textbf{Protocole :}
\begin{enumerate}
\item Régler le vibreur sur une fréquence $f$ connue (par exemple $f = \SI{20}{Hz}$).
\item Affleurer les deux pointes $S_1$ et $S_2$ à la surface de l'eau : chacune engendre des ondes circulaires transversales qui se propagent à la célérité $c$.
\item Observer la surface de l'eau (ou son image sur l'écran) en éclairage normal, puis en éclairage stroboscopique.
\end{enumerate}

\textbf{Observations :} on voit apparaître une alternance régulière de \emph{zones agitées} (l'eau y vibre avec une grande amplitude) et de \emph{zones calmes} (l'eau y reste pratiquement immobile), dessinant des lignes courbes stables autour des deux sources. Cette répartition stable constitue la \cle{figure d'interférences}.

\textbf{Interprétation :} en chaque point $M$ de la surface, deux ondes se superposent. Selon le point considéré, elles arrivent \emph{en phase} (leurs effets s'ajoutent : zone agitée) ou \emph{en opposition de phase} (leurs effets s'annulent : zone calme).
\end{experience}

\begin{attention}[Pourquoi deux pointes sur le même vibreur ?]
Pour obtenir une figure d'interférences \emph{stable}, les deux sources doivent être \cle{cohérentes} : même fréquence et déphasage constant. En fixant les deux pointes $S_1$ et $S_2$ sur le \emph{même} vibreur, on garantit automatiquement cette cohérence : les deux pointes vibrent ensemble, en phase, à la fréquence $f$ imposée. Deux vibreurs indépendants, même réglés sur la même fréquence nominale, dériveraient l'un par rapport à l'autre (instabilité thermique, électronique) et la figure bougerait sans cesse, la rendant invisible à l'œil nu.
\end{attention}

\begin{popfigure}
\begin{center}
\begin{tikzpicture}[scale=0.9]
% Cuve
\draw[popblue, thick, fill=popblueL] (-5,-2.6) rectangle (5,2.6);
\node[popblue, below] at (0,-2.6) {\small Cuve à ondes (surface de l'eau vue de dessus)};
% Vibreur
\draw[popdark, thick, fill=popmuted] (-0.9,3.4) rectangle (0.9,3.0);
\node[popdark, above] at (0,3.4) {\small Vibreur (fréquence $f$)};
\draw[popdark, thick] (-0.45,3.0) -- (-0.45,0.55);
\draw[popdark, thick] (0.45,3.0) -- (0.45,0.55);
% Sources
\fill[popink] (-0.45,0.4) circle (0.08) node[below left=1pt, popink] {\small $S_1$};
\fill[popink] (0.45,0.4) circle (0.08) node[below right=1pt, popink] {\small $S_2$};
% Ondes circulaires S1
\foreach \r in {0.5,1.0,1.5,2.0,2.5}{
  \draw[poporange, opacity=0.75] (-0.45,0.4) circle (\r);
}
% Ondes circulaires S2
\foreach \r in {0.5,1.0,1.5,2.0,2.5}{
  \draw[popgreen, opacity=0.75] (0.45,0.4) circle (\r);
}
\node[poporange] at (-3.6,1.9) {\small crêtes issues de $S_1$};
\node[popgreen] at (3.6,1.9) {\small crêtes issues de $S_2$};
\end{tikzpicture}
\end{center}
\legende{Dispositif expérimental : deux pointes $S_1$ et $S_2$ solidaires du même vibreur engendrent des ondes circulaires cohérentes qui se superposent à la surface de l'eau.}
\end{popfigure}

\courssub{Différence de marche en un point}

Considérons un point $M$ quelconque de la surface de l'eau. L'onde issue de $S_1$ parcourt la distance $d_1 = S_1M$ pour atteindre $M$ ; l'onde issue de $S_2$ parcourt la distance $d_2 = S_2M$. Comme les deux sources vibrent en phase (déphasage initial nul), tout l'écart de comportement en $M$ provient de la différence des chemins parcourus.

\begin{definition}[Différence de marche]
On appelle \cle{différence de marche} en un point $M$ la différence des distances séparant $M$ des deux sources :
\[ \delta = d_2 - d_1 = S_2M - S_1M \]
Elle représente le \emph{chemin supplémentaire} parcouru par l'onde issue de $S_2$ par rapport à celle issue de $S_1$ pour arriver en $M$. Elle s'exprime en mètres (m). Lorsque seul l'état vibratoire (maximal ou minimal) nous intéresse, on utilise sa valeur absolue $|\delta| = |S_2M - S_1M|$.
\end{definition}

\begin{popfigure}
\begin{center}
\begin{tikzpicture}[scale=1.0]
\fill[popink] (0,0) circle (0.07) node[below left] {$S_1$};
\fill[popink] (2.4,0) circle (0.07) node[below right] {$S_2$};
\fill[poporange] (4.6,2.6) circle (0.07) node[above] {$M$};
\draw[popblue, thick, ->] (0,0) -- (4.6,2.6) node[midway, above left, popblue] {$d_1 = S_1M$};
\draw[popgreen, thick, ->] (2.4,0) -- (4.6,2.6) node[midway, above right, popgreen] {$d_2 = S_2M$};
\draw[popdark, dashed] (0,0) -- (2,4,0) node[midway, below] {$a = S_1S_2$};
\node[popdark, align=center] at (2.3,-1.3) {$\delta = d_2 - d_1$ : chemin supplémentaire\\ parcouru par l'onde issue de $S_2$};
\end{tikzpicture}
\end{center}
\legende{Géométrie de la différence de marche : en $M$, l'onde issue de $S_2$ a parcouru $\delta = S_2M - S_1M$ de plus que celle issue de $S_1$.}
\end{popfigure}

\courssub{Du déphasage aux conditions d'interférences}

Établissons maintenant rigoureusement les conditions d'amplitude maximale et minimale. Les deux sources vibrent en phase avec la période $T = 1/f$. L'onde issue de $S_1$ met la durée $t_1 = d_1/c$ pour atteindre $M$ ; celle issue de $S_2$ met $t_2 = d_2/c$. L'onde issue de $S_2$ arrive donc en $M$ avec le \emph{retard} :
\[ \tau = t_2 - t_1 = \frac{d_2 - d_1}{c} = \frac{\delta}{c} \]
À ce retard temporel correspond un \cle{déphasage} (on rappelle qu'un retard d'une période $T$ correspond à un déphasage de $2\pi$) :
\[ \Delta\varphi = 2\pi\,\frac{\tau}{T} = \frac{2\pi\,\delta}{cT} = \frac{2\pi\,\delta}{\lambda} \]
car $\lambda = cT$. Vérifions l'homogénéité : $\delta$ et $\lambda$ sont des longueurs, donc $\Delta\varphi$ est bien sans dimension (un angle, en radians).

\begin{propriete}[Conditions d'interférences constructives et destructives — démonstration exigible au Bac]
En un point $M$ où se superposent deux ondes issues de sources synchrones (en phase) et cohérentes, le déphasage en $M$ vaut $\Delta\varphi = \dfrac{2\pi\,\delta}{\lambda}$.

\textbf{Vibrations en phase (amplitude maximale) :} les deux vibrations s'ajoutent lorsque $\Delta\varphi = 2k\pi$ avec $k \in \mathbb{Z}$, c'est-à-dire :
\[ \boxed{\ \delta = k\lambda,\quad k \in \mathbb{Z}\ } \qquad \text{(interférences constructives)} \]

\textbf{Vibrations en opposition de phase (amplitude minimale, nulle si les amplitudes sont égales) :} les vibrations s'annulent lorsque $\Delta\varphi = (2k+1)\pi$ avec $k \in \mathbb{Z}$, c'est-à-dire :
\[ \boxed{\ \delta = (2k+1)\,\frac{\lambda}{2},\quad k \in \mathbb{Z}\ } \qquad \text{(interférences destructives)} \]
\end{propriete}

\begin{experience}[Démonstration guidée pas à pas]
\begin{enumerate}
\item Les sources $S_1$ et $S_2$ vibrent en phase : $y_{S_1}(t) = y_{S_2}(t) = a\cos(\omega t)$ (même amplitude $a$ pour simplifier).
\item En $M$, chaque onde reproduit la vibration de sa source avec le retard du trajet : $y_1(M,t) = a\cos(\omega t - \frac{2\pi d_1}{\lambda})$ et $y_2(M,t) = a\cos(\omega t - \frac{2\pi d_2}{\lambda})$.
\item Le déphasage de l'onde 2 par rapport à l'onde 1 en $M$ est donc $\Delta\varphi = \dfrac{2\pi(d_2 - d_1)}{\lambda} = \dfrac{2\pi\delta}{\lambda}$.
\item Si $\Delta\varphi = 2k\pi$, alors $\cos(\omega t - \Delta\varphi) = \cos(\omega t)$ : les deux vibrations montent et descendent \emph{ensemble} ; l'amplitude résultante est $2a$ (maximale). La condition $\frac{2\pi\delta}{\lambda} = 2k\pi$ se simplifie en $\delta = k\lambda$.
\item Si $\Delta\varphi = (2k+1)\pi$, alors $\cos(\omega t - \Delta\varphi) = -\cos(\omega t)$ : quand l'une monte, l'autre descend ; l'amplitude résultante est $0$ (minimale, nulle ici car amplitudes égales). La condition devient $\delta = (2k+1)\frac{\lambda}{2}$.
\end{enumerate}
\end{experience}

\begin{methode}[Un point $M$ vibre-t-il avec une amplitude maximale ou minimale ?]
\begin{enumerate}
\item Mesurer (ou calculer) les distances $d_1 = S_1M$ et $d_2 = S_2M$.
\item Calculer la différence de marche $\delta = |d_2 - d_1|$ (valeur absolue).
\item Calculer la longueur d'onde $\lambda = c/f = cT$ (dans la même unité que $\delta$).
\item Former le rapport $\dfrac{\delta}{\lambda}$ :
\begin{itemize}
\item si $\dfrac{\delta}{\lambda} = k$ est un \textbf{entier} ($0 ; 1 ; 2 ; \ldots$) : amplitude \textbf{maximale} ;
\item si $\dfrac{\delta}{\lambda} = k + 0{,}5$ est un \textbf{demi-entier} ($0{,}5 ; 1{,}5 ; 2{,}5 ; \ldots$) : amplitude \textbf{minimale} ;
\item sinon : amplitude intermédiaire.
\end{itemize}
\end{enumerate}
\end{methode}

\begin{exoresolu}[État vibratoire d'un point de la cuve]
\textbf{Énoncé.} Sur une cuve à ondes, deux pointes $S_1$ et $S_2$ vibrent en phase à la fréquence $f = \SI{20}{Hz}$. Les ondes se propagent à la célérité $c = \SI{40}{cm/s}$. Un point $M$ de la surface est situé à $S_1M = \SI{12}{cm}$ de la première source et $S_2M = \SI{17}{cm}$ de la seconde. Le point $M$ vibre-t-il avec une amplitude maximale ou minimale ?
\tcblower
\textbf{Solution détaillée.}
\begin{enumerate}
\item \textbf{Longueur d'onde :} $\lambda = \dfrac{c}{f} = \dfrac{40}{20} = \SI{2,0}{cm}$.
\item \textbf{Différence de marche :} $\delta = |S_2M - S_1M| = |17 - 12| = \SI{5,0}{cm}$.
\item \textbf{Rapport :} $\dfrac{\delta}{\lambda} = \dfrac{5{,}0}{2{,}0} = 2{,}5 = \dfrac{5}{2}$.
\end{enumerate}
Le rapport est un \emph{demi-entier} : $\delta = 2{,}5\,\lambda = (2\times 2 + 1)\dfrac{\lambda}{2}$, ce qui correspond à la condition d'interférences destructives avec $k = 2$.

\textbf{Conclusion :} les deux ondes arrivent en $M$ en opposition de phase ; le point $M$ vibre avec une amplitude \textbf{minimale} (pratiquement nulle si les deux ondes y ont la même amplitude) : $M$ appartient à une zone calme de la figure.
\end{exoresolu}

\begin{attention}[Le piège classique du Bac]
L'erreur la plus fréquente est d'\emph{inverser} les deux conditions. Retiens ce moyen mnémotechnique : si le chemin supplémentaire est un nombre \textbf{entier de longueurs d'onde} ($\delta = k\lambda$), les ondes « retombent ensemble » : elles sont \textbf{en phase}, l'amplitude est \textbf{maximale}. Si le chemin supplémentaire est un nombre \textbf{demi-entier de longueurs d'onde} ($\delta = (2k+1)\lambda/2$), elles arrivent « à contre-temps » : \textbf{opposition de phase}, amplitude \textbf{minimale}. Autre piège : $k$ appartient à $\mathbb{Z}$ tout entier (positif, nul \emph{ou négatif}) ; en pratique, avec $\delta = |d_2 - d_1| \geq 0$, on se limite à $k \in \mathbb{N}$. Enfin, n'oublie pas de convertir toutes les longueurs dans la même unité avant de former le rapport $\delta/\lambda$.
\end{attention}

\courssub{Forme géométrique des franges : des hyperboles}

Cherchons le lieu géométrique des points où l'amplitude est maximale : ce sont les points tels que $\delta = |S_2M - S_1M| = k\lambda$, c'est-à-dire les points dont la \emph{différence des distances à deux points fixes est constante}. Or c'est précisément la définition géométrique d'une \cle{hyperbole} de foyers $S_1$ et $S_2$.

\begin{propriete}[Nature des franges d'interférences]
Les points de la surface de l'eau vibrant avec une amplitude maximale (resp. minimale) se répartissent sur des courbes appelées \cle{franges d'interférences}, qui sont des \textbf{hyperboles de foyers} $S_1$ et $S_2$ :
\begin{itemize}
\item franges d'amplitude \textbf{maximale} : $|S_2M - S_1M| = k\lambda$, $k \in \mathbb{N}$ ;
\item franges d'amplitude \textbf{minimale} : $|S_2M - S_1M| = (2k+1)\dfrac{\lambda}{2}$, $k \in \mathbb{N}$.
\end{itemize}
Cas particulier fondamental : pour $k = 0$, la condition $|S_2M - S_1M| = 0$ signifie $S_1M = S_2M$ : c'est la \textbf{médiatrice du segment $[S_1S_2]$}. C'est donc toujours une frange d'amplitude maximale (la \emph{frange centrale}, d'ordre 0) lorsque les sources vibrent en phase.
\end{propriete}

\begin{popfigure}
\begin{center}
\begin{tikzpicture}[scale=0.85]
% Foyers
\coordinate (S1) at (-1.2,0);
\coordinate (S2) at (1.2,0);
% Hyperboles maxima (pleines) : delta = k*lambda avec lambda = 0.8
\foreach \a/\col in {0.4/poporange, 0.8/poporange, 1.2/poporange}{
  \draw[\col, thick, domain=-1.6:1.6, samples=60, smooth, variable=\t]
    plot ({\a*cosh(\t)}, {sqrt(max(0,\a*\a + 1.44))*sinh(\t)*0.9});
  \draw[\col, thick, domain=-1.6:1.6, samples=60, smooth, variable=\t]
    plot ({-\a*cosh(\t)}, {sqrt(max(0,\a*\a + 1.44))*sinh(\t)*0.9});
}
% Frange centrale (médiatrice)
\draw[poporange, very thick] (0,-2.6) -- (0,2.6);
\node[poporange, rotate=90, left=2pt] at (0,0) {\small frange centrale $\delta = 0$};
% Hyperboles minima (pointillés)
\foreach \a in {0.2, 0.6, 1.0}{
  \draw[popblue, thick, dashed, domain=-1.5:1.5, samples=60, smooth, variable=\t]
    plot ({\a*cosh(\t)}, {sqrt(max(0,\a*\a + 1.44))*sinh(\t)*0.9});
  \draw[popblue, thick, dashed, domain=-1.5:1.5, samples=60, smooth, variable=\t]
    plot ({-\a*cosh(\t)}, {sqrt(max(0,\a*\a + 1.44))*sinh(\t)*0.9});
}
% Sources
\fill[popink] (S1) circle (0.09) node[below] {$S_1$};
\fill[popink] (S2) circle (0.09) node[below] {$S_2$};
% Légende
\node[poporange] at (4.6,1.6) {\small maxima : $\delta = k\lambda$};
\node[popblue] at (4.6,0.9) {\small minima : $\delta = (2k+1)\frac{\lambda}{2}$};
\end{tikzpicture}
\end{center}
\legende{Figure d'interférences : les franges d'amplitude maximale (traits pleins) et minimale (pointillés) sont des hyperboles de foyers $S_1$ et $S_2$. La médiatrice de $[S_1S_2]$ est la frange maximale centrale ($\delta = 0$).}
\end{popfigure}

\courssub{Interfrange et exploitation de la figure d'interférences}

Lorsque l'on observe la figure d'interférences \emph{loin des sources} (distance $D$ à un écran ou zone d'observation telle que $D \gg a = S_1S_2$), les hyperboles deviennent pratiquement des droites parallèles à la médiatrice de $[S_1S_2]$. L'espacement entre deux franges brillantes (ou deux franges sombres) consécutives devient constant : c'est l'\cle{interfrange}, notée $i$.

\begin{propriete}[Interfrange en champ lointain — formule exigible]
Dans l'approximation $D \gg a$ (petits angles), la différence de marche pour un point $M$ d'abscisse $y$ sur un écran parallèle à $S_1S_2$ vaut $\delta \approx \dfrac{a\,y}{D}$.
Les maxima ($k$-ième frange brillante) sont obtenus pour $\delta = k\lambda$, soit :
\[ y_k = k\,\frac{\lambda D}{a} \]
L'interfrange $i$ (distance entre deux franges brillantes consécutives) est donc :
\[ \boxed{\ i = \frac{\lambda D}{a}\ } \]
On en déduit une méthode classique pour mesurer la longueur d'onde : $\lambda = i\,\dfrac{a}{D}$.
\end{propriete}

\begin{methode}[Exploiter une figure d'interférences pour déterminer $\lambda$]
\begin{enumerate}
\item Mesurer l'écartement des sources $a = S_1S_2$.
\item Mesurer la distance $D$ sources--écran (ou zone d'observation).
\item Sur l'écran, mesurer la distance $L$ qui sépare deux franges brillantes d'ordres connus $k_1$ et $k_2$ (par exemple 5 franges : $L = 5i$).
\item Calculer l'interfrange $i = L / (k_2 - k_1)$.
\item En déduire $\lambda = i\,a/D$.
\end{enumerate}
\end{methode}

\begin{exoresolu}[Mesure de la longueur d'onde par interfrange]
\textbf{Énoncé.} Dans une cuve à ondes, $S_1S_2 = \SI{3,0}{cm}$. On observe la figure sur un écran placé à $D = \SI{50}{cm}$ des sources. On compte 10 franges brillantes sur une longueur $L = \SI{12,0}{cm}$. Déterminer $\lambda$.
\tcblower
\textbf{Solution.}
L'interfrange est $i = L / 10 = \SI{1,2}{cm}$.
La formule $i = \lambda D / a$ donne $\lambda = i\,a/D = 1,2 \times 3,0 / 50 = \SI{0,072}{cm} = \SI{0,72}{mm}$.
\end{exoresolu}

\courssub{Observation en éclairage normal et en éclairage stroboscopique}

\textbf{En éclairage normal (continu) :} l'œil ne perçoit pas le mouvement rapide de la surface (la fréquence $f$ est de l'ordre de quelques dizaines de hertz), mais il distingue nettement l'\emph{alternance de zones agitées et de zones calmes} : les zones agitées apparaissent floues ou brillantes sur l'écran (effet de lentille variable), les zones calmes restent nettes et sombres. On observe directement la figure d'interférences et l'on peut repérer les franges.

\textbf{En éclairage stroboscopique :} le stroboscope émet des éclairs brefs à la fréquence $f_s$.
\begin{itemize}
\item Si $f_s = f$ (ou un sous-multiple exact $f/n$), la surface est toujours éclairée dans le même état vibratoire : la figure paraît \textbf{figée} ; on voit les crêtes et les creux immobiles, ce qui permet de mesurer $\lambda$ directement (distance entre deux crêtes consécutives d'une même onde).
\item Si $f_s$ est légèrement différente de $f$ (par exemple $f_s = f \pm 1\,\text{Hz}$), la figure semble évoluer au \textbf{ralenti}, avec un mouvement apparent d'autant plus lent que l'écart $|f_s - f|$ est faible : c'est idéal pour analyser la propagation des ondes et visualiser le sens de déplacement des franges.
\end{itemize}

\begin{savaistu}[Des vagues de la Sanaga aux ondes lumineuses]
Le phénomène d'interférences n'est pas qu'une curiosité de laboratoire : à l'embouchure de la Sanaga ou près des côtes de Kribi, on peut observer les houles diffractées par des obstacles (rochers, jetées) se superposer et dessiner des figures d'interférences géantes. Surtout, c'est exactement le même raisonnement — différence de marche, déphasage $\Delta\varphi = 2\pi\delta/\lambda$ — qui explique les interférences \emph{lumineuses} (franges d'Young) étudiées plus tard dans l'année : la mécanique des ondes sur l'eau est le modèle intuitif de toute l'optique ondulatoire. Les radars de la météo nationale (Douala, Yaoundé) et les sonars des navires de la marine camerounaise exploitent aussi ces principes pour la télémétrie.
\end{savaistu}

\begin{aretenir}[L'essentiel sur les interférences à la surface de l'eau]
\begin{itemize}
\item Deux sources \textbf{cohérentes} (même vibreur, donc même fréquence, en phase) créent une figure d'interférences \textbf{stable}.
\item En un point $M$, tout est gouverné par la différence de marche $\delta = |S_2M - S_1M|$.
\item Amplitude \textbf{maximale} si $\delta = k\lambda$ ; amplitude \textbf{minimale} si $\delta = (2k+1)\dfrac{\lambda}{2}$, avec $k \in \mathbb{Z}$.
\item Les franges sont des \textbf{hyperboles de foyers $S_1$ et $S_2$} ; la médiatrice de $[S_1S_2]$ est la frange maximale centrale.
\item En champ lointain ($D \gg a$), les franges sont rectilignes, parallèles et équidistantes : l'interfrange vaut $i = \lambda D / a$.
\end{itemize}
\end{aretenir}

\begin{exercice}[À toi de jouer]
Sur une cuve à ondes, deux pointes synchrones vibrent à $f = \SI{25}{Hz}$ ; la célérité des ondes est $c = \SI{30}{cm/s}$.
\begin{enumerate}
\item Calculer la longueur d'onde $\lambda$.
\item Un point $A$ vérifie $S_1A = \SI{8,4}{cm}$ et $S_2A = \SI{11,4}{cm}$. Déterminer son état vibratoire.
\item Un point $B$ est situé sur la médiatrice de $[S_1S_2]$. Que vaut $\delta$ en $B$ ? Quel est son état vibratoire ?
\item On place un écran à $D = \SI{60}{cm}$ des sources écartées de $a = \SI{4,0}{cm}$. Calculer l'interfrange $i$ attendue sur cet écran.
\end{enumerate}
\end{exercice}

\begin{corrige}[Exercice : états vibratoires et interfrange sur la cuve à ondes]
\begin{enumerate}
\item $\lambda = \dfrac{c}{f} = \dfrac{30}{25} = \SI{1,2}{cm}$.
\item $\delta = |11{,}4 - 8{,}4| = \SI{3,0}{cm}$ ; $\dfrac{\delta}{\lambda} = \dfrac{3{,}0}{1{,}2} = 2{,}5$ : demi-entier, donc $\delta = (2\times 2 + 1)\dfrac{\lambda}{2}$ : $A$ vibre avec une amplitude \textbf{minimale}.
\item Sur la médiatrice, $S_1B = S_2B$ donc $\delta = 0 = 0\times\lambda$ : $B$ vibre avec une amplitude \textbf{maximale} (frange centrale d'ordre 0).
\item $i = \dfrac{\lambda D}{a} = \dfrac{1,2 \times 60}{4,0} = \SI{18}{cm}$.
\end{enumerate}
\end{corrige}

\coursec{Interfrange et exploitation de la figure d'interférences}

Dans la section précédente, nous avons établi la géométrie de la figure d'interférences à la surface de l'eau : les points de vibration maximale et les points de repos se répartissent sur des familles d'hyperboles de foyers $S_1$ et $S_2$, caractérisées par la différence de marche $\delta = d_2 - d_1$. Sur un écran d'observation placé loin des sources, ces hyperboles sont localement assimilables à des droites parallèles : on observe alors des \cle{franges} rectilignes, alternativement brillantes (vibration maximale) et sombres (repos). La question qui se pose maintenant est toute simple : \emph{quelle est la distance entre deux franges consécutives, et que peut-on en tirer expérimentalement ?} Cette distance, appelée \cle{interfrange}, est la clé qui permet de remonter à la longueur d'onde $\lambda$, puis à la célérité des ondes. C'est l'outil de mesure par excellence de ce chapitre.

\courssub{Définition de l'interfrange}

\begin{definition}[Interfrange]
On appelle \cle{interfrange}, notée $i$, la distance entre deux franges consécutives \textbf{de même nature} (deux franges de vibration maximale consécutives, ou deux franges de repos consécutives), mesurée sur un axe \textbf{perpendiculaire aux franges}. L'interfrange s'exprime en mètres (m), ou en pratique en centimètres à la surface de l'eau.
\end{definition}

L'insistance sur « de même nature » n'est pas un détail de langage : une frange sombre et la frange brillante voisine sont séparées de $i/2$, pas de $i$. Nous y reviendrons dans la boîte « Attention ».

\begin{attention}[Le piège classique de la mesure de l'interfrange]
L'erreur la plus fréquente au Baccalauréat consiste à mesurer la distance entre une frange de vibration maximale et la frange de repos \emph{voisine} : on obtient alors $i/2$ et non $i$, ce qui fausse tous les calculs qui suivent (on trouve une longueur d'onde deux fois trop petite). Deux bonnes habitudes pour l'éviter :
\begin{itemize}
\item mesurer la distance $L$ occupée par $n$ interfranges entières (par exemple entre la frange brillante d'ordre $0$ et celle d'ordre $n$), puis calculer $i = L/n$ : cela réduit en outre l'incertitude de mesure ;
\item vérifier que les deux franges repérées sont bien de \emph{même nature} (toutes les deux brillantes, ou toutes les deux sombres).
\end{itemize}
\end{attention}

\courssub{Expression de l'interfrange : $i = \lambda D / a$}

Notons $a = S_1S_2$ la distance entre les deux sources et $D$ la distance des sources à l'écran d'observation (ou à la zone où l'on mesure les franges, perpendiculairement à la médiatrice de $S_1S_2$).

\begin{propriete}[Interfrange (admise)]
Dans les conditions usuelles d'observation ($D \gg a$ et observation à faible distance de l'axe), l'interfrange est donnée par :
\[ i = \frac{\lambda\, D}{a} \]
où $\lambda$ est la longueur d'onde, $D$ la distance sources--écran et $a = S_1S_2$ l'écartement des sources. Cette formule est \textbf{admise} au programme ; une démonstration guidée est proposée en complément ci-dessous.
\end{propriete}

Vérifions l'homogénéité, réflexe indispensable : $\lambda$, $D$ et $a$ sont toutes trois des longueurs, donc $\dfrac{[\lambda]\times[D]}{[a]} = \dfrac{\text{m}\times\text{m}}{\text{m}} = \text{m}$, qui est bien une longueur. La formule est dimensionnellement correcte.

\begin{savaistu}[La même formule en optique]
Cette relation est exactement celle des interférences lumineuses dans l'expérience des fentes d'Young : $i = \lambda D/a$. Ce n'est pas une coïncidence : la forme de la figure d'interférences ne dépend que de la géométrie (deux sources cohérentes séparées de $a$, observation à distance $D$) et de la longueur d'onde, pas de la nature de l'onde. C'est une illustration magnifique de l'universalité des phénomènes ondulatoires : la cuve à ondes du laboratoire est un « modèle réduit » des expériences d'optique qui ont établi la nature ondulatoire de la lumière au XIX\textsuperscript{e} siècle.
\end{savaistu}

\courssub{Démonstration guidée (complément Bac)}

Cette démonstration n'est pas exigible, mais la suivre est très formateur : elle repose sur une approximation géométrique simple.

\begin{experience}[Démonstration guidée de $i = \lambda D/a$]
On se place dans le plan contenant les sources et l'écran. Les sources $S_1$ et $S_2$ sont distantes de $a$ ; l'écran est situé à la distance $D$ du segment $[S_1S_2]$, perpendiculairement à sa médiatrice. Un point $M$ de l'écran est repéré par son abscisse $x$ comptée depuis l'axe de symétrie.

\textbf{Étape 1 : différence de marche en $M$.} Les distances de $M$ aux deux sources sont :
\[ d_1 = \sqrt{D^2 + \left(x - \tfrac{a}{2}\right)^2} \qquad \text{et} \qquad d_2 = \sqrt{D^2 + \left(x + \tfrac{a}{2}\right)^2}. \]
Comme $D \gg a$ et $D \gg x$, on utilise l'approximation $\sqrt{1+\varepsilon} \approx 1 + \varepsilon/2$ pour $\varepsilon \ll 1$ :
\[ d_2 - d_1 \approx \frac{1}{2D}\left[\left(x+\tfrac{a}{2}\right)^2 - \left(x-\tfrac{a}{2}\right)^2\right] = \frac{1}{2D}\times 2ax = \frac{ax}{D}. \]
La différence de marche vaut donc $\boxed{\delta \approx \dfrac{ax}{D}}$.

\textbf{Étape 2 : positions des franges brillantes.} La condition de vibration maximale est $\delta = k\lambda$ avec $k \in \mathbb{Z}$, soit :
\[ x_k = k\,\frac{\lambda D}{a}. \]
Les franges brillantes sont donc équidistantes, d'abscisses proportionnelles à $k$.

\textbf{Étape 3 : interfrange.} La distance entre deux franges brillantes consécutives (ordres $k$ et $k+1$) est :
\[ i = x_{k+1} - x_k = (k+1)\frac{\lambda D}{a} - k\frac{\lambda D}{a} = \frac{\lambda D}{a}. \]
Le résultat est indépendant de $k$ : les franges sont régulièrement espacées. \emph{CQFD.}
\end{experience}

\begin{popfigure}
\begin{center}
\begin{tikzpicture}[scale=1.0, >=Stealth]
% Sources
\coordinate (S1) at (0,1.0);
\coordinate (S2) at (0,-1.0);
\coordinate (O) at (0,0);
% Ecran
\draw[line width=1.5pt, popblue] (7,-2.6) -- (7,2.6);
\node[popblue, anchor=west] at (7.05,2.4) {\small écran};
% Point M
\coordinate (M) at (7,1.6);
\fill[poporange] (M) circle (2pt) node[anchor=south west, popdark]{\small $M$};
% Rayons
\draw[popdark] (S1) -- (M) node[midway, above, sloped]{\small $d_1$};
\draw[popdark] (S2) -- (M) node[midway, below, sloped]{\small $d_2$};
% Axe
\draw[dashed, popmuted] (0,0) -- (7,0);
\coordinate (H) at (7,0);
\fill (H) circle (1.5pt) node[below, popdark]{\small $O'$};
% Sources points
\fill[popgreen] (S1) circle (2.5pt) node[left, popdark]{\small $S_1$};
\fill[popgreen] (S2) circle (2.5pt) node[left, popdark]{\small $S_2$};
% Cotes a, D, x
\draw[<->, poppurple] (-0.5,1.0) -- (-0.5,-1.0) node[midway, left]{\small $a$};
\draw[<->, poppurple] (0,-2.9) -- (7,-2.9) node[midway, below]{\small $D$};
\draw[<->, poporange] (7.35,0) -- (7.35,1.6) node[midway, right]{\small $x$};
% Angle
\draw[popmuted] (2.2,0) arc (0:12.9:2.2);
\node[popmuted] at (2.75,0.28) {\small $\theta$};
\end{tikzpicture}
\end{center}
\legende{Géométrie de la démonstration : deux sources cohérentes $S_1$ et $S_2$ distantes de $a$, écran placé à la distance $D$. Au point $M$ d'abscisse $x$, la différence de marche est $\delta = d_2 - d_1 \approx ax/D$ pour $D \gg a$.}
\end{popfigure}

\courssub{Influence des paramètres sur l'interfrange}

La relation $i = \dfrac{\lambda D}{a}$ se lit directement :

\begin{itemize}
\item \textbf{Influence de $\lambda$ (donc de la fréquence)} : $i$ est proportionnelle à $\lambda$. À célérité $c$ fixée, $\lambda = c/f$ : \emph{plus la fréquence du vibreur est faible, plus les franges sont écartées}. À la cuve à ondes, diminuer la fréquence « étale » la figure.
\item \textbf{Influence de $D$} : $i$ est proportionnelle à $D$. Plus on observe loin des sources, plus la figure est grande : c'est pourquoi on projette la figure sur un écran éloigné pour la mesurer confortablement.
\item \textbf{Influence de $a$} : $i$ est inversement proportionnelle à $a$. \emph{Rapprocher les deux pointes du vibreur écarte les franges} ; les écarter resserre la figure. Si $a$ devient trop grand, les franges deviennent si fines qu'elles ne sont plus discernables.
\end{itemize}

\begin{popfigure}
\begin{center}
\begin{tikzpicture}
\begin{axis}[
  width=0.85\linewidth, height=6.5cm,
  xlabel={$1/a$ \quad (cm$^{-1}$)},
  ylabel={$i$ \quad (cm)},
  xmin=0, xmax=0.5, ymin=0, ymax=5,
  xtick={0,0.1,0.2,0.3,0.4,0.5},
  ytick={0,1,2,3,4,5},
  grid=both, grid style={popline!40},
  axis line style={popdark},
  label style={popdark},
  tick label style={popdark},
  legend style={at={(0.03,0.97)}, anchor=north west, draw=popline}
]
\addplot[popcurve, domain=0:0.5, samples=50] {9*x};
\addlegendentry{$i = \lambda D \times \dfrac{1}{a}$ \; ($\lambda D = 9$ cm$^2$)}
\addplot[only marks, mark=*, mark size=2pt, poporange] coordinates {(0.1,0.9) (0.2,1.8) (0.25,2.25) (0.333,3.0) (0.4,3.6)};
\addlegendentry{points expérimentaux}
\end{axis}
\end{tikzpicture}
\end{center}
\legende{Interfrange $i$ en fonction de $1/a$ : les points expérimentaux s'alignent sur une droite passant par l'origine, dont la pente vaut $\lambda D$. C'est une vérification expérimentale directe de la loi $i = \lambda D/a$.}
\end{popfigure}

\courssub{Exploitation expérimentale : mesurer $\lambda$ puis $c$}

Renversons la perspective : au lieu de prédire $i$, utilisons la figure d'interférences comme un \emph{instrument de mesure}.

\begin{methode}[Détermination de la longueur d'onde et de la célérité]
\begin{enumerate}
\item Réaliser la figure d'interférences à la cuve à ondes avec deux pointes solidaires du même vibreur (sources cohérentes), de fréquence $f$ connue (lue sur le GBF ou la notice du vibreur).
\item Mesurer l'écartement des pointes $a = S_1S_2$ et la distance $D$ entre les sources et la zone d'observation (écran ou règle posée sur la cuve).
\item Mesurer la longueur $L$ occupée par $n$ interfranges entières (entre deux franges de même nature) ; en déduire $i = L/n$.
\item Calculer la longueur d'onde : $\lambda = \dfrac{i\,a}{D}$.
\item Calculer la célérité des ondes à la surface de l'eau : $c = \lambda f$.
\end{enumerate}
Toutes les distances doivent être exprimées dans la \emph{même unité} (par exemple le cm) ; $\lambda$ s'obtient alors dans cette unité, et $c$ en cm/s si $f$ est en Hz.
\end{methode}

\begin{exemplebox}[Exemple chiffré]
Sur une cuve à ondes, deux pointes distantes de $a = \SI{4}{cm}$ vibrent en phase à la fréquence $f = \SI{25}{Hz}$. On observe la figure sur un écran situé à $D = \SI{30}{cm}$ des sources et l'on mesure un interfrange $i = \SI{1,5}{cm}$.

Longueur d'onde :
\[ \lambda = \frac{i\,a}{D} = \frac{1,5 \times 4}{30} = \SI{0,2}{cm}. \]
Célérité :
\[ c = \lambda f = 0,2 \times 25 = \SI{5}{cm/s} = \SI{5e-2}{m/s}. \]
Vérification de cohérence : $i = \lambda D/a = 0,2 \times 30 / 4 = \SI{1,5}{cm}$ : on retrouve bien la mesure.
\end{exemplebox}

\begin{exoresolu}[Mesure de la célérité à la cuve à ondes]
Un élève de Terminale C réalise des interférences à la surface de l'eau avec un vibreur de fréquence $f = \SI{20}{Hz}$. Les deux pointes sont distantes de $a = \SI{5,0}{cm}$. Sur l'écran placé à $D = \SI{40}{cm}$, il mesure une distance $L = \SI{9,6}{cm}$ entre la frange brillante centrale et la sixième frange brillante (soit $n = 6$ interfranges). Déterminer l'interfrange, la longueur d'onde et la célérité des ondes.
\tcblower
\textbf{Solution détaillée.}

\textbf{1. Interfrange.} La distance $L$ couvre $n = 6$ interfranges entières :
\[ i = \frac{L}{n} = \frac{9,6}{6} = \SI{1,6}{cm}. \]

\textbf{2. Longueur d'onde.} D'après $i = \lambda D/a$ :
\[ \lambda = \frac{i\,a}{D} = \frac{1,6 \times 5,0}{40} = \SI{0,20}{cm} = \SI{2,0e-3}{m}. \]

\textbf{3. Célérité.} 
\[ c = \lambda f = 2,0\times 10^{-3} \times 20 = \SI{4,0e-2}{m/s} = \SI{4,0}{cm/s}. \]

\textbf{Vérification.} Les ordres de grandeur sont typiques des ondes à la surface de l'eau (quelques cm/s à quelques dizaines de cm/s selon la profondeur) : le résultat est physiquement plausible.
\end{exoresolu}

\courssub{Observation stroboscopique de la figure}

À l'œil nu, en éclairage continu, la surface de l'eau vibre trop vite pour que l'on distingue le détail du mouvement : on perçoit seulement des zones agitées et des zones calmes. Le \cle{stroboscope} — source lumineuse émettant des éclairs brefs à fréquence réglable $f_s$ — permet de « figer » le mouvement.

\begin{propriete}[Observation stroboscopique]
Soit $f$ la fréquence du vibreur et $f_s$ la fréquence des éclairs du stroboscope.
\begin{itemize}
\item Si $f_s = f$ : entre deux éclairs, la surface a effectué exactement une période de vibration ; chaque éclair éclaire la surface dans le même état. La figure d'interférences apparaît \textbf{immobile} : franges brillantes et sombres nettement figées.
\item Si $f_s$ est \textbf{légèrement différente} de $f$ : d'un éclair au suivant, la surface a effectué un peu plus (ou un peu moins) d'une période ; on observe un \textbf{mouvement apparent ralenti} des ondes, dans le sens réel si $f_s \lesssim f$, en sens inverse si $f_s \gtrsim f$.
\item Si $f_s = f/k$ ($k$ entier) : la figure paraît également immobile, mais une onde sur $k$ est éclairée.
\end{itemize}
\end{propriete}

\begin{attention}[Immobilité apparente : une condition à bien rédiger]
Au Baccalauréat, il ne suffit pas d'écrire « le stroboscope fige le mouvement ». Il faut préciser la \emph{condition} : l'immobilité apparente est obtenue lorsque la fréquence des éclairs est égale à celle du vibreur (ou à un sous-multiple $f/k$). Si $f_s = 2f$, chaque point est éclairé deux fois par période dans deux états opposés : la figure peut paraître brouillée. Et si $f_s$ est proche de $f$ sans lui être égale, c'est le mouvement \emph{ralenti} qui est observé — c'est d'ailleurs ainsi que l'on mesure finement la fréquence du vibreur : on règle $f_s$ jusqu'à obtenir l'immobilité.
\end{attention}

\begin{aretenir}[L'essentiel de la section]
\begin{itemize}
\item L'\cle{interfrange} $i$ est la distance entre deux franges consécutives \textbf{de même nature}, mesurée perpendiculairement aux franges.
\item Formule (admise) : $i = \dfrac{\lambda D}{a}$, avec $a = S_1S_2$ et $D$ la distance sources--écran. Elle découle de $\delta \approx ax/D$ et de la condition $\delta = k\lambda$.
\item Exploitation : mesurer $i$, $a$, $D$ donne $\lambda = \dfrac{ia}{D}$, puis $c = \lambda f$.
\item $i$ augmente avec $\lambda$ (fréquence faible) et avec $D$ ; $i$ diminue quand $a$ augmente.
\item Stroboscope : $f_s = f$ donne une figure immobile ; $f_s \approx f$ donne un mouvement apparent ralenti.
\end{itemize}
\end{aretenir}

\coursec{Ondes stationnaires : expérience de Melde}

Après avoir étudié les interférences entre deux ondes émises par des sources cohérentes distinctes, nous allons maintenant analyser une situation particulière d'interférence : celle où une onde interfère avec sa propre réflexion. C'est le phénomène des \emph{ondes stationnaires}, mis en évidence de façon spectaculaire par l'expérience de Melde. Tu verras que, contrairement aux ondes progressives, l'énergie ne se propage plus le long du milieu : elle reste « coincée » entre des points immobiles.

\courssub{Dispositif et observation de l'expérience de Melde}

\begin{experience}[Expérience de Melde — Ondes stationnaires sur une corde]
\textbf{Matériel :} un vibreur électromagnétique (fréquence $f$ réglable, typiquement \SIrange{20}{200}{Hz}), une corde légère et souple (longueur $L \approx \SI{1,5}{m}$), une poulie fixe, un jeu de masses marquées (pour régler la tension $F$), un chronomètre, une règle graduée.

\textbf{Montage :} La corde est tendue horizontalement. Une extrémité est fixée sur le vibreur (qui impose un mouvement vertical sinusoïdal $y(0,t)=a\cos(2\pi f t)$), l'autre passe sur la poulie et supporte une masse $m$ qui assure la tension $F=mg$. La longueur vibrante $L$ (distance vibreur--poulie) est mesurée.

\textbf{Protocole :}
\begin{enumerate}
\item Fixer $m$ (donc $F$) et $L$. Faire varier la fréquence $f$ du vibreur en partant de valeurs basses.
\item Observer la corde : pour la plupart des fréquences, la corde vibre de façon confuse, sans forme nette (amplitude faible).
\item Pour certaines fréquences bien précises $f_1, f_2, f_3,\dots$, la corde adopte une forme stable, composée de $n$ fuseaux (ou boucles) bien nets. Les points de la corde semblent soit immobiles, soit vibrer avec une grande amplitude.
\item Noter le nombre $n$ de fuseaux pour chaque fréquence de résonance. Répéter l'expérience en changeant $m$ (donc $F$) ou $L$.
\end{enumerate}

\textbf{Observations clés :}
\begin{itemize}
\item Les points immobiles sont appelés \textbf{nœuds} (notés N). Ils sont toujours au nombre de $n+1$ (y compris les extrémités).
\item Les points d'amplitude maximale sont les \textbf{ventres} (notés V). Ils sont au nombre de $n$ (un par fuseau).
\item La distance entre deux nœuds consécutifs (ou deux ventres consécutifs) est constante pour une résonance donnée.
\item Si on éclaire la corde avec une lampe stroboscopique réglée sur la fréquence $f$, la corde apparaît \emph{figée} dans sa forme d'enveloppe ; sans stroboscope, on voit un flou entre les nœuds.
\end{itemize}

\textbf{Interprétation physique :} Le vibreur génère une \textbf{onde incidente} qui se propage vers la poulie. À l'extrémité fixe (poulie), l'onde se réfléchit en \textbf{onde réfléchie} de même fréquence, même amplitude (si on néglige l'amortissement), mais en \textbf{opposition de phase} (inversion de l'élongation) car l'extrémité est bloquée. La superposition de ces deux ondes progressives de même fréquence, se propageant en sens opposés, crée une \textbf{onde stationnaire}.
\end{experience}

\begin{popfigure}
\begin{tikzpicture}[scale=0.9, >=Stealth]
% Base
\draw[popdark, thick] (-0.5,-1.5) rectangle (14.5,-1.3);
% Vibreur (gauche)
\draw[popblue, thick, fill=popblueL] (0.5,-1.3) rectangle (2.5,0.5);
\draw[popblue, thick] (1.5,-1.3) -- (1.5,-1.8) node[below] {Vibreur};
\node[align=center] at (1.5,0.8) {Générateur\\de fréquence $f$};
% Flèche mouvement vertical vibreur
\draw[poporange, thick, <->] (1.5,0.7) -- (1.5,1.2) node[above] {$y(0,t)$};
% Corde
\draw[popdark, very thick] (2.5,0) -- (12.5,0);
% Poulie (droite)
\draw[poppurple, thick] (12.5,-0.3) circle (0.8);
\draw[poppurple, thick, fill=poppurpleL] (12.5,-0.3) circle (0.3);
\draw[poppurple, thick] (12.5,0.5) -- (12.5,1.5);
\node[align=center] at (12.5,1.9) {Poulie\\libre};
% Masse suspendue
\draw[popdark, thick, fill=popgoldL] (12.0,1.5) rectangle (13.0,3.0);
\node at (12.5,2.25) {$m$};
\draw[popdark, thick, ->] (12.5,3.0) -- (12.5,3.5) node[right] {$F=mg$};
% Nœuds et ventres (exemple 3 fuseaux)
\foreach \x/\label in {2.5/N_0, 5.83/N_1, 9.17/N_2, 12.5/N_3} {
    \draw[popgreen, thick] (\x,-0.2) -- (\x,0.2);
    \node[popgreen, below] at (\x,-0.4) {$\label$};
}
\foreach \x/\label in {4.17/V_1, 7.5/V_2, 10.83/V_3} {
    \draw[poppink, thick] (\x,-0.2) -- (\x,0.2);
    \node[poppink, above] at (\x,0.6) {$\label$};
}
% Enveloppe (forme à t=0)
\draw[popcurve, thick, domain=2.5:12.5, samples=200, variable=\x] plot ({\x}, {0.8*sin(deg(3*pi*(\x-2.5)/10))});
\draw[popcurve, thick, domain=2.5:12.5, samples=200, variable=\x, dashed] plot ({\x}, {-0.8*sin(deg(3*pi*(\x-2.5)/10))});
% Longueur L
\draw[<->, popmuted] (2.5,-1.0) -- (12.5,-1.0) node[midway, below] {$L$};
% Lambda/2
\draw[<->, popblue] (2.5,0.9) -- (5.83,0.9) node[midway, above] {$\lambda/2$};
\draw[<->, popblue] (5.83,0.9) -- (9.17,0.9) node[midway, above] {$\lambda/2$};
% Lambda/4
\draw[<->, poporange] (2.5,1.3) -- (4.17,1.3) node[midway, above] {$\lambda/4$};
% Légendes
\node[align=left, anchor=north west] at (0.2,1.5) {
    \textbf{Légende :}\\
    \textcolor{popgreen}{■} Nœuds (N) : amplitude nulle\\
    \textcolor{poppink}{■} Ventres (V) : amplitude max\\
    \textcolor{popcurve}{—} Enveloppe de l'onde stationnaire
};
\end{tikzpicture}
\legende{Dispositif de Melde : corde tendue entre un vibreur (gauche) et une poulie avec masse $m$ (droite). Pour $f=f_3$, la corde vibre en $n=3$ fuseaux. Les nœuds (N) sont espacés de $\lambda/2$, un nœud et le ventre adjacent de $\lambda/4$.}
\end{popfigure}

\courssub{De la superposition à l'onde stationnaire : modélisation mathématique}

Rappelons le \cle{principe de superposition} : si deux ondes $y_1(x,t)$ et $y_2(x,t)$ se propagent dans le même milieu, l'élongation résultante est $y(x,t)=y_1(x,t)+y_2(x,t)$. Ici, l'onde incidente (vers $+x$) et l'onde réfléchie (vers $-x$) s'écrivent (en négligeant l'amortissement et en choisissant l'origine des temps pour simplifier l'écriture) :
\[
y_i(x,t) = a \sin\!\left( \frac{2\pi}{\lambda}x - \omega t \right), \qquad
y_r(x,t) = a \sin\!\left( -\frac{2\pi}{\lambda}x - \omega t + \pi \right) \quad (\omega = 2\pi/T)
\]
Le terme $+\pi$ traduit l'inversion de phase (changement de signe) à la réflexion sur point fixe. En utilisant $\sin(\alpha+\pi)=-\sin\alpha$ et la parité du sinus ($\sin(-\theta)=-\sin\theta$), l'onde réfléchie devient :
\[
y_r(x,t) = -a \sin\!\left( -\frac{2\pi}{\lambda}x - \omega t \right) = a \sin\!\left( \frac{2\pi}{\lambda}x + \omega t \right).
\]
La superposition donne alors :
\begin{align*}
y(x,t) &= a\left[ \sin\!\left( \frac{2\pi}{\lambda}x - \omega t \right) + \sin\!\left( \frac{2\pi}{\lambda}x + \omega t \right) \right] \\
&= 2a \underbrace{\sin\!\left( \frac{2\pi}{\lambda}x \right)}_{\text{enveloppe spatiale } A(x)} \underbrace{\cos(\omega t)}_{\text{oscillation temporelle}}.
\end{align*}

\begin{definition}[Onde stationnaire]
Une \cle{onde stationnaire} résulte de la superposition de deux ondes progressives sinusoïdales de \textbf{même fréquence} (donc même $\lambda$ et même $T$), de \textbf{même amplitude}, se propageant en \textbf{sens opposés}. Son élongation s'écrit sous la forme factorisée :
\[
\boxed{y(x,t) = A(x) \cos(\omega t + \varphi_0) \quad \text{avec} \quad A(x) = 2a \sin\!\left( \frac{2\pi}{\lambda}x \right)}
\]
Chaque point $x$ oscille en régime sinusoïdal forcé de pulsation $\omega$, mais son \textbf{amplitude $A(x)$ dépend de sa position}. Il n'y a \textbf{pas de propagation} de la forme d'onde le long de $x$.
\end{definition}

\begin{definition}[Nœud et Ventre]
\begin{itemize}
\item Un \cle{nœud} est un point où l'amplitude $A(x)$ est \textbf{nulle} : le point reste \textbf{permanemment au repos} ($y=0\ \forall t$).
\item Un \cle{ventre} (ou antinœud) est un point où l'amplitude $|A(x)|$ est \textbf{maximale} (égale à $2a$) : le point vibre avec la \textbf{plus grande amplitude possible}.
\end{itemize}
\end{definition}

\courssub{Localisation des nœuds et des ventres — Démonstration guidée}

À partir de l'expression admise $A(x) = 2a \sin\!\left( \dfrac{2\pi}{\lambda}x \right)$, déterminons les positions des nœuds et des ventres.

\begin{propriete}[Positions des nœuds et des ventres]
Pour une onde stationnaire sur une corde fixée en $x=0$ (vibreur) et $x=L$ (poulie) :
\begin{itemize}
\item \textbf{Nœuds :} $A(x)=0 \;\Leftrightarrow\; \sin\!\left( \frac{2\pi}{\lambda}x \right)=0 \;\Leftrightarrow\; \frac{2\pi}{\lambda}x = k\pi \;\Leftrightarrow\; \boxed{x_N = k\,\frac{\lambda}{2},\quad k\in\mathbb{N}}$.
\item \textbf{Ventres :} $|A(x)|$ maximum $\Leftrightarrow \left|\sin\!\left( \frac{2\pi}{\lambda}x \right)\right|=1 \;\Leftrightarrow\; \frac{2\pi}{\lambda}x = \frac{\pi}{2}+k\pi \;\Leftrightarrow\; \boxed{x_V = \left(k+\frac{1}{2}\right)\frac{\lambda}{2},\quad k\in\mathbb{N}}$.
\end{itemize}
\end{propriete}

\begin{aretenir}[Distances caractéristiques]
\begin{itemize}
\item Distance entre deux nœuds consécutifs : $\Delta x_N = \frac{\lambda}{2}$.
\item Distance entre deux ventres consécutifs : $\Delta x_V = \frac{\lambda}{2}$.
\item Distance entre un nœud et le ventre adjacent : $\Delta x_{NV} = \frac{\lambda}{4}$.
\end{itemize}
Ces relations sont \textbf{universelles} pour toute onde stationnaire (corde, colonne d'air, cavité électromagnétique).
\end{aretenir}

\begin{methode}[Déterminer $\lambda$ sur une corde de Melde]
\begin{enumerate}
\item Compter le nombre $n$ de fuseaux (ventres) observés sur la longueur vibrante $L$.
\item La longueur $L$ contient $n$ demi-longueurs d'onde : $\boxed{L = n\,\frac{\lambda}{2}}$.
\item En déduire $\lambda = \dfrac{2L}{n}$.
\item Connaître la fréquence $f$ imposée par le vibreur (affichée ou mesurée au fréquencemètre).
\item Calculer la célérité de l'onde sur la corde : $v = \lambda f = \dfrac{2Lf}{n}$.
\item Vérifier la loi de la corde : $v = \sqrt{\dfrac{F}{\mu}}$ où $\mu$ est la masse linéique de la corde.
\end{enumerate}
\end{methode}

\courssub{Allure temporelle : immobilité des nœuds, oscillation des ventres}

La factorisation $y(x,t)=A(x)\cos(\omega t)$ montre que \textbf{tous les points d'un même fuseau passent par leur position d'équilibre en même temps} et atteignent leur élongation maximale en même temps : ils vibrent \textbf{en phase}. Deux fuseaux voisins ont des amplitudes $A(x)$ de signes opposés (car $\sin$ change de signe entre deux zéros) : ils vibrent en \textbf{opposition de phase} (déphasage $\pi$).

\begin{popfigure}
\begin{tikzpicture}[scale=1.1]
\begin{axis}[
    width=\linewidth, height=6cm,
    axis lines=middle,
    xlabel={$x$ (m)}, ylabel={$y$ (m)},
    xmin=0, xmax=1.5, ymin=-1.2, ymax=1.2,
    xtick={0,0.25,0.5,0.75,1.0,1.25,1.5},
    xticklabels={$0$, $\lambda/4$, $\lambda/2$, $3\lambda/4$, $\lambda$, $5\lambda/4$, $3\lambda/2$},
    ytick={-1,0,1},
    yticklabels={$-2a$, $0$, $+2a$},
    domain=0:1.5, samples=300,
    legend style={at={(0.5,-0.18)}, anchor=north, legend columns=3, fill=popcream, draw=popline, font=\small},
    grid=major, grid style={popline!50},
]
% Enveloppes
\addplot[popgreen!50, thick, domain=0:1.5] {sin(deg(2*pi*x))};
\addplot[popgreen!50, thick, domain=0:1.5, forget plot] {-sin(deg(2*pi*x))};
% Instants
\addplot[popblue, very thick, domain=0:1.5] {sin(deg(2*pi*x))};
\addplot[poporange, very thick, domain=0:1.5, dashed] {0*x};
\addplot[poppurple, thick, domain=0:1.5] {0.707*sin(deg(2*pi*x))};
\addplot[poppink, thick, domain=0:1.5] {-0.707*sin(deg(2*pi*x))};
% Nœuds
\addplot[only marks, mark=*, mark size=2.5, popgreen] coordinates {(0,0) (0.5,0) (1.0,0) (1.5,0)};
\legend{Enveloppes $\pm A(x)$, $t=0$ ($y=A(x)$), $t=T/4$ ($y=0$), $t=T/8$, $t=3T/8$, Nœuds};
\end{axis}
\end{tikzpicture}
\legende{Allure de la corde à différents instants pour une onde stationnaire à 3 fuseaux ($\lambda=1,00\,\text{m}$, $L=\dfrac{3\lambda}{2}=1,50\,\text{m}$). Les nœuds (points verts) restent fixes. À $t=0$, la corde coïncide avec l'enveloppe supérieure ; à $t=T/4$, elle est plate (tous les points passent par l'équilibre simultanément).}
\end{popfigure}

\courssub{Différence fondamentale avec une onde progressive : aspect énergétique}

\begin{attention}[Piège majeur : pas de propagation d'énergie !]
Dans une \textbf{onde progressive}, l'énergie (cinétique + potentielle) est transportée par l'onde à la vitesse $v$. Dans une \textbf{onde stationnaire}, \textbf{l'énergie ne se propage pas le long de la corde}. Elle reste localisée entre deux nœuds consécutifs (à l'intérieur d'un fuseau) : elle oscille entre énergie potentielle d'élongation (corde déformée au maximum, $v=0$) et énergie cinétique (corde plate, vitesse maximale). \textbf{La puissance moyenne transportée à travers une section est nulle.}
\end{attention}

\begin{propriete}[Comportement en phase]
\begin{itemize}
\item Tous les points situés \textbf{entre deux nœuds consécutifs} (un même fuseau) vibrent \textbf{en phase} (même signe de $A(x)$).
\item Deux fuseaux voisins vibrent en \textbf{opposition de phase} (signes opposés de $A(x)$).
\item À tout instant, la corde a une forme homothétique de l'enveloppe $A(x)$ (facteur $\cos\omega t$).
\end{itemize}
\end{propriete}

\courssub{Conditions aux limites et modes propres (fréquences de résonance)}

La corde est fixée en $x=0$ (vibreur) et $x=L$ (poulie). Ces points \textbf{doivent être des nœuds} : $y(0,t)=0$ et $y(L,t)=0\ \forall t$. La condition $x=0$ est satisfaite par $\sin(2\pi x/\lambda)$. La condition $x=L$ impose :
\[
\sin\!\left( \frac{2\pi}{\lambda}L \right) = 0 \;\Longrightarrow\; \frac{2\pi}{\lambda}L = n\pi \;\Longrightarrow\; \boxed{L = n\,\frac{\lambda_n}{2},\quad n\in\mathbb{N}^*}
\]
Chaque entier $n\ge 1$ définit un \textbf{mode propre} (ou harmonique de rang $n$) de la corde.

\begin{definition}[Fréquences propres et modes de vibration]
\begin{itemize}
\item \textbf{Longueur d'onde du mode $n$ :} $\lambda_n = \dfrac{2L}{n}$.
\item \textbf{Fréquence propre (résonance) :} $f_n = \dfrac{v}{\lambda_n} = n\,\dfrac{v}{2L} = n\,f_1$.
\item \textbf{Fréquence fondamentale (premier harmonique) :} $f_1 = \dfrac{v}{2L} = \dfrac{1}{2L}\sqrt{\dfrac{F}{\mu}}$.
\item Le mode $n$ présente $n$ fuseaux (ventres) et $n+1$ nœuds (extrémités comprises).
\end{itemize}
\end{definition}

\begin{exemplebox}[Exemple chiffré : Corde de guitare — Accordage]
Une corde de guitare a une longueur vibrante $L=\SI{0,65}{m}$ et une masse linéique $\mu=\SI{4,0e-4}{kg.m^{-1}}$. Elle est tendue avec une force $F=\SI{80}{N}$.
\begin{enumerate}
\item Calculer la célérité $v$ des ondes sur cette corde.
\item Déterminer la fréquence fondamentale $f_1$.
\item Donner les fréquences des deux premiers harmoniques ($n=2,3$).
\item Si on pince la corde au tiers de sa longueur ($x=L/3$), quels harmoniques sont favorisés ?
\end{enumerate}
\textbf{Solution :}
\begin{enumerate}
\item $v = \sqrt{\dfrac{F}{\mu}} = \sqrt{\dfrac{80}{4,0\times10^{-4}}} = \sqrt{2,0\times10^5} \approx \SI{447}{m.s^{-1}}$.
\item $f_1 = \dfrac{v}{2L} = \dfrac{447}{2\times0,65} \approx \SI{344}{Hz}$ (cette note est proche du \textbf{Fa$_4$} standard à \SI{349}{Hz} ; le \textbf{La$_3$} standard (diapason) est à \SI{440}{Hz}).
\item $f_2 = 2f_1 \approx \SI{688}{Hz}$, $f_3 = 3f_1 \approx \SI{1032}{Hz}$.
\item Pincer en $x=L/3$ crée un antinœud (ventre) à cet endroit. Les modes qui ont un ventre en $x=L/3$ sont ceux pour lesquels $\sin(2\pi x/\lambda_n) \neq 0$ et de préférence maximal. $x=L/3 = n\lambda_n/6$. Pour $n=3$, $\lambda_3=2L/3$, $x=L/3 = \lambda_3/2$ : c'est un \textbf{nœud} pour $n=3$ (mode non excité). Pour $n=1$, $x=L/3 = \lambda_1/6$ : ventre. Pour $n=2$, $x=L/3 = \lambda_2/3$ : ventre. \textbf{Les modes $n=1,2,4,5,\dots$ sont excités ; les modes multiples de 3 ($n=3,6,\dots$) sont absents.}
\end{enumerate}
\end{exemplebox}

\courssub{Exercice résolu : Détermination de la masse linéique par Melde}

\begin{exoresolu}[Mesure de $\mu$ par l'expérience de Melde]
\textbf{Énoncé :} On réalise l'expérience de Melde avec une corde de longueur vibrante $L=\SI{1,20}{m}$. La tension est assurée par une masse $m=\SI{0,500}{kg}$. On observe les résonances successives en faisant varier la fréquence $f$ du vibreur. On relève les fréquences des trois premiers modes : $f_1=\SI{18,2}{Hz}$, $f_2=\SI{36,5}{Hz}$, $f_3=\SI{54,7}{Hz}$. Déterminer la masse linéique $\mu$ de la corde et son incertitude-type (données : $u(m)=\SI{0,5}{g}$, $u(L)=\SI{0,5}{cm}$, $u(f)=\SI{0,1}{Hz}$).

\tcblower
\textbf{Solution :}
\begin{enumerate}
\item \textbf{Vérification de la cohérence :} $f_2/f_1 = 36,5/18,2 \approx 2,01 \approx 2$ ; $f_3/f_1 \approx 3,01 \approx 3$. Les fréquences sont bien proportionnelles aux entiers : on est bien sur les modes propres $n=1,2,3$.
\item \textbf{Célérité $v$ :} Pour chaque mode, $v = \lambda_n f_n = \frac{2L}{n} f_n = 2L \frac{f_n}{n}$.
    \begin{align*}
    v_1 &= 2\times1,20\times18,2 = \SI{43,7}{m.s^{-1}} \\
    v_2 &= 2\times1,20\times\frac{36,5}{2} = \SI{43,8}{m.s^{-1}} \\
    v_3 &= 2\times1,20\times\frac{54,7}{3} = \SI{43,8}{m.s^{-1}}
    \end{align*}
    Moyenne : $\bar{v} = \SI{43,8}{m.s^{-1}}$.
\item \textbf{Masse linéique :} $v = \sqrt{F/\mu} \Rightarrow \mu = F/v^2 = mg/v^2$.
    $\mu = \dfrac{0,500\times9,81}{(43,8)^2} = \dfrac{4,905}{1918} \approx \SI{2,56e-3}{kg.m^{-1}}$.
\item \textbf{Incertitude-type (loi de propagation, variables indépendantes) :}
    $\mu = \dfrac{mg}{v^2}$, $v = 2L f_1$ (on utilise le mode fondamental pour minimiser l'incertitude relative sur $f$).
    $\dfrac{u(\mu)}{\mu} = \sqrt{ \left(\dfrac{u(m)}{m}\right)^2 + \left(2\dfrac{u(v)}{v}\right)^2 }$, avec $\dfrac{u(v)}{v} = \sqrt{ \left(\dfrac{u(L)}{L}\right)^2 + \left(\dfrac{u(f_1)}{f_1}\right)^2 }$.
    \begin{align*}
    \frac{u(m)}{m} &= \frac{0,0005}{0,500} = 0,001 \\
    \frac{u(L)}{L} &= \frac{0,005}{1,20} \approx 0,00417 \\
    \frac{u(f_1)}{f_1} &= \frac{0,1}{18,2} \approx 0,00549 \\
    \frac{u(v)}{v} &= \sqrt{0,00417^2 + 0,00549^2} \approx 0,0069 \\
    \frac{u(\mu)}{\mu} &= \sqrt{0,001^2 + (2\times0,0069)^2} \approx 0,0138 \\
    u(\mu) &= 0,0138 \times 2,56\times10^{-3} \approx 0,035\times10^{-3}\,\text{kg}\cdot\text{m}^{-1}
    \end{align*}
    \textbf{Résultat :} $\boxed{\mu = (2,56 \pm 0,04)\times10^{-3}\,\text{kg}\cdot\text{m}^{-1}}$.
\end{enumerate}
\end{exoresolu}

\courssub{Applications technologiques et prolongements}

\begin{savaistu}[Ondes stationnaires partout !]
\begin{itemize}
\item \textbf{Instruments à vent (flûte, clarinette, saxophone, orgue) :} La colonne d'air dans le tube vibre en ondes stationnaires. Les nœuds sont des points de pression maximale (pression) ou de déplacement nul (vitesse particulaire), selon la grandeur considérée. La longueur du tube détermine les fréquences jouées ($L=n\lambda/2$ tube ouvert, $L=(2n-1)\lambda/4$ tube fermé).
\item \textbf{Instruments traditionnels camerounais (Mvet, Ntong) :} Le Mvet (cithare sur calebasse) utilise des cordes tendues sur un résonateur ; les modes propres de la corde déterminent la gamme pentatonique. Le Ntong (tambour à fente) exploite les modes propres de vibration du bois creux.
\item \textbf{Cavités laser (optique) :} Deux miroirs parallèles forment une cavité de Fabry-Pérot. La lumière se réfléchit en aller-retour : seules les longueurs d'onde vérifiant $L=n\lambda/2$ (onde stationnaire) sont amplifiées. C'est la sélectivité spectrale du laser.
\item \textbf{Four à micro-ondes :} La cavité du four est un résonateur électromagnétique. Les ondes stationnaires créent des zones chaudes (ventres du champ électrique) et froides (nœuds) — d'où le plateau tournant pour chauffer uniformément.
\item \textbf{Échographie médicale (mode A, B) :} Bien que basée sur l'écho (onde progressive réfléchie), l'analyse des signaux stationnaires formés dans les couches minces (ex: œil) permet des mesures d'épaisseur précises.
\item \textbf{Antennes radio (dipôle $\lambda/2$) :} Le courant sur l'antenne forme une onde stationnaire. L'adaptation d'impédance (absence d'onde réfléchie vers l'émetteur) se fait quand l'antenne est accordée (longueur = multiple de $\lambda/2$). C'est le principe des antennes relais MTN/Orange/CAMTEL au Cameroun.
\end{itemize}
\end{savaistu}

\begin{exercice}[Entraînement : Corde de Melde à 4 fuseaux]
Une corde de longueur $L=\SI{1,60}{m}$, de masse linéique $\mu=\SI{2,5e-3}{kg.m^{-1}}$, est tendue par une masse $m=\SI{0,800}{kg}$. Le vibreur est réglé sur une fréquence $f$ telle que la corde vibre en $n=4$ fuseaux.
\begin{enumerate}
\item Calculer la tension $F$ et la célérité $v$ des ondes sur la corde.
\item Déterminer la fréquence $f$ du vibreur.
\item Quelle est la distance entre deux nœuds consécutifs ? Entre un nœud et le ventre adjacent ?
\item Tracer l'allure de la corde à l'instant $t=0$ (enveloppe maximale) en indiquant les positions des nœuds et ventres sur l'axe $x$ (origine au vibreur).
\item Si on double la masse $m$ (autres paramètres inchangés), quel sera le nouveau nombre de fuseaux pour la même fréquence $f$ ? Justifier.
\end{enumerate}
\end{exercice}


\begin{corrige}[Exercice : Corde de Melde à 4 fuseaux]
\begin{enumerate}
\item $F = mg = 0,800\times9,81 = \SI{7,85}{N}$. $v = \sqrt{F/\mu} = \sqrt{7,85 / 2,5\times10^{-3}} = \sqrt{3140} \approx \SI{56,0}{m.s^{-1}}$.
\item $\lambda = 2L/n = 2\times1,60/4 = \SI{0,80}{m}$. $f = v/\lambda = 56,0/0,80 = \SI{70,0}{Hz}$.
\item Distance nœud--nœud = $\lambda/2 = \SI{0,40}{m}$. Distance nœud--ventre = $\lambda/4 = \SI{0,20}{m}$.
\item Nœuds en $x = 0,\ 0,40,\ 0,80,\ 1,20,\ 1,60\,\text{m}$. Ventres en $x = 0,20,\ 0,60,\ 1,00,\ 1,40\,\text{m}$. Enveloppe : $y(x,0) = 2a\sin(2\pi x/\lambda)$.
\item Nouvelle tension $F' = 2F = \SI{15,7}{N}$. Nouvelle célérité $v' = \sqrt{2}v \approx \SI{79,2}{m.s^{-1}}$. Nouvelle longueur d'onde $\lambda' = v'/f = \sqrt{2}\lambda \approx \SI{1,13}{m}$. Nouveau nombre de fuseaux $n' = 2L/\lambda' = 2L/(\sqrt{2}\lambda) = n/\sqrt{2} = 4/\sqrt{2} \approx 2,83$. \textbf{Le nombre de fuseaux doit être un entier.} La corde ne résonnera plus exactement à cette fréquence $f$ avec la nouvelle tension ; il faudra changer $f$ pour retomber sur un mode propre (probablement $n'=3$ si on augmente légèrement $f$, ou $n'=2$ si on la diminue).
\end{enumerate}
\end{corrige}

\begin{attention}[Erreurs fréquentes au Bac — Check-list]
\begin{itemize}
\item \textbf{Confondre nœud et ventre :} Nœud = amplitude nulle (point fixe) ; Ventre = amplitude max. Sur une corde fixée aux deux bouts, les extrémités sont \textbf{toujours des nœuds}.
\item \textbf{Oublier le facteur 2 :} $L = n\lambda/2$ (et non $n\lambda$). Un fuseau = un demi-longueur d'onde.
\item \textbf{Croire que l'onde stationnaire se propage :} Non, l'enveloppe est fixe. L'énergie ne voyage pas.
\item \textbf{Mélanger phase spatiale et phase temporelle :} Deux points d'un même fuseau sont \textbf{en phase} (même signe de $A(x)$) ; deux fuseaux voisins sont en \textbf{opposition de phase} (signes opposés). Mais \textbf{tous} les points passent par l'équilibre en même temps ($\cos\omega t=0$).
\item \textbf{Unités :} $v$ en $\text{m}\cdot\text{s}^{-1}$, $f$ en $\text{Hz}$, $\lambda$ en $\text{m}$, $\mu$ en $\text{kg}\cdot\text{m}^{-1}$, $F$ en $\text{N}$. Vérifier l'homogénéité : $[v] = \sqrt{[F]/[\mu]} = \sqrt{\text{N}/(\text{kg}\cdot\text{m}^{-1})} = \sqrt{\text{kg}\cdot\text{m}\cdot\text{s}^{-2}\cdot\text{m}\cdot\text{kg}^{-1}} = \text{m}\cdot\text{s}^{-1}$. ✓
\end{itemize}
\end{attention}

\begin{methode}[Résoudre un problème d'ondes stationnaires sur corde]
\begin{enumerate}
\item \textbf{Identifier les conditions aux limites :} Corde fixée aux deux bouts $\Rightarrow$ nœuds aux extrémités. Corde libre à un bout $\Rightarrow$ ventre à cet endroit.
\item \textbf{Écrire la relation $L = n\lambda/2$ (ou $(2n-1)\lambda/4$ si une extrémité libre).}
\item \textbf{Utiliser la loi de la corde :} $v = \sqrt{F/\mu}$ avec $F=mg$ si tension par masse pendante.
\item \textbf{Relier $v$, $f$, $\lambda$ :} $v = \lambda f = 2Lf/n$.
\item \textbf{Calculer la grandeur demandée} ($f$, $m$, $L$, $\mu$, $n$...).
\item \textbf{Vérifier l'ordre de grandeur} et l'homogénéité des unités.
\end{enumerate}
\end{methode}

Cette section t'a permis de maîtriser la physique des ondes stationnaires sur corde : du dispositif de Melde à la modélisation mathématique, en passant par l'analyse énergétique et les applications concrètes (instruments de musique, lasers, antennes, four micro-ondes). La prochaine section (5/6) généralisera ces concepts aux \textbf{modes propres des instruments de musique à cordes et à vent} et à la \textbf{résonance}.

\coursec{Modes propres d'une corde et instruments de musique}

Dans la section précédente, l'expérience de Melde nous a appris qu'une corde vibrante ne prend un aspect stable et spectaculaire --- les fameux \cle{fuseaux} --- que pour certaines fréquences bien précises d'excitation. Pour toutes les autres fréquences, la corde s'agite de façon confuse et désordonnée. Pourquoi la corde est-elle si \emph{sélective} ? Pourquoi n'existe-t-il qu'une liste discrète de fréquences « autorisées » ? C'est toute la physique des \cle{modes propres}, et c'est aussi le secret de fabrication de tous les instruments de musique à cordes, de la guitare du lycéen de Yaoundé au \cle{mvet} traditionnel des Beti et des Fang. Cette section répond à ces questions avec rigueur, puis les applique concrètement.

\courssub{Conditions aux limites : la clé de la quantification}

Reprenons l'expérience de Melde : la corde est fixée d'un côté à la lame du vibreur (dont l'amplitude est si faible qu'on la considère comme immobile) et de l'autre à une poulie, où le poids suspendu impose la tension. Les deux extrémités de la corde sont donc pratiquement \textbf{fixes}.

\begin{propriete}[Conditions aux limites d'une corde fixée à ses deux extrémités]
Une corde de longueur $L$ fixée à ses deux extrémités ne peut vibrer en régime stationnaire que si chaque extrémité est un \cle{nœud} de vibration, c'est-à-dire un point dont l'élongation reste nulle à tout instant :
\[
y(0,t) = 0 \qquad \text{et} \qquad y(L,t) = 0 \quad \forall\, t.
\]
\end{propriete}

Cette contrainte, en apparence anodine, a une conséquence spectaculaire. Une onde stationnaire possède une structure périodique en espace : les nœuds sont équidistants, séparés les uns des autres de $\lambda/2$ (un fuseau mesure exactement une demi-longueur d'onde). Si les deux extrémités doivent être des nœuds, alors la longueur $L$ de la corde doit contenir un \textbf{nombre entier de fuseaux}. On ne peut pas « couper » un fuseau : la corde ne peut pas s'arrêter au milieu d'un ventre, car son extrémité serait alors en mouvement, ce qui est impossible pour une extrémité fixe.

\begin{attention}[Le piège classique : $L = n\lambda$ ou $L = n\lambda/2$ ?]
L'erreur la plus fréquente au Baccalauréat est d'écrire $L = n\lambda$. C'est \textbf{faux} : chaque fuseau mesure $\lambda/2$, pas $\lambda$. Retiens l'image : un fuseau, c'est une \emph{demi}-boucle de sinusoïde, comprise entre deux nœuds consécutifs. Une longueur d'onde complète correspond à \emph{deux} fuseaux. Donc $L = n\,\dfrac{\lambda}{2}$ avec $n$ égal au nombre de fuseaux observés.
\end{attention}

\courssub{Condition de quantification et fréquences propres}

\begin{propriete}[Théorème des modes propres d'une corde — démonstration exigible]
Pour une corde de longueur $L$ fixée à ses deux extrémités, les vibrations stationnaires stables n'existent que si
\[
L = n\,\frac{\lambda_n}{2}, \qquad n \in \mathbb{N}^{*},
\]
et les fréquences correspondantes, appelées \cle{fréquences propres}, sont
\[
f_n = n\,\frac{c}{2L} \qquad \text{avec} \qquad c = \sqrt{\frac{F}{\mu}},
\]
où $F$ est la tension de la corde (en N) et $\mu$ sa masse linéique (en kg/m).
\end{propriete}

\begin{experience}[Démonstration guidée du théorème des modes propres]
Suivons le raisonnement pas à pas.
\begin{enumerate}
\item \textbf{Condition aux limites.} Les deux extrémités étant des nœuds, la corde contient un nombre entier $n$ de fuseaux, chacun de longueur $\lambda_n/2$ :
\[
L = n\,\frac{\lambda_n}{2} \qquad (n = 1, 2, 3, \dots)
\]
\item \textbf{Relation onde-célérité.} Pour chaque mode, l'onde se propage sur la corde à la célérité $c$, liée à la longueur d'onde et à la fréquence par la relation fondamentale :
\[
c = \lambda_n f_n \quad \Longrightarrow \quad \lambda_n = \frac{c}{f_n}.
\]
\item \textbf{Substitution.} En reportant dans la condition de quantification :
\[
L = n\,\frac{c}{2f_n} \quad \Longrightarrow \quad \boxed{\,f_n = n\,\frac{c}{2L}\,}.
\]
\item \textbf{Célérité sur la corde.} L'étude mécanique d'un élément de corde (bilan des forces de tension projeté sur l'axe transversal) montre que la célérité des ondes transversales est $c = \sqrt{F/\mu}$. Vérifions l'homogénéité : $[F/\mu] = \dfrac{\mathrm{kg\,m\,s^{-2}}}{\mathrm{kg\,m^{-1}}} = \mathrm{m^2\,s^{-2}}$, donc $\sqrt{F/\mu}$ est bien en $\mathrm{m\,s^{-1}}$. \textbf{CQFD.}
\end{enumerate}
\end{experience}

\begin{definition}[Mode fondamental et harmoniques]
\begin{itemize}
\item Le mode $n = 1$, de fréquence $f_1 = \dfrac{c}{2L}$, est le \cle{mode fondamental} : la corde vibre en un seul fuseau. C'est la plus basse fréquence propre ; elle détermine la \textbf{hauteur} du son perçu.
\item Les modes $n \geq 2$, de fréquences $f_n = n f_1$, sont les \cle{harmoniques} de rang $n$ : la corde vibre alors en $n$ fuseaux.
\end{itemize}
Les fréquences propres d'une corde forment ainsi une suite arithmétique : $f_1,\ 2f_1,\ 3f_1,\ \dots$ On dit que le spectre d'une corde est \emph{harmonique}.
\end{definition}

\begin{popfigure}
\begin{center}
\begin{tikzpicture}[scale=1.0]
% Mode n=1
\draw[popmuted, dashed] (0,4.6) -- (6,4.6);
\draw[popmuted, dashed] (0,3.4) -- (6,3.4);
\draw[popdark, thick] (0,4.6) -- (0,3.4) (6,4.6) -- (6,3.4);
\draw[popblue, very thick, domain=0:6, samples=100] plot (\x, {4.0 + 0.6*sin(30*\x)});
\draw[popblue, thick, dashed, domain=0:6, samples=100] plot (\x, {4.0 - 0.6*sin(30*\x)});
\fill[popink] (0,4.0) circle (2pt) (6,4.0) circle (2pt);
\node[left, popdark] at (-0.2,4.0) {\small $n=1$ : fondamental};
\node[right, popdark] at (6.3,4.0) {\small $L = \dfrac{\lambda_1}{2}$, \ $f_1 = \dfrac{c}{2L}$};
% Mode n=2
\draw[popmuted, dashed] (0,2.6) -- (6,2.6);
\draw[popmuted, dashed] (0,1.4) -- (6,1.4);
\draw[popdark, thick] (0,2.6) -- (0,1.4) (6,2.6) -- (6,1.4);
\draw[poporange, very thick, domain=0:6, samples=150] plot (\x, {2.0 + 0.6*sin(60*\x)});
\draw[poporange, thick, dashed, domain=0:6, samples=150] plot (\x, {2.0 - 0.6*sin(60*\x)});
\fill[popink] (0,2.0) circle (2pt) (3,2.0) circle (2pt) (6,2.0) circle (2pt);
\node[left, popdark] at (-0.2,2.0) {\small $n=2$ : harmonique 2};
\node[right, popdark] at (6.3,2.0) {\small $L = 2\,\dfrac{\lambda_2}{2}$, \ $f_2 = 2f_1$};
% Mode n=3
\draw[popmuted, dashed] (0,0.6) -- (6,0.6);
\draw[popmuted, dashed] (0,-0.6) -- (6,-0.6);
\draw[popdark, thick] (0,0.6) -- (0,-0.6) (6,0.6) -- (6,-0.6);
\draw[popgreen, very thick, domain=0:6, samples=200] plot (\x, {0.0 + 0.6*sin(90*\x)});
\draw[popgreen, thick, dashed, domain=0:6, samples=200] plot (\x, {0.0 - 0.6*sin(90*\x)});
\fill[popink] (0,0.0) circle (2pt) (2,0.0) circle (2pt) (4,0.0) circle (2pt) (6,0.0) circle (2pt);
\node[left, popdark] at (-0.2,0.0) {\small $n=3$ : harmonique 3};
\node[right, popdark] at (6.3,0.0) {\small $L = 3\,\dfrac{\lambda_3}{2}$, \ $f_3 = 3f_1$};
% Longueur L
\draw[<->, popdark, thick] (0,-1.0) -- (6,-1.0) node[midway, below] {$L$};
\end{tikzpicture}
\end{center}
\legende{Les trois premiers modes propres d'une corde fixée à ses deux extrémités. Les points noirs sont les nœuds ; les positions extrêmes de la corde (traits plein et pointillé) délimitent les fuseaux. Le mode $n$ contient $n$ fuseaux : $L = n\lambda_n/2$.}
\end{popfigure}

\courssub{Méthode de résolution et exemple chiffré}

\begin{methode}[Résoudre un problème de corde vibrante en 3 étapes]
\begin{enumerate}
\item \textbf{Calculer la célérité} des ondes sur la corde : $c = \sqrt{\dfrac{F}{\mu}}$, en convertissant $\mu$ en kg/m et $F$ en N.
\item \textbf{Appliquer la condition de quantification} : $L = n\,\dfrac{\lambda_n}{2}$, où $n$ est le nombre de fuseaux observés. On en tire $\lambda_n = \dfrac{2L}{n}$.
\item \textbf{En déduire la fréquence} : $f_n = \dfrac{c}{\lambda_n} = n\,\dfrac{c}{2L} = n f_1$.
\end{enumerate}
Vérifie toujours la cohérence : plus il y a de fuseaux, plus la fréquence est élevée.
\end{methode}

\begin{exemplebox}[Corde vibrante : calcul complet]
Une corde de longueur $L = \SI{80}{cm} = \SI{0,80}{m}$ est tendue par une force $F = \SI{100}{N}$. Sa masse linéique vaut $\mu = \SI{1}{g/m} = \SI{1e-3}{kg/m}$.
\begin{enumerate}
\item Célérité : $c = \sqrt{\dfrac{F}{\mu}} = \sqrt{\dfrac{100}{10^{-3}}} = \sqrt{10^{5}} \approx \SI{316}{m/s}$.
\item Fondamental : $f_1 = \dfrac{c}{2L} = \dfrac{316}{2 \times {0,80}} \approx \SI{198}{Hz}$ (proche du \emph{sol}\textsubscript{3} des musiciens).
\item Troisième harmonique : $f_3 = 3 f_1 \approx \SI{593}{Hz}$ ; la corde vibre alors en 3 fuseaux, avec $\lambda_3 = \dfrac{2L}{3} \approx \SI{0,53}{m}$.
\end{enumerate}
\end{exemplebox}

\begin{exoresolu}[Accorder une corde de mvet]
Une corde d'un instrument à cordes mesure $L = \SI{60}{cm}$ et a une masse linéique $\mu = \SI{2,5}{g/m}$. Le musicien veut qu'elle émette un fondamental de fréquence $f_1 = \SI{220}{Hz}$ (le \emph{la}\textsubscript{3}). Quelle tension $F$ doit-il appliquer à la corde ?
\tcblower
\textbf{Solution détaillée.}
\begin{enumerate}
\item La fréquence fondamentale est $f_1 = \dfrac{c}{2L}$, donc la célérité requise est :
\[
c = 2L f_1 = 2 \times {0,60} \times 220 = \SI{264}{m/s}.
\]
\item Or $c = \sqrt{\dfrac{F}{\mu}}$, donc $F = \mu c^{2}$. Avec $\mu = \SI{2,5e-3}{kg/m}$ :
\[
F = {2,5} \times 10^{-3} \times (264)^{2} = {2,5} \times 10^{-3} \times 69\,696 \approx \SI{174}{N}.
\]
\item \textbf{Vérification dimensionnelle} : $[\mu c^2] = \mathrm{kg\,m^{-1}} \times \mathrm{m^2\,s^{-2}} = \mathrm{kg\,m\,s^{-2}} = \mathrm{N}$. Correct.
\end{enumerate}
Le musicien doit tendre la corde avec une force d'environ $\SI{174}{N}$, soit l'équivalent du poids d'une masse d'environ $\SI{17,7}{kg}$.
\end{exoresolu}

\courssub{Influence des paramètres : comment le musicien joue}

La formule $f_n = \dfrac{n}{2L}\sqrt{\dfrac{F}{\mu}}$ est la « notice d'emploi » de tout instrument à cordes. Elle montre que la hauteur du son dépend de trois paramètres que l'instrumentiste peut modifier.

\textbf{1. La longueur vibrante $L$.} En appuyant un doigt sur la touche d'une guitare, le musicien raccourcit la longueur de corde libre de vibrer. Comme $f_1 = c/(2L)$, diminuer $L$ \textbf{augmente} la fréquence : le son devient plus aigu. C'est le geste de base du guitariste, du harpiste et du joueur de mvet.

\textbf{2. La tension $F$.} En tournant une cheville, on tend davantage la corde : $F$ augmente, donc $c = \sqrt{F/\mu}$ augmente, et le son monte dans les aigus. C'est ainsi qu'on \emph{accorde} un instrument. Attention : la fréquence varie comme $\sqrt{F}$, il faut donc \emph{quadrupler} la tension pour doubler la fréquence.

\textbf{3. La masse linéique $\mu$.} Les cordes graves d'une guitare ou d'une harpe sont plus grosses (souvent filées) : $\mu$ est grand, donc $c$ est faible et le son est grave. Les cordes aiguës sont fines : $\mu$ petit, fréquence élevée.

\begin{popfigure}
\begin{center}
\begin{tikzpicture}[scale=1.0]
% Manche de guitare stylisé
\draw[fill=popgoldL, draw=popdark, thick] (0,-0.35) rectangle (8,0.35);
\foreach \x in {2,4,5.6,6.9} \draw[popdark, thick] (\x,-0.35) -- (\x,0.35);
% Corde entière
\draw[popblue, very thick] (0,0.55) -- (8,0.55);
\fill[popink] (0,0.55) circle (2pt) (8,0.55) circle (2pt);
\draw[<->, popblue, thick] (0,0.95) -- (8,0.95) node[midway, above, popblue] {$L$ : corde entière, son grave, $f_1 = \dfrac{c}{2L}$};
% Corde raccourcie
\draw[poporange, very thick] (4,-0.75) -- (8,-0.75);
\fill[popink] (4,-0.75) circle (2pt) (8,-0.75) circle (2pt);
\draw[popdark, thick] (4,-1.05) -- (4,-0.45);
\node[above, popdark] at (4,-0.42) {\small doigt};
\draw[<->, poporange, thick] (4,-1.35) -- (8,-1.35) node[midway, below, poporange] {$L' = \dfrac{L}{2}$ : fréquence doublée, $f'_1 = 2f_1$};
\end{tikzpicture}
\end{center}
\legende{Effet du raccourcissement de la longueur vibrante sur une corde de guitare ou de mvet : en bloquant la corde au milieu du manche ($L' = L/2$), le musicien double la fréquence fondamentale --- il obtient la même note une octave plus haut.}
\end{popfigure}

\begin{exemplebox}[Le doigt qui double la fréquence]
Sur la corde de l'exemple précédent ($f_1 \approx \SI{198}{Hz}$ pour $L = \SI{80}{cm}$), le musicien appuie son doigt de sorte que la longueur vibrante passe à $L' = \SI{40}{cm}$. La tension et la masse linéique sont inchangées, donc $c \approx \SI{316}{m/s}$ :
\[
f'_1 = \frac{c}{2L'} = \frac{316}{2 \times {0,40}} \approx \SI{395}{Hz} \approx 2 f_1.
\]
Diviser la longueur par deux double la fréquence : le son monte d'une octave. C'est exactement ce que fait un guitariste en jouant à la 12\textsuperscript{e} case.
\end{exemplebox}

\begin{attention}[Deux confusions à bannir]
\begin{itemize}
\item Ne confonds pas la \textbf{fréquence propre} (qui dépend de $L$, $F$, $\mu$) avec la \textbf{fréquence d'excitation} du vibreur : en régime stationnaire, la corde vibre à la fréquence imposée, mais les fuseaux n'apparaissent que si cette fréquence coïncide avec une fréquence propre (phénomène de \cle{résonance}).
\item La formule $f_n = \dfrac{nc}{2L}$ suppose une corde fixée aux \emph{deux} bouts. Pour un tuyau sonore ou une corde libre à une extrémité, les conditions aux limites changent et la quantification aussi.
\end{itemize}
\end{attention}

\courssub{Complément Bac : du fondamental au timbre}

\begin{aretenir}[Timbre d'un instrument — complément exigible]
Lorsqu'on pince une corde de guitare ou de mvet, elle ne vibre jamais dans un seul mode : la vibration réelle est la \textbf{superposition} du fondamental et de plusieurs harmoniques,
\[
y(x,t) = \sum_{n} A_n \sin\!\left(\frac{n\pi x}{L}\right) \cos(2\pi f_n t + \varphi_n).
\]
Le fondamental fixe la \textbf{hauteur} du son ; les amplitudes relatives $A_n$ des harmoniques définissent le \cle{timbre}, qui permet de distinguer une guitare d'un mvet jouant la même note. Pincer la corde \emph{en son milieu} favorise le fondamental (son rond et doux) ; la pincer près du chevalet enrichit les harmoniques (son brillant et métallique).
\end{aretenir}

\begin{savaistu}[Le mvet, instrument savant des Beti et des Fang]
Le \cle{mvet}, harpe-cithare traditionnelle du Sud-Cameroun, illustre parfaitement cette leçon. Le musicien (le \emph{mbomo mvet}) accorde ses cordes de raphia ou de fil métallique en ajustant leur \textbf{tension} et leur \textbf{longueur vibrante} : il applique sans le savoir la formule $f_1 = \dfrac{1}{2L}\sqrt{\dfrac{F}{\mu}}$ ! Les cordes les plus épaisses, de grande masse linéique, produisent les sons graves qui accompagnent le récit épique, tandis que les cordes fines et tendues donnent les notes aiguës. La physique des modes propres était ainsi maîtrisée empiriquement bien avant d'être écrite dans les manuels.
\end{savaistu}

\begin{exercice}[Pour t'entraîner]
Une corde de guitare de longueur $L = \SI{65}{cm}$, de masse linéique $\mu = \SI{1,0}{g/m}$, émet un fondamental de fréquence $f_1 = \SI{330}{Hz}$ (note \emph{mi}\textsubscript{4}).
\begin{enumerate}
\item Calculer la célérité $c$ des ondes sur cette corde, puis la tension $F$.
\item On observe, en éclairage stroboscopique, la corde vibrant en 4 fuseaux. Quelle est la fréquence de vibration ? Quel est le rang de l'harmonique ?
\item Le guitariste appuie sur la touche de sorte que $L' = \SI{48,75}{cm}$. Quelle est la nouvelle fréquence fondamentale ?
\end{enumerate}
\end{exercice}

\begin{corrige}[Exercice : corde de guitare]
\begin{enumerate}
\item $c = 2Lf_1 = 2 \times {0,65} \times 330 = \SI{429}{m/s}$. Puis $F = \mu c^2 = 10^{-3} \times (429)^2 \approx \SI{184}{N}$.
\item 4 fuseaux $\Rightarrow n = 4$ (harmonique de rang 4) : $f_4 = 4f_1 = \SI{1320}{Hz}$.
\item $f'_1 = \dfrac{c}{2L'} = \dfrac{429}{2 \times {0,4875}} \approx \SI{440}{Hz}$ : c'est le \emph{la}\textsubscript{3} du diapason !
\end{enumerate}
\end{corrige}

\begin{aretenir}[L'essentiel de la section]
\begin{itemize}
\item Corde fixée aux deux bouts $\Rightarrow$ extrémités = nœuds $\Rightarrow$ $L = n\dfrac{\lambda_n}{2}$, $n \in \mathbb{N}^{*}$.
\item Fréquences propres : $f_n = n\dfrac{c}{2L} = n f_1$, avec $c = \sqrt{\dfrac{F}{\mu}}$.
\item Fondamental ($n=1$, un fuseau) = hauteur du son ; harmoniques ($f_n = nf_1$) = timbre.
\item Le musicien joue sur $L$ (doigt sur la touche), $F$ (accordage) et $\mu$ (choix de la corde).
\end{itemize}
\end{aretenir}

\coursec{Systèmes technologiques à ondes}

Dans la section précédente, nous avons vu comment une corde de guitare ou de \emph{mvet} transforme la superposition d'ondes en modes propres, c'est-à-dire en sons musicaux bien définis. Mais les ondes ne servent pas qu'à faire vibrer l'air et émouvoir nos oreilles. Depuis un siècle, l'humanité a appris à les utiliser comme de véritables \cle{messagers} : on envoie une onde vers un obstacle, elle se réfléchit, revient, et son temps de voyage nous renseigne sur la distance, la forme ou la nature de cet obstacle. C'est le principe de l'\cle{écho}, que la nature exploite depuis des millions d'années (chauves-souris, dauphins) et que la technologie a décliné en quatre grandes familles de systèmes : l'\cle{échographie}, la \cle{télémétrie}, le \cle{sonar} et le \cle{radar}. Cette dernière section te donne les clés physiques de ces dispositifs, très prisés au Baccalauréat.

\courssub{Le principe général : l'écho comme messager}

\textbf{Une situation concrète pour commencer.}\par
Tu es au pied d'une falaise dans les monts Bamboutos. Tu cries « Oh ! » et, quelques instants plus tard, tu entends ton propre « Oh ! » te revenir. Entre l'émission et la réception, le son a fait un aller-retour jusqu'à la falaise. Si tu mesures la durée $\Delta t$ de ce voyage et si tu connais la célérité $c$ du son dans l'air, tu peux en déduire la distance $d$ de la falaise \emph{sans jamais t'en approcher}. C'est exactement ce que font, de manière industrialisée et précise, tous les systèmes à ondes que nous allons étudier.

\begin{definition}[Principe de la mesure par écho]
Un système de mesure par écho comporte trois étapes :
\begin{enumerate}
\item \textbf{Émission} : un émetteur envoie une salve (impulsion brève) d'onde dans le milieu ;
\item \textbf{Réflexion} : l'onde rencontre un obstacle (ou une interface entre deux milieux) et une partie de son énergie est \cle{réfléchie} vers la source ;
\item \textbf{Réception et mesure} : un récepteur (souvent confondu avec l'émetteur) capte l'écho et un chronomètre électronique mesure la durée $\Delta t$ de l'\textbf{aller-retour}.
\end{enumerate}
La connaissance de la célérité $c$ de l'onde dans le milieu permet alors de calculer la distance $d$ de l'obstacle.
\end{definition}

\begin{propriete}[Relation fondamentale de la télémétrie]
Soit une onde de célérité $c$ dans le milieu considéré, émise vers un obstacle situé à la distance $d$. L'écho revient au récepteur après une durée $\Delta t$. Alors :
\[ d = \frac{c\,\Delta t}{2} \]
\textbf{Démonstration (exigible).} Pendant la durée $\Delta t$, l'onde parcourt le trajet \emph{aller} (distance $d$) puis le trajet \emph{retour} (distance $d$), soit une distance totale $2d$. À célérité constante, la distance parcourue vaut $c\,\Delta t$. Donc :
\[ 2d = c\,\Delta t \quad\Longrightarrow\quad d = \frac{c\,\Delta t}{2}. \]
\textit{Vérification dimensionnelle :} $[c\,\Delta t] = \mathrm{m\,s^{-1}}\times\mathrm{s} = \mathrm{m}$, homogène à une longueur. La relation est cohérente.
\end{propriete}

\begin{attention}[Le piège du facteur 2 — erreur classique au Bac]
La durée $\Delta t$ mesurée correspond à un \textbf{aller-retour}, pas à un aller simple ! Oublier de diviser par 2 double la distance calculée : c'est l'erreur la plus fréquente dans les exercices de télémétrie. Astuce de rédaction : écris toujours d'abord $2d = c\,\Delta t$ sur ta copie, puis isole $d$. De même, si l'énoncé donne la distance et demande $\Delta t$, n'oublie pas que $\Delta t = \dfrac{2d}{c}$.
\end{attention}

\begin{popfigure}
\begin{center}
\begin{tikzpicture}[scale=1.0, >=Stealth]
% Emetteur/recepteur
\fill[popblue] (0,0) rectangle (0.8,1.6);
\node[below, popblue, font=\small] at (0.4,-0.05) {Émetteur--récepteur};
% Obstacle
\fill[poppurple!60] (7,-0.6) rectangle (7.5,2.2);
\node[below, poppurple, font=\small] at (7.25,-0.65) {Obstacle};
% Distance d
\draw[<->, popdark, thick] (0.8,-1.3) -- (7,-1.3) node[midway, below, font=\small]{$d$};
% Aller
\draw[->, poporange, very thick, decorate, decoration={snake, amplitude=1.2pt, segment length=6pt}] (0.9,1.2) -- (6.9,1.2);
\node[above, poporange, font=\small] at (3.9,1.25) {onde incidente (aller), célérité $c$};
% Retour
\draw[->, popgreen, very thick, decorate, decoration={snake, amplitude=1.2pt, segment length=6pt}] (6.9,0.3) -- (0.9,0.3);
\node[below, popgreen, font=\small] at (3.9,0.22) {écho (retour)};
% Formule
\node[draw=popdark, rounded corners, fill=popcream, font=\small] at (3.9,2.6) {$2d = c\,\Delta t \;\Rightarrow\; d = \dfrac{c\,\Delta t}{2}$};
\end{tikzpicture}
\end{center}
\legende{Principe de la mesure par écho : l'onde effectue un aller-retour de longueur totale $2d$ pendant la durée $\Delta t$.}
\end{popfigure}

\begin{methode}[Résoudre un problème de télémétrie en quatre étapes]
\begin{enumerate}
\item \textbf{Identifier la nature de l'onde} (son, ultrason, lumière, micro-onde) et en déduire la \textbf{célérité $c$ dans le milieu} : environ $\SI{340}{m.s^{-1}}$ pour le son dans l'air, $\SI{1500}{m.s^{-1}}$ pour les ultrasons dans l'eau de mer, $c = \SI{3,0e8}{m.s^{-1}}$ pour les ondes électromagnétiques dans l'air ou le vide.
\item \textbf{Relever la durée $\Delta t$} de l'aller-retour (attention à convertir en secondes : ms, \textmu s...).
\item \textbf{Écrire la relation} $2d = c\,\Delta t$, puis $d = \dfrac{c\,\Delta t}{2}$.
\item \textbf{Calculer et conclure} avec l'unité SI et un nombre raisonnable de chiffres significatifs ; vérifier l'ordre de grandeur (un fond marin de $\SI{3}{km}$ pour un écho de $\SI{4}{s}$, c'est plausible ; $\SI{300}{m}$ pour un écho de $\SI{0,4}{s}$ aussi).
\end{enumerate}
\end{methode}

\courssub{L'échographie médicale : voir à l'intérieur du corps sans l'ouvrir}

\textbf{Le choix des ultrasons.}\par
Les \cle{ultrasons} sont des ondes mécaniques de fréquence supérieure à $\SI{20}{kHz}$, inaudibles pour l'oreille humaine. En échographie médicale, on utilise typiquement des fréquences de $2$ à $\SI{15}{MHz}$. Pourquoi de si hautes fréquences ? Parce que la \cle{résolution} — la capacité à distinguer deux points proches — est d'autant meilleure que la longueur d'onde $\lambda = c/f$ est petite. Avec $c \approx \SI{1500}{m.s^{-1}}$ dans les tissus mous et $f = \SI{5}{MHz}$, on obtient $\lambda = \dfrac{1500}{5\times 10^{6}} = \SI{3e-4}{m} = \SI{0,3}{mm}$ : on peut donc « voir » des détails millimétriques.

\textbf{Principe de formation de l'image.}\par
La sonde de l'échographe, posée sur la peau avec un gel (qui assure le contact acoustique, car les ultrasons ne traversent pas l'air), joue le rôle d'émetteur \emph{et} de récepteur. Elle envoie de brèves salves d'ultrasons. À chaque \cle{interface} entre deux tissus de nature différente (peau/graisse, muscle/organe, liquide amniotique/fœtus), une partie de l'onde est réfléchie : c'est un écho. L'appareil mesure pour chaque écho :
\begin{itemize}
\item sa \textbf{durée d'aller-retour} $\Delta t$, qui donne la profondeur de l'interface par $d = \dfrac{c\,\Delta t}{2}$ ;
\item son \textbf{amplitude}, qui renseigne sur la nature des tissus (un os réfléchit beaucoup plus qu'un muscle).
\end{itemize}
En balayant la sonde et en traitant des milliers d'échos par seconde, l'ordinateur construit une image en coupe de l'intérieur du corps en temps réel.

\begin{savaistu}[L'échographie obstétricale au Cameroun]
Dans les hôpitaux de Yaoundé, Douala ou Bafoussam, l'échographie obstétricale est devenue un examen de routine du suivi de grossesse : elle permet de vérifier la croissance du fœtus, sa position, le rythme de son cœur, sans aucun danger connu — contrairement aux rayons X, les ultrasons sont des ondes mécaniques non ionisantes. C'est l'un des plus beaux succès de la physique des ondes mise au service de la santé publique.
\end{savaistu}

\courssub{La télémétrie : mesurer les distances sans mètre}

La \cle{télémétrie} est la mesure de distance par écho. Selon la portée visée, on choisit le type d'onde :

\begin{itemize}
\item \textbf{Télémètre à ultrasons} : portée de quelques mètres ; utilisé en robotique, dans les radars de recul des voitures ou pour mesurer des distances en travaux pratiques. Avec $c \approx \SI{340}{m.s^{-1}}$ dans l'air, un écho revenant après $\Delta t = \SI{20}{ms}$ signale un obstacle à $d = \dfrac{340\times 0{,}020}{2} = \SI{3,4}{m}$.
\item \textbf{Télémètre laser} : on utilise une impulsion lumineuse (onde électromagnétique, $c = \SI{3,0e8}{m.s^{-1}}$). Les géomètres qui bornent les parcelles à Douala ou les ouvriers des chantiers de construction l'utilisent quotidiennement. Comme la lumière va très vite, les durées à mesurer sont infimes : pour $d = \SI{150}{m}$, $\Delta t = \dfrac{2\times 150}{3\times 10^{8}} = \SI{1,0e-6}{s}$, soit une microseconde ! D'où la nécessité d'une électronique très rapide.
\end{itemize}

\courssub{Le sonar : l'oreille des navires}

Le \cle{sonar} (acronyme de \emph{SOund Navigation And Ranging}) applique le principe de l'écho sous l'eau, où les ondes électromagnétiques s'atténuent très vite mais où les ultrasons se propagent excellemment ($c \approx \SI{1500}{m.s^{-1}}$ dans l'eau de mer). Le navire émet une salve d'ultrasons vers le fond ; l'écho revient après réflexion sur le fond marin ou sur un banc de poissons (la vessie natatoire des poissons, remplie de gaz, réfléchit fortement les ultrasons).

Applications concrètes :
\begin{itemize}
\item \textbf{Sondage des fonds} : mesure de la profondeur sous la quille, cartographie des fonds marins ;
\item \textbf{Pêche} : les chalutiers au large de Kribi et de Limbé repèrent les bancs de sardinelles grâce à leurs sonars, ce qui optimise les campagnes de pêche ;
\item \textbf{Navigation et sécurité} : détection d'obstacles sous-marins, épaves, sous-marins.
\end{itemize}

\begin{exemplebox}[Sondage d'un fond marin au large de Kribi]
Un chalutier émet une salve ultrasonore verticalement vers le fond. L'écho est reçu après $\Delta t = \SI{0,4}{s}$. On donne $c = \SI{1500}{m.s^{-1}}$ dans l'eau de mer. Quelle est la profondeur du fond ?
\begin{align*}
2d &= c\,\Delta t = 1500 \times 0{,}4 = \SI{600}{m} \\
d &= \frac{c\,\Delta t}{2} = \frac{600}{2} = \SI{300}{m}.
\end{align*}
Le fond marin se trouve à $\SI{300}{m}$ sous le navire. (Si l'on avait oublié le facteur 2, on aurait annoncé $\SI{600}{m}$ : le chalutier croirait pêcher au-dessus d'un abîme !)
\end{exemplebox}

\courssub{Le radar : surveiller le ciel et la route}

Le \cle{radar} (\emph{RAdio Detection And Ranging}) utilise des \cle{ondes électromagnétiques} de la famille des micro-ondes (fréquences de l'ordre du gigahertz, longueurs d'onde de quelques centimètres). Ces ondes se propagent dans l'air et même dans le vide à la célérité $c = \SI{3,0e8}{m.s^{-1}}$, et sont réfléchies par les objets métalliques : avions, navires, véhicules.

\begin{itemize}
\item \textbf{Contrôle aérien} : les radars de l'aéroport international de Douala suivent la position des avions en approche, à des distances de plusieurs dizaines de kilomètres. Pour un avion à $d = \SI{60}{km}$, l'écho revient après $\Delta t = \dfrac{2\times 60\times 10^{3}}{3\times 10^{8}} = \SI{4e-4}{s}$, soit $\SI{0,4}{ms}$.
\item \textbf{Contrôle de vitesse routière} : les radars de gendarmerie sur l'axe Yaoundé--Douala mesurent la distance du véhicule à deux instants très proches ; la variation de distance divisée par l'intervalle de temps donne la vitesse (effet Doppler en réalité, mais le principe de l'écho reste la base).
\item \textbf{Météorologie} : les radars météorologiques détectent les gouttes de pluie et suivent les zones d'orage.
\end{itemize}

\begin{exoresolu}[Radar de contrôle routier]
Un radar fixe envoie une impulsion micro-onde vers une voiture et reçoit l'écho après $\Delta t = \SI{2,0}{\micro s}$. On donne $c = \SI{3,0e8}{m.s^{-1}}$.
\begin{enumerate}
\item À quelle distance du radar se trouve la voiture ?
\item Une seconde mesure, effectuée $\SI{0,50}{s}$ plus tard, donne une distance de $\SI{285}{m}$. Le véhicule respecte-t-il la limitation à $\SI{90}{km/h}$ ?
\end{enumerate}
\tcblower
\textbf{1.} Distance de la voiture :
\[ d = \frac{c\,\Delta t}{2} = \frac{3{,}0\times 10^{8}\times 2{,}0\times 10^{-6}}{2} = \SI{300}{m}. \]
\textbf{2.} La voiture s'est rapprochée de $\Delta d = 300 - 285 = \SI{15}{m}$ en $\SI{0,50}{s}$, donc :
\[ v = \frac{\Delta d}{\Delta t'} = \frac{15}{0{,}50} = \SI{30}{m.s^{-1}} = 30\times 3{,}6 = \SI{108}{km/h}. \]
Le véhicule roule à $\SI{108}{km/h}$ : il dépasse la limitation de $\SI{90}{km/h}$. Verdict : amende !
\end{exoresolu}

\courssub{Comparaison des quatre systèmes : le bon choix d'onde}

\begin{popfigure}
\begin{center}
\begin{tikzpicture}[font=\small]
% Quatre blocs
\node[draw=popblue, fill=popblueL, rounded corners, minimum width=3.1cm, minimum height=2.3cm, align=center] (echo) at (0,0) {\textbf{Échographie}\\[2pt] ultrasons\\ (mécaniques)\\ milieu : tissus\\ $c \approx 1500$ m/s};
\node[draw=popgreen, fill=popgreenL, rounded corners, minimum width=3.1cm, minimum height=2.3cm, align=center] (tele) at (4.1,0) {\textbf{Télémètre}\\[2pt] lumière laser ou\\ ultrasons\\ $c = 3\times10^{8}$ m/s\\ ou $340$ m/s};
\node[draw=poporange, fill=poporangeL, rounded corners, minimum width=3.1cm, minimum height=2.3cm, align=center] (sonar) at (0,-3.3) {\textbf{Sonar}\\[2pt] ultrasons\\ (mécaniques)\\ milieu : eau de mer\\ $c \approx 1500$ m/s};
\node[draw=poppurple, fill=poppurpleL, rounded corners, minimum width=3.1cm, minimum height=2.3cm, align=center] (radar) at (4.1,-3.3) {\textbf{Radar}\\[2pt] micro-ondes\\ (électromagnétiques)\\ air ou vide\\ $c = 3\times10^{8}$ m/s};
% Principe commun
\node[draw=popdark, fill=popcream, rounded corners, align=center, font=\small] (principe) at (2.05,2.2) {Principe commun : $2d = c\,\Delta t$};
\draw[->, popdark, thick] (principe) -- (echo);
\draw[->, popdark, thick] (principe) -- (tele);
\draw[->, popdark, thick] (principe.south west) to[bend right=10] (sonar.north west);
\draw[->, popdark, thick] (principe.south east) to[bend left=10] (radar.north east);
\end{tikzpicture}
\end{center}
\legende{Les quatre grands systèmes à ondes reposent sur le même principe de l'écho ; seuls changent la nature de l'onde et le milieu de propagation.}
\end{popfigure}

Le choix du type d'onde n'est jamais anodin ; il obéit à une logique physique claire :

\begin{itemize}
\item \textbf{Ondes mécaniques (sons, ultrasons)} : elles ont \emph{besoin d'un milieu matériel} pour se propager. Elles sont parfaites dans l'eau (sonar) ou dans les tissus biologiques (échographie), mais inutilisables dans le vide.
\item \textbf{Ondes électromagnétiques (lumière laser, micro-ondes)} : elles se propagent dans le vide comme dans l'air, à $c = \SI{3,0e8}{m.s^{-1}}$. Idéales pour les grandes distances dans l'atmosphère (radar) ou l'espace. En revanche, elles pénètrent mal l'eau de mer : c'est pourquoi aucun sous-marin n'utilise de radar sous l'eau.
\end{itemize}

\begin{aretenir}[L'essentiel des systèmes à ondes]
\begin{itemize}
\item Tous les systèmes à écho reposent sur la même relation : $d = \dfrac{c\,\Delta t}{2}$, où $\Delta t$ est la durée de l'\textbf{aller-retour}.
\item Échographie : ultrasons, interfaces des organes, image médicale non invasive.
\item Télémétrie : mesure de distances par écho lumineux (laser) ou ultrasonore.
\item Sonar : ultrasons sous l'eau ($c \approx \SI{1500}{m.s^{-1}}$), sondage et pêche.
\item Radar : micro-ondes ($c = \SI{3,0e8}{m.s^{-1}}$), contrôle aérien et routier.
\item Les ondes mécaniques exigent un milieu ; les ondes électromagnétiques se propagent dans le vide.
\end{itemize}
\end{aretenir}

\courssub{Limites et précautions d'usage}

\textbf{La résolution est limitée par la longueur d'onde.}\par
Un système à ondes ne peut pas détecter correctement des détails beaucoup plus petits que la longueur d'onde utilisée : un obstacle de taille inférieure à $\lambda$ diffracte l'onde au lieu de la réfléchir proprement. C'est pourquoi l'échographie monte en fréquence (donc descend en longueur d'onde) pour observer de fins détails, et pourquoi un radar à micro-ondes de $\lambda \approx \SI{3}{cm}$ ne peut pas « voir » un insecte.

\textbf{L'atténuation limite la portée.}\par
En se propageant, l'onde perd de l'énergie (absorption par le milieu, étalement géométrique de la salve) : l'écho revient de plus en plus faible. Au-delà d'une certaine distance, il devient indétectable. Dans les tissus biologiques, l'atténuation augmente avec la fréquence : il faut donc \emph{compromettre} entre résolution (haute fréquence) et profondeur d'exploration (basse fréquence). C'est pourquoi l'échographie cardiaque (profonde) utilise des fréquences plus basses que l'échographie superficielle de la thyroïde.

\begin{savaistu}[Les pionniers de l'écholocation : la chauve-souris et le dauphin]
Bien avant les ingénieurs, la nature a inventé le sonar ! La chauve-souris émet des ultrasons (jusqu'à $\SI{100}{kHz}$) par la bouche ou le nez et analyse les échos pour cartographier son environnement dans l'obscurité totale : elle localise un insecte au vol avec une précision stupéfiante. Le dauphin fait de même sous l'eau grâce à un organe frontal, le \emph{melon}, qui focalise ses clics ultrasonores. Le mot « écholocation » désigne ce sonar biologique. Les physiciens n'ont fait, au $XX^{\text{ème}}$ siècle, qu'imiter et industrialiser un principe vieux de dizaines de millions d'années.
\end{savaistu}

\begin{exercice}[Échographie et télémètre laser]
\textbf{1.} Lors d'une échographie, un écho revient à la sonde après $\Delta t = \SI{80}{\micro s}$. On prend $c = \SI{1500}{m.s^{-1}}$ dans les tissus. À quelle profondeur se trouve l'interface réfléchissante ?

\textbf{2.} Un télémètre laser mesure la distance d'un mur situé à $d = \SI{45}{m}$. Quelle est la durée de l'aller-retour de l'impulsion lumineuse ? ($c = \SI{3,0e8}{m.s^{-1}}$)

\textbf{3.} Pourquoi un sonar ne peut-il pas utiliser des ondes radio comme un radar ?
\end{exercice}

\begin{corrige}[Exercice : Échographie et télémètre laser]
\textbf{1.} $d = \dfrac{c\,\Delta t}{2} = \dfrac{1500\times 80\times 10^{-6}}{2} = \dfrac{0{,}12}{2} = \SI{0,060}{m} = \SI{6,0}{cm}$. L'interface se trouve à $\SI{6,0}{cm}$ de profondeur.

\textbf{2.} $\Delta t = \dfrac{2d}{c} = \dfrac{2\times 45}{3{,}0\times 10^{8}} = \SI{3,0e-7}{s}$, soit $\SI{0,30}{\micro s}$.

\textbf{3.} Les ondes radio (électromagnétiques) sont très fortement absorbées par l'eau de mer, qui est conductrice : elles ne s'y propagent que sur quelques mètres. Les ultrasons, ondes mécaniques, s'y propagent au contraire sur des kilomètres. Le choix de l'onde dépend donc du milieu.
\end{corrige}

Ainsi s'achève notre étude de la superposition des ondes mécaniques : des interférences sur la cuve à ondes aux modes propres de la corde de Melde, puis aux sonars des chalutiers de Kribi et aux radars de l'aéroport de Douala, un même fil conducteur relie tous ces phénomènes — une onde émise transporte de l'énergie et de l'\emph{information}, et sa réflexion ou sa superposition avec d'autres ondes nous permet de mesurer, d'imager et de comprendre le monde. Retiens la relation $d = \dfrac{c\,\Delta t}{2}$, garde-toi du piège du facteur 2, et tu seras prêt pour toutes les questions de télémétrie du Baccalauréat !

\newpage

\coursec{Activité d'intégration}

\coursec{Activité d'intégration : Télémétrie par sonar et ondes stationnaires sur corde de \emph{mvet}}

\courssub{Objectifs de l'activité}
Cette activité d'intégration vise à mobiliser l'ensemble des savoir-faire de la leçon sur la superposition des ondes mécaniques dans une situation authentique ancrée au Cameroun. Elle permet de :
\begin{itemize}
    \item Appliquer la relation fondamentale de la télémétrie par écho (sonar) pour déterminer une distance et une longueur d'onde.
    \item Modéliser un système vibratoire réel (corde de \emph{mvet}) par l'expérience de Melde : calculer la célérité des ondes transverses, déterminer les fréquences propres et représenter graphiquement un mode stationnaire.
    \item Établir le lien conceptuel unificateur entre ces deux situations : le rôle essentiel de la \textbf{réflexion des ondes} dans la formation d'ondes stationnaires (corde) et dans la détection par écho (sonar).
\end{itemize}

\courssub{Situation}
Au port en eau profonde de \textbf{Kribi}, un bateau de pêche artisanale utilise un sonar pour localiser des bancs de poissons. Le sonar émet des impulsions d'ultrasons de fréquence $f = \SI{50}{kHz}$ qui se propagent dans l'eau de mer à la célérité $v_{\text{eau}} = \SI{1500}{m.s^{-1}}$. Un écho revient vers le transducteur après un intervalle de temps $\Delta t = \SI{0,12}{s}$.

En classe de Terminale C, pour modéliser le principe de résonance et de modes propres mis en jeu dans les capteurs piézoélectriques du sonar, l'enseignant propose l'expérience de \textbf{Melde} avec une corde de \emph{mvet} (instrument traditionnel à cordes pincées). La corde a une longueur $L = \SI{1,20}{m}$, une masse linéique $\mu = \SI{2}{g.m^{-1}}$ et est tendue par une force de tension $F = \SI{50}{N}$ (masse suspendue d'environ $\SI{5,1}{kg}$). Le vibreur est réglé sur une fréquence correspondant à un mode propre de la corde.

\courssub{Consignes}

\textbf{Partie 1 : Télémétrie par sonar au large de Kribi}
\begin{enumerate}
    \item Rappeler la relation liant la distance $d$ parcourue par l'onde (aller-retour), la célérité $v$ et le temps de propagation $\Delta t$. En déduire la profondeur $h$ du banc de poissons.
    \item Calculer la longueur d'onde $\lambda_{\text{eau}}$ des ultrasons dans l'eau. Vérifier l'homogénéité dimensionnelle du résultat.
\end{enumerate}

\textbf{Partie 2 : Modélisation par ondes stationnaires sur une corde (Expérience de Melde)}
\begin{enumerate}
    \item Calculer la célérité $v_{\text{corde}}$ des ondes transverses progressives sur la corde tendue.
    \item Déterminer les fréquences $f_1$, $f_2$, $f_3$ des trois premiers modes propres (harmoniques) de la corde fixe aux deux extrémités.
    \item Représenter sur papier millimétré (échelle suggérée : \SI{1}{cm} pour \SI{0,10}{m}) l'allure de la corde pour le mode $n=3$ (trois ventres). Placer précisément les nœuds et les ventres en indiquant leurs abscisses exactes (en m) par rapport à une extrémité fixe.
    \item Vérifier numériquement que la condition $L = 3\frac{\lambda_3}{2}$ est satisfaite pour ce mode.
\end{enumerate}

\textbf{Partie 3 : Synthèse unificatrice}
\begin{enumerate}
    \item Expliquer en quelques lignes rigoureuses le point commun physique fondamental entre le fonctionnement du sonar (détection par écho) et la formation des fuseaux (ondes stationnaires) sur la corde de Melde. Utiliser les termes \emph{réflexion}, \emph{superposition}, \emph{onde incidente} et \emph{onde réfléchie}.
\end{enumerate}

\courssub{Ressources à mobiliser}
\begin{propriete}[Relations utiles]
\begin{itemize}
    \item \textbf{Télémétrie par écho (aller-retour) :} $2d = v \Delta t \quad \Rightarrow \quad d = \dfrac{v \Delta t}{2}$.
    \item \textbf{Relation fondamentale des ondes :} $v = \lambda f$ (avec $\lambda$ en m, $f$ en Hz, $v$ en $\text{m.s}^{-1}$).
    \item \textbf{Célérité des ondes transverses sur une corde tendue :} $v = \sqrt{\dfrac{F}{\mu}}$ ($F$ en N, $\mu$ en $\text{kg.m}^{-1}$).
    \item \textbf{Condition aux limites (corde fixe aux deux bouts) :} Les extrémités sont des \textbf{nœuds}. Longueur $L = n\dfrac{\lambda_n}{2}$ avec $n \in \mathbb{N}^*$.
    \item \textbf{Fréquences propres :} $f_n = n\dfrac{v}{2L} = n f_1$ ($f_1$ : fréquence fondamentale).
    \item \textbf{Position des nœuds (mode $n$) :} $x_k = k\dfrac{\lambda_n}{2} = k\dfrac{L}{n}$ pour $k = 0, 1, \dots, n$.
    \item \textbf{Position des ventres (mode $n$) :} $x_{k+\frac{1}{2}} = \left(k+\frac{1}{2}\right)\dfrac{\lambda_n}{2} = \left(k+\frac{1}{2}\right)\dfrac{L}{n}$ pour $k = 0, 1, \dots, n-1$.
\end{itemize}
\end{propriete}

\courssub{Démarche et corrigé commenté}
\begin{exoresolu}[Corrigé détaillé de l'activité d'intégration]
\tcblower
\textbf{Partie 1 : Télémétrie par sonar au large de Kribi}

\textit{Question 1.1 : Profondeur du banc de poissons.}
L'onde sonore effectue un aller-retour : elle parcourt la distance $2h$ (profondeur $h$ aller + profondeur $h$ retour) en un temps $\Delta t = \SI{0,12}{s}$.
La relation de la télémétrie est :
\[
2h = v_{\text{eau}} \times \Delta t \quad \Rightarrow \quad h = \frac{v_{\text{eau}} \Delta t}{2}
\]
Substitution numérique :
\[
h = \frac{\SI{1500}{m.s^{-1}} \times \SI{0,12}{s}}{2} = \frac{\SI{180}{m}}{2} = \SI{90}{m}
\]
\resultat{La profondeur du banc de poissons est $h = \SI{90}{m}$.}

\textit{Question 1.2 : Longueur d'onde des ultrasons dans l'eau.}
On utilise la relation fondamentale $v = \lambda f$, d'où $\lambda = \dfrac{v}{f}$.
Attention : $f = \SI{50}{kHz} = \SI{50\,000}{Hz} = \num{5,0e4}\,\text{Hz}$.
\[
\lambda_{\text{eau}} = \frac{\SI{1500}{m.s^{-1}}}{\num{5,0e4}\,\text{Hz}} = \frac{1500}{50000}\,\text{m} = \SI{0,03}{m} = \SI{3}{cm}
\]
\textbf{Vérification dimensionnelle :} $[\lambda] = \dfrac{[v]}{[f]} = \dfrac{\text{L.T}^{-1}}{\text{T}^{-1}} = \text{L}$ (longueur en mètre). \textbf{Correct.}
\resultat{La longueur d'onde des ultrasons dans l'eau est $\lambda_{\text{eau}} = \SI{3}{cm}$.}

\bigskip
\textbf{Partie 2 : Modélisation par ondes stationnaires sur une corde (Melde)}

\textit{Question 2.1 : Célérité des ondes sur la corde.}
Formule : $v_{\text{corde}} = \sqrt{\dfrac{F}{\mu}}$.
Conversion de la masse linéique : $\mu = \SI{2}{g.m^{-1}} = \SI{2e-3}{kg.m^{-1}}$.
\[
v_{\text{corde}} = \sqrt{\frac{\SI{50}{N}}{\SI{2e-3}{kg.m^{-1}}}} = \sqrt{\num{25000}}\,\text{m.s}^{-1} = \sqrt{25 \times 10^3}\,\text{m.s}^{-1} = 50\sqrt{10}\,\text{m.s}^{-1}
\]
Valeur approchée : $\sqrt{10} \approx 3,162$, donc $v_{\text{corde}} \approx \SI{158,1}{m.s^{-1}}$.
\resultat{La célérité des ondes transverses sur la corde est $v_{\text{corde}} = 50\sqrt{10} \approx \SI{158}{m.s^{-1}}$.}

\textit{Question 2.2 : Fréquences des trois premiers modes propres.}
Pour une corde fixe aux deux extrémités, les fréquences propres sont $f_n = n \dfrac{v}{2L}$.
Calcul de la fondamentale ($n=1$) :
\[
f_1 = \frac{v_{\text{corde}}}{2L} = \frac{50\sqrt{10}}{2 \times 1,20} = \frac{50\sqrt{10}}{2,4} = \frac{125\sqrt{10}}{6} \approx \SI{65,9}{Hz}
\]
Les harmoniques sont des multiples entiers :
\[
f_2 = 2 f_1 = \frac{125\sqrt{10}}{3} \approx \SI{131,8}{Hz}
\]
\[
f_3 = 3 f_1 = \frac{125\sqrt{10}}{2} \approx \SI{197,6}{Hz}
\]
\resultat{$f_1 \approx \SI{65,9}{Hz}$ ; $f_2 \approx \SI{132}{Hz}$ ; $f_3 \approx \SI{198}{Hz}$.}

\textit{Question 2.3 : Représentation graphique du mode $n=3$ (trois ventres).}
Pour $n=3$, la longueur d'onde est $\lambda_3 = \dfrac{2L}{3} = \dfrac{2 \times 1,20}{3} = \SI{0,80}{m}$.
La distance entre deux nœuds consécutifs est $\dfrac{\lambda_3}{2} = \dfrac{L}{3} = \SI{0,40}{m}$.
Les nœuds (points fixes, amplitude nulle) sont aux abscisses :
\[
x_0 = 0,\quad x_1 = \frac{L}{3} = \SI{0,40}{m},\quad x_2 = \frac{2L}{3} = \SI{0,80}{m},\quad x_3 = L = \SI{1,20}{m}
\]
Les ventres (amplitude maximale) sont au milieu des intervalles entre nœuds :
\[
x_{v1} = \frac{L}{6} = \SI{0,20}{m},\quad x_{v2} = \frac{L}{2} = \SI{0,60}{m},\quad x_{v3} = \frac{5L}{6} = \SI{1,00}{m}
\]

\begin{popfigure}
\begin{tikzpicture}[scale=0.8, every node/.style={font=\small}]
    % Axes : 1 unité TikZ = 1 cm papier. Échelle : 1 cm papier = 0,10 m réel.
    % L_réel = 1,20 m = 120 cm réel -> 12 cm papier.
    % Amplitude arbitraire : 2 cm papier.
    \draw[->, thick] (0,0) -- (13,0) node[right] {$x\,\text{(cm papier)}$};
    \draw[->, thick] (0,-3) -- (0,3) node[above] {$y\,\text{(cm papier)}$};
    % Graduations
    \foreach \x in {0,1,...,12} {
        \draw (\x, -0.1) -- (\x, 0.1) node[below=2pt] {\x};
    }
    \foreach \y in {-2,-1,1,2} {
        \draw (-0.1, \y) -- (0.1, \y) node[left=2pt] {\y};
    }
    % Corde au repos
    \draw[dashed, gray] (0,0) -- (12,0);
    
    % Mode n=3 : y = A sin(3 pi x / L). x en m. x_cm = 10*x_m. L_cm = 12.
    % y = 2 sin(3 pi * x_cm / 12) = 2 sin(pi * x_cm / 4)
    \draw[popblue, very thick, domain=0:12, samples=200, smooth] plot (\x, {2*sin(deg(\x*pi/4))});
    
    % Nœuds (cercles pleins rouges)
    \foreach \x in {0,4,8,12} {
        \fill[popred] (\x,0) circle (2.5pt);
        \node[below=4pt, popred] at (\x,0) {N};
    }
    % Ventres (cercles verts)
    \foreach \x/\y in {2/2, 6/-2, 10/2} {
        \draw[popgreen, thick] (\x,\y) circle (3pt);
        \node[above=4pt, popgreen] at (\x,\y) {V};
    }
    
    % Annotations distances réelles
    \draw[<->, poporange] (0,-2.5) -- (4,-2.5) node[midway, below] {$0,40\,\text{m}$};
    \draw[<->, poporange] (4,-2.5) -- (8,-2.5) node[midway, below] {$0,40\,\text{m}$};
    \draw[<->, poporange] (8,-2.5) -- (12,-2.5) node[midway, below] {$0,40\,\text{m}$};
    
    \node[align=center, popdark] at (6, -3.5) {Mode $n=3$ : $\lambda_3 = 0,80\,\text{m}$, $L = 3\lambda_3/2$};
\end{tikzpicture}
\legende{Allure de la corde pour le mode $n=3$ (échelle : 1 cm sur le graphique = 0,10 m réel). N = Nœud, V = Ventre.}
\end{popfigure}

\textit{Question 2.4 : Vérification de la condition $L = 3\frac{\lambda_3}{2}$.}
Nous avons trouvé $\lambda_3 = \SI{0,80}{m}$.
Calculons $3\dfrac{\lambda_3}{2} = 3 \times \dfrac{0,80}{2} = 3 \times 0,40 = \SI{1,20}{m}$.
Or $L = \SI{1,20}{m}$.
\resultat{L'égalité $L = 3\dfrac{\lambda_3}{2}$ est parfaitement vérifiée.}

\bigskip
\textbf{Partie 3 : Synthèse unificatrice}

\textit{Question 3.1 : Point commun physique.}
Dans les deux situations, le phénomène central est la \textbf{réflexion d'une onde progressive sur un obstacle (ou une discontinuité de milieu)} et la \textbf{superposition} de l'onde incidente et de l'onde réfléchie.

\begin{itemize}
    \item \textbf{Sur la corde de Melde :} L'onde transversale émise par le vibreur se propage jusqu'à la poulie (extrémité fixe), où elle se réfléchit avec \textbf{inversion de phase} (changement de signe de l'amplitude). La superposition de l'onde incidente et de l'onde réfléchie (de même fréquence, même amplitude, sens opposés) crée une \textbf{onde stationnaire}. Les interférences constructives forment les ventres, les interférences destructives forment les nœuds.
    \item \textbf{Dans le sonar de Kribi :} L'impulsion ultrasonore (onde incidente) se propage dans l'eau, rencontre l'interface eau/poisson (changement d'impédance acoustique), et se réfléchit partiellement sous forme d'un \textbf{écho} (onde réfléchie). Le transducteur capte la superposition de l'émission et de la réception. Le temps de vol $\Delta t$ de cet aller-retour permet de déduire la distance (télémétrie).
\end{itemize}
\resultat{\textbf{Point commun :} Dans les deux cas, l'information (position d'un nœud/ventre ou distance d'un obstacle) est codée dans la structure spatiale ou temporelle résultant de la \textbf{superposition d'une onde incidente et de son onde réfléchie}. La réflexion est le processus physique indispensable à la formation de l'onde stationnaire (résonance) comme à la détection par écho (télémétrie).}
\end{exoresolu}

\courssub{Bilan de l'activité}
Cette activité a permis de mettre en évidence plusieurs compétences clés du programme :
\begin{enumerate}
    \item \textbf{Modélisation unifiée :} Une même formalisation mathématique (équation des ondes, superposition, réflexion) décrit des phénomènes physiques a priori distincts : un instrument de musique traditionnel (\emph{mvet}) et un système de navigation moderne (sonar du port de Kribi).
    \item \textbf{Maîtrise des grandeurs et unités :} Le passage des unités usuelles (kHz, g/m, cm) aux unités SI (Hz, kg/m, m) est crucial pour éviter les erreurs d'ordre de grandeur (facteur 1000 sur la fréquence, facteur 1000 sur la masse linéique).
    \item \textbf{Condition de résonance :} La quantification des longueurs d'onde ($\lambda_n = 2L/n$) et des fréquences ($f_n = n v/2L$) sur une corde fixe aux deux bouts est une application directe des conditions aux limites (nœuds aux extrémités).
    \item \textbf{Représentation graphique rigoureuse :} Le tracé du mode $n=3$ impose le respect des positions exactes des nœuds (espacés de $\lambda/2 = L/3$) et des ventres (au milieu), validant la relation $L = n\lambda/2$.
    \item \textbf{Analyse dimensionnelle :} Vérifiée explicitement pour $\lambda = v/f$, elle est un garde-fou indispensable au Bac.
    \item \textbf{Concept de réflexion :} L'activité ancre la compréhension que la \textbf{réflexion} n'est pas une simple "rebond", mais la condition \emph{sine qua non} de l'interférence stationnaire (corde) et de la télémétrie active (sonar, radar, échographie).
\end{enumerate}

\begin{aretenir}[À retenir pour le Bac]
\begin{itemize}
    \item \textbf{Sonar / Télémétrie :} $d = \dfrac{v \Delta t}{2}$ (aller-retour). $\Delta t$ est le temps entre émission et réception de l'écho.
    \item \textbf{Corde de Melde (fixe-fixe) :} $v = \sqrt{F/\mu}$, $\lambda_n = 2L/n$, $f_n = n v/2L$. Nœuds aux extrémités et tous les $\lambda/2$. Ventres au milieu des intervalles.
    \item \textbf{Unificateur :} \textbf{Réflexion + Superposition} $\rightarrow$ Onde stationnaire (résonance) \textbf{OU} Écho (télémétrie).
    \item \textbf{Attention :} $\mu$ en $\text{kg.m}^{-1}$ (pas g/m) ! $f$ en Hz (pas kHz) ! $\Delta t$ aller-retour $\neq$ temps aller simple.
\end{itemize}
\end{aretenir}

\newpage

\coursec{Méthodes et exercices résolus}

\coursec{Superposition des ondes mécaniques : interférences et ondes stationnaires}

\courssub{Principe de superposition et sources cohérentes}

\begin{definition}[Principe de superposition]
Lorsque plusieurs ondes mécaniques se propagent dans un même milieu, l'élongation réelle en un point $M$ et à l'instant $t$ est la somme algébrique (ou vectorielle pour les ondes transversales) des élongations que chaque onde produirait séparément en ce point à cet instant :
\[ y_M(t) = y_1(M,t) + y_2(M,t) + \dots \]
Ce principe est valable tant que le milieu reste linéaire (amplitudes faibles).
\end{definition}

\begin{propriete}[Sources cohérentes]
Deux sources $S_1$ et $S_2$ sont dites \cle{cohérentes} si elles vibrent à la \textbf{même fréquence} $f$ (donc même période $T$ et même longueur d'onde $\lambda$ dans un même milieu) et si leur \textbf{différence de phase initiale reste constante} au cours du temps.
\begin{itemize}
\item \textbf{Condition nécessaire} pour observer une figure d'interférences \emph{stable} dans le temps.
\item \textbf{Réalisation pratique} : les deux sources sont dérivées d'un \emph{même} vibreur (ex. deux pointes sur un même vibreur de cuve à ondes, deux fentes éclairées par une même source lumineuse).
\end{itemize}
\end{propriete}

\begin{experience}[Mise en évidence de la superposition sur une corde]
\textbf{Matériel :} Corde élastique tendue, vibreur à basse fréquence, générateur de fonction basse fréquence (GBF), stroboscope.
\textbf{Protocole :}
\begin{enumerate}
\item Exciter la corde en un point par une impulsion unique (coup sec) : observer la propagation d'un créneau.
\item Exciter la corde par deux impulsions simultanées de part et d'autre : observer le croisement. Pendant le recouvrement, la forme de la corde est la somme des deux créneaux ; après le croisement, chaque créneau reprend sa forme initiale inchangée.
\item Exciter la corde en régime sinusoïdal permanent par le vibreur : observer la formation d'ondes stationnaires (voir §4).
\end{enumerate}
\textbf{Interprétation :} Le croisement sans déformation permanente valide le principe de superposition pour les ondes de faible amplitude sur une corde.
\end{experience}

\courssub{Interférences à la surface de l'eau}

\begin{methode}[Décrire l'expérience des interférences mécaniques]
Pour décrire et schématiser correctement l'expérience (question très fréquente au Bac) :
\begin{enumerate}
\item \textbf{Citer le dispositif} : une cuve à ondes (cuve à eau peu profonde) éclairée par une source lumineuse ; deux pointes $S_1$ et $S_2$ fixées sur un même vibreur frappent la surface de l'eau en $S_1$ et $S_2$.
\item \textbf{Justifier la cohérence} : les deux pointes sont solidaires du \emph{même} vibreur, donc les sources vibrent à la même fréquence $f$ et avec une différence de phase constante (souvent en phase). Ce sont des \cle{sources cohérentes}, condition indispensable pour observer des interférences stables.
\item \textbf{Décrire les observations} : à la surface, on observe un réseau de franges fixes en forme d'\cle{hyperboles} de foyers $S_1$ et $S_2$ : des lignes d'amplitude maximale (vibration forte) alternant avec des lignes d'amplitude quasi nulle (surface immobile).
\item \textbf{Comparer les éclairages} : en \cle{éclairage normal} (continu), on voit le contraste global de la figure ; en \cle{éclairage stroboscopique} de fréquence voisine de $f$, on fige ou ralentit le mouvement et on distingue les crêtes et les creux qui se propagent entre les franges.
\item \textbf{Schématiser} : représenter $S_1$ et $S_2$, des arcs de cercles concentriques (crêtes en trait plein, creux en pointillés) ; les intersections crête--crête ou creux--creux donnent les ventres (amplitude max), les intersections crête--creux donnent les nœuds (repos).
\end{enumerate}
\end{methode}

\begin{propriete}[Conditions d'interférences pour deux sources en phase]
Soit un point $M$ tel que $S_1M = d_1$ et $S_2M = d_2$. La différence de marche est $\delta = |d_2 - d_1|$.
\begin{itemize}
\item \textbf{Interférence constructive (amplitude maximale, ventre)} : $\delta = k\lambda \quad (k \in \mathbb{Z})$.
\item \textbf{Interférence destructive (amplitude nulle, nœud, point de repos)} : $\delta = (2k+1)\dfrac{\lambda}{2} \quad (k \in \mathbb{Z})$.
\end{itemize}
\textbf{Attention :} Si les sources vibrent en \emph{opposition de phase} (déphasage initial $\pi$), les conditions sont inversées.
\end{propriete}

\begin{propriete}[Géométrie des franges]
Le lieu des points $M$ tels que $|S_2M - S_1M| = \text{constante}$ est une \textbf{hyperbole} de foyers $S_1$ et $S_2$. Les franges d'interférences (lignes d'amplitude constante) sont donc des branches d'hyperboles. La médiatrice de $[S_1S_2]$ correspond à $\delta = 0$ (frange centrale, ordre $k=0$).
\end{propriete}

\begin{exoresolu}[Cuve à ondes : description et observation]
\textbf{Énoncé.} Deux pointes $S_1$ et $S_2$, fixées sur un même vibreur de fréquence $f = \SI{20}{Hz}$, affleurent la surface de l'eau d'une cuve à ondes. La célérité des ondes à la surface de l'eau est $c = \SI{0,30}{m.s^{-1}}$.

1. Pourquoi dit-on que $S_1$ et $S_2$ sont des sources cohérentes ? Pourquoi cette condition est-elle nécessaire ?
2. Calculer la longueur d'onde $\lambda$.
3. Décrire ce qu'observe un élève en éclairage normal, puis en éclairage stroboscopique de fréquence $f_s = \SI{20}{Hz}$.

\tcblower
\textbf{Solution détaillée.}

\textbf{1. Cohérence des sources.} Les deux pointes sont solidaires du \emph{même} vibreur : elles sont donc mises en mouvement par la même excitation mécanique. Elles vibrent à la même fréquence $f = \SI{20}{Hz}$ et leur différence de phase reste constante au cours du temps (ici nulle : elles vibrent en phase). Par définition, ce sont des \cle{sources cohérentes}.

Cette condition est nécessaire car le phénomène d'interférences résulte de la superposition en chaque point $M$ de deux ondes dont le déphasage doit être \emph{stable dans le temps} : si les sources étaient indépendantes (fréquences légèrement différentes ou phase aléatoire), le déphasage en $M$ varierait sans cesse, les franges se déplaceraient et se brouilleraient ; aucune figure stable ne serait observable.

\textbf{2. Longueur d'onde.} La relation entre célérité, fréquence et longueur d'onde s'écrit :
\[ \lambda = \frac{c}{f} = \frac{\num{0,30}}{20} = \SI{0,015}{m} = \SI{1,5}{cm}. \]
Vérification dimensionnelle : $[c/f] = \dfrac{\mathrm{m\,s^{-1}}}{\mathrm{s^{-1}}} = \mathrm{m}$. Cohérent.

\textbf{3. Observations.}
\begin{itemize}
\item \emph{En éclairage normal (continu)} : l'œil ne peut suivre la vibration à \SI{20}{Hz} ; on observe une figure stable constituée de franges en forme d'hyperboles de foyers $S_1$ et $S_2$ : des zones où la surface vibre fortement (amplitude maximale) alternent avec des zones où la surface reste pratiquement immobile (amplitude nulle).
\item \emph{En éclairage stroboscopique} avec $f_s = f = \SI{20}{Hz}$ : chaque point de la surface est éclairé toujours à la même phase de son mouvement ; le mouvement paraît \textbf{figé}. On distingue alors nettement les crêtes et les creux immobiles en apparence, et l'on vérifie que les franges de repos correspondent aux intersections crête--creux.
\end{itemize}

\textbf{Conclusion.} Les deux pointes, alimentées par le même vibreur, forment des sources cohérentes de longueur d'onde $\lambda = \SI{1,5}{cm}$ ; la superposition de leurs ondes produit une figure d'interférences stable, figée par le stroboscope réglé à la même fréquence.
\end{exoresolu}

\courssub{Interfrange et exploitation de la figure d'interférences}

\begin{methode}[Exploiter la différence de marche]
Soit un point $M$ de la surface, à distances $d_1 = S_1M$ et $d_2 = S_2M$ des deux sources vibrant \textbf{en phase}. On pose la différence de marche $\delta = |d_2 - d_1|$.
\begin{enumerate}
\item \textbf{Calculer} $\delta$ en mètres (convertir toutes les distances dans la même unité !).
\item \textbf{Calculer} $\lambda = c/f$ ou $\lambda = c\,T$.
\item \textbf{Former le rapport} $k = \delta/\lambda$ et conclure :
\begin{itemize}
\item si $\delta = k\lambda$ avec $k \in \mathbb{Z}$ (c'est-à-dire $\delta/\lambda$ entier) : ondes \textbf{en phase} en $M$ $\Rightarrow$ \cle{amplitude maximale} (frange brillante, ventre) ;
\item si $\delta = (2k+1)\dfrac{\lambda}{2}$ avec $k \in \mathbb{Z}$ (c'est-à-dire $\delta/\lambda$ demi-entier : $0,5$ ; $1,5$ ; $2,5$...) : ondes \textbf{en opposition de phase} en $M$ $\Rightarrow$ \cle{amplitude nulle} (point de repos, nœud).
\end{itemize}
\item \textbf{Géométrie} : le lieu des points de même état vibratoire est une \cle{hyperbole} de foyers $S_1$ et $S_2$. La médiatrice de $[S_1S_2]$ ($\delta = 0$) est une frange d'amplitude maximale si les sources vibrent en phase.
\item \textbf{Attention} : si les sources vibrent en \emph{opposition de phase}, les conditions sont inversées.
\end{enumerate}
\end{methode}

\begin{exoresolu}[État vibratoire d'un point de la surface]
\textbf{Énoncé.} Sur une cuve à ondes, deux sources $S_1$ et $S_2$ vibrent en phase avec la fréquence $f = \SI{25}{Hz}$. La célérité des ondes est $c = \SI{50}{cm.s^{-1}}$. On considère trois points de la surface :
\begin{itemize}
\item point $A$ : $S_1A = \SI{7,0}{cm}$ et $S_2A = \SI{13,0}{cm}$ ;
\item point $B$ : $S_1B = \SI{9,0}{cm}$ et $S_2B = \SI{14,0}{cm}$ ;
\item point $C$ situé sur la médiatrice du segment $[S_1S_2]$.
\end{itemize}
1. Calculer la longueur d'onde $\lambda$.
2. Déterminer, en justifiant, l'état vibratoire de chacun des points $A$, $B$ et $C$.
3. Préciser la nature géométrique des lignes joignant les points de même état vibratoire.

\tcblower
\textbf{Solution détaillée.}

\textbf{1. Longueur d'onde.}
\[ \lambda = \frac{c}{f} = \frac{50}{25} = \SI{2,0}{cm}. \]

\textbf{2. États vibratoires.} Les sources vibrent en phase ; on compare donc la différence de marche $\delta = |d_2 - d_1|$ à $\lambda$.

\emph{Point $A$ :} $\delta_A = |13,0 - 7,0| = \SI{6,0}{cm}$. Rapport :
\[ \frac{\delta_A}{\lambda} = \frac{\num{6,0}}{\num{2,0}} = 3 \quad \text{entier, donc } \delta_A = 3\lambda. \]
Les deux ondes arrivent en $A$ \textbf{en phase} : elles s'ajoutent. Le point $A$ vibre avec une \textbf{amplitude maximale} (frange d'ordre $k = 3$).

\emph{Point $B$ :} $\delta_B = |14,0 - 9,0| = \SI{5,0}{cm}$. Rapport :
\[ \frac{\delta_B}{\lambda} = \frac{\num{5,0}}{\num{2,0}} = 2,5 \quad \text{demi-entier, donc } \delta_B = (2\times 2 + 1)\frac{\lambda}{2} = \frac{5\lambda}{2}. \]
Les deux ondes arrivent en $B$ \textbf{en opposition de phase} : elles s'annulent. Le point $B$ est un \textbf{point de repos} (amplitude nulle).

\emph{Point $C$ :} sur la médiatrice de $[S_1S_2]$, $S_1C = S_2C$ donc $\delta_C = 0 = 0\times\lambda$, cas particulier de $\delta = k\lambda$ avec $k = 0$. Les ondes arrivent en phase : $C$ vibre avec une \textbf{amplitude maximale}. Toute la médiatrice est la frange d'amplitude maximale d'ordre $0$.

\textbf{3. Nature des lignes.} L'ensemble des points $M$ tels que $|S_2M - S_1M| = \text{constante}$ est une \textbf{hyperbole de foyers $S_1$ et $S_2$}. Les franges d'amplitude maximale ($\delta = k\lambda$) et les franges de repos ($\delta = (2k+1)\lambda/2$) forment donc deux familles d'hyperboles alternées, la médiatrice étant l'hyperbole dégénérée correspondant à $k = 0$.

\textbf{Conclusion.} $A$ et $C$ vibrent avec une amplitude maximale ($\delta_A = 3\lambda$, $\delta_C = 0$), tandis que $B$ est au repos ($\delta_B = 5\lambda/2$) ; les franges sont des hyperboles de foyers $S_1$ et $S_2$.
\end{exoresolu}

\begin{methode}[Mesurer $\lambda$ grâce à l'interfrange]
L'\cle{interfrange} $i$ est la distance, le long du segment $[S_1S_2]$, entre deux franges consécutives de même nature (deux maxima ou deux minima).
\begin{enumerate}
\item \textbf{Retenir la formule} : sur l'axe des sources, deux franges de même nature consécutives correspondent à une variation de $\delta$ égale à $\lambda$, donc
\[ i = \frac{\lambda}{2}. \]
\item \textbf{Compter les franges} : si $n$ interfranges occupent une longueur $L$ sur l'axe, alors $L = n\,i = n\,\dfrac{\lambda}{2}$, d'où $\lambda = \dfrac{2L}{n}$.
\item \textbf{Compter correctement} : $n$ est le nombre d'\emph{intervalles} entre franges, pas le nombre de franges (piège classique : 6 franges $\Rightarrow$ 5 interfranges).
\item \textbf{En déduire} la célérité $c = \lambda f$ ou la fréquence $f = c/\lambda$.
\item \textbf{Vérifier} l'homogénéité et donner le résultat avec un nombre raisonnable de chiffres significatifs.
\end{enumerate}
\end{methode}

\begin{exoresolu}[Détermination de la célérité d'une onde à la surface de l'eau]
\textbf{Énoncé.} Deux sources $S_1$ et $S_2$ vibrant en phase à la fréquence $f = \SI{40}{Hz}$ produisent une figure d'interférences à la surface de l'eau. Sur le segment $[S_1S_2]$, on compte 9 franges d'amplitude maximale (médiatrice comprise) réparties sur une longueur $L = \SI{6,4}{cm}$ entre les deux franges extrêmes.

1. Déterminer la valeur de l'interfrange $i$.
2. En déduire la longueur d'onde $\lambda$ puis la célérité $c$ des ondes.
3. On double la fréquence du vibreur. Que devient le nombre de franges d'amplitude maximale sur la même longueur $L$ ?

\tcblower
\textbf{Solution détaillée.}

\textbf{1. Interfrange.} 9 franges délimitent $n = 9 - 1 = 8$ intervalles (interfranges) sur la longueur $L$. Donc :
\[ i = \frac{L}{n} = \frac{\num{6,4}}{8} = \SI{0,80}{cm}. \]

\textbf{2. Longueur d'onde et célérité.} Sur l'axe des sources, l'interfrange vaut la moitié de la longueur d'onde :
\[ i = \frac{\lambda}{2} \quad \Rightarrow \quad \lambda = 2i = \SI{1,6}{cm} = \num{1,6e-2}\,\mathrm{m}. \]
La célérité s'en déduit :
\[ c = \lambda f = \num{1,6e-2} \times 40 = \SI{0,64}{m.s^{-1}}. \]
Vérification dimensionnelle : $[\lambda f] = \mathrm{m} \times \mathrm{s^{-1}} = \mathrm{m\,s^{-1}}$. Correct.

\textbf{3. Doublement de la fréquence.} La célérité $c$ ne dépend que du milieu (eau de la cuve) : elle reste égale à $\SI{0,64}{m.s^{-1}}$. La nouvelle longueur d'onde est :
\[ \lambda' = \frac{c}{f'} = \frac{0,64}{80} = \num{8,0e-3}\,\mathrm{m} = \SI{0,80}{cm}, \qquad i' = \frac{\lambda'}{2} = \SI{0,40}{cm}. \]
Le nombre d'interfranges sur $L$ devient $n' = \dfrac{L}{i'} = \dfrac{\num{6,4}}{\num{0,40}} = 16$, soit $16 + 1 = 17$ franges d'amplitude maximale.

\textbf{Conclusion.} L'interfrange mesurée $i = \SI{0,80}{cm}$ donne $\lambda = \SI{1,6}{cm}$ et $c = \SI{0,64}{m.s^{-1}}$ ; en doublant la fréquence, on double le nombre de franges (17 maxima sur la même longueur), car la célérité, propriété du milieu, est inchangée.
\end{exoresolu}

\courssub{Ondes stationnaires : expérience de Melde}

\begin{definition}[Onde stationnaire]
Une \cle{onde stationnaire} résulte de la superposition de deux ondes progressives de même fréquence, même amplitude, se propageant en sens inverse dans un milieu borné (réflexion sur un obstacle). Contrairement à une onde progressive, le profil d'une onde stationnaire ne se déplace pas : il pulse sur place. L'énergie ne se propage pas, elle est stockée (énergie potentielle élastique et cinétique) entre deux nœuds.
\end{definition}

\begin{propriete}[Nœuds, ventres et distances caractéristiques]
Pour une onde stationnaire unidimensionnelle (corde, colonne d'air) :
\begin{itemize}
\item \textbf{Nœud (N)} : point d'amplitude nulle en permanence (élongation toujours nulle). Pression/force maximale.
\item \textbf{Ventre (V)} : point d'amplitude maximale. Élongation maximale, pression/force nulle.
\item \textbf{Distances} :
  \begin{align*}
  \text{Nœud--Nœud consécutifs} &= \text{Ventre--Ventre consécutifs} = \frac{\lambda}{2}, \\
  \text{Nœud--Ventre voisins} &= \frac{\lambda}{4}.
  \end{align*}
\end{itemize}
\end{propriete}

\begin{methode}[Nœuds, ventres et longueur d'onde (Expérience de Melde)]
Pour une \cle{onde stationnaire} sur une corde (expérience de Melde) :
\begin{enumerate}
\item \textbf{Identifier les points remarquables} : les \cle{nœuds} (points immobiles, amplitude nulle) et les \cle{ventres} (points d'amplitude maximale).
\item \textbf{Appliquer les distances caractéristiques} :
\begin{itemize}
\item distance entre deux nœuds consécutifs (ou deux ventres consécutifs) : $\dfrac{\lambda}{2}$ ;
\item distance entre un nœud et le ventre voisin : $\dfrac{\lambda}{4}$.
\end{itemize}
\item \textbf{Exploiter les conditions aux limites} : une extrémité fixe est toujours un nœud ; une extrémité libre est un ventre. Pour une corde de longueur $L$ fixée aux deux bouts et vibrant en $n$ fuseaux (ventres) :
\[ L = n\,\frac{\lambda}{2} \quad (n \in \mathbb{N}^{*}). \]
\item \textbf{Relier à la célérité} : $v = \lambda f$, avec pour une corde tendue $v = \sqrt{F/\mu}$ ($F$ tension en newtons, $\mu$ masse linéique en $\mathrm{kg\,m^{-1}}$).
\item \textbf{Rédiger proprement} : justifier chaque étape par le comptage des fuseaux sur le schéma ou dans l'énoncé.
\end{enumerate}
\end{methode}

\begin{exoresolu}[Expérience de Melde : corde vibrante en fuseaux]
\textbf{Énoncé.} Une corde de guitare de longueur utile $L = \SI{1,20}{m}$, de masse linéique $\mu = \num{2,0e-3}\,\mathrm{kg.m^{-1}}$, est fixée à ses deux extrémités. Elle est tendue par une force de norme $F = \SI{80}{N}$. Un vibreur l'excite et l'on observe une onde stationnaire formée de $n = 4$ fuseaux.

1. Calculer la célérité $v$ des ondes transversales sur cette corde.
2. Déterminer la longueur d'onde $\lambda$ de l'onde stationnaire observée.
3. En déduire la fréquence $f$ de vibration de la corde.
4. Quelle est la distance entre un nœud et le ventre le plus proche ?

\tcblower
\textbf{Solution détaillée.}

\textbf{1. Célérité des ondes sur la corde.} Pour une corde tendue :
\[ v = \sqrt{\frac{F}{\mu}} = \sqrt{\frac{80}{\num{2,0e-3}}} = \sqrt{\num{4,0e4}} = \num{2,0e2}\,\mathrm{m.s^{-1}} = \SI{200}{m.s^{-1}}. \]
Vérification dimensionnelle : $\left[\sqrt{F/\mu}\right] = \sqrt{\dfrac{\mathrm{kg\,m\,s^{-2}}}{\mathrm{kg\,m^{-1}}}} = \sqrt{\mathrm{m^2\,s^{-2}}} = \mathrm{m\,s^{-1}}$. Correct.

\textbf{2. Longueur d'onde.} Les deux extrémités sont fixes : ce sont des nœuds. La corde forme $n = 4$ fuseaux, chaque fuseau ayant pour longueur $\lambda/2$ :
\[ L = n\,\frac{\lambda}{2} \quad \Rightarrow \quad \lambda = \frac{2L}{n} = \frac{2 \times \num{1,20}}{4} = \SI{0,60}{m}. \]

\textbf{3. Fréquence de vibration.}
\[ f = \frac{v}{\lambda} = \frac{200}{\num{0,60}} \approx \SI{3,3e2}{Hz} = \SI{333}{Hz}. \]

\textbf{4. Distance nœud--ventre.} Un ventre est situé au milieu de chaque fuseau ; la distance entre un nœud et le ventre voisin est donc :
\[ \frac{\lambda}{4} = \frac{\num{0,60}}{4} = \SI{0,15}{m} = \SI{15}{cm}. \]

\textbf{Conclusion.} La corde vibre en 4 fuseaux avec $\lambda = \SI{0,60}{m}$ à la fréquence $f \approx \SI{333}{Hz}$ ; les nœuds sont espacés de $\SI{30}{cm}$ et chaque ventre se trouve à $\SI{15}{cm}$ des nœuds voisins.
\end{exoresolu}

\courssub{Modes propres d'une corde et instruments de musique}

\begin{propriete}[Modes propres d'une corde fixée aux deux extrémités]
Les seules vibrations stationnaires possibles (modes propres) vérifient la condition de quantification :
\[ L = n\,\frac{\lambda_n}{2} \quad \Rightarrow \quad \lambda_n = \frac{2L}{n}, \qquad f_n = \frac{v}{\lambda_n} = n\,\frac{v}{2L} = n\,f_1, \quad n \in \mathbb{N}^{*}. \]
\begin{itemize}
\item $f_1 = \dfrac{v}{2L}$ est le \cle{mode fondamental} (1 fuseau, 2 nœuds aux extrémités).
\item $f_n = n f_1$ sont les \cle{harmoniques} de rang $n$ ($n$ fuseaux, $n+1$ nœuds).
\item La fréquence fondamentale s'exprime en fonction des paramètres physiques : $f_1 = \dfrac{1}{2L}\sqrt{\dfrac{F}{\mu}}$.
\end{itemize}
\end{propriete}

\begin{methode}[Modes propres et harmoniques]
Pour une corde de longueur $L$ fixée à ses deux extrémités (guitare, harpe, \emph{ngoni}...) :
\begin{enumerate}
\item \textbf{Écrire la condition de quantification} : $L = n\,\dfrac{\lambda_n}{2}$, $\lambda_n = \dfrac{2L}{n}$, $f_n = n\,f_1$.
\item \textbf{Identifier} : $f_1$ fondamental ; $f_n$ harmoniques.
\item \textbf{Exploiter les paramètres du musicien} : avec $v = \sqrt{F/\mu}$, on obtient $f_1 = \dfrac{1}{2L}\sqrt{\dfrac{F}{\mu}}$. La fréquence augmente si l'on tend la corde ($F$ augmente), si l'on raccourcit la corde ($L$ diminue, doigt sur le manche) ou si la corde est plus fine ($\mu$ diminue).
\item \textbf{Répondre aux questions inverses} : calculer $L$, $F$ ou $\mu$ connaissant la note (fréquence) émise.
\end{enumerate}
\end{methode}

\begin{savaistu}[Le ngoni et la physique des cordes]
Le \emph{ngoni} (ou \emph{n'goni}), luth traditionnel d'Afrique de l'Ouest très présent au Cameroun (notamment chez les peuples du Nord et dans la musique moderne), illustre parfaitement la physique des cordes. Son manche sans frettes permet au musicien de modifier $L$ en continu pour jouer les notes de la gamme. La tension $F$ est ajustée par les chevilles (accordage). Les cordes, traditionnellement en boyau ou en nylon, ont une masse linéique $\mu$ choisie pour obtenir les bonnes fréquences avec des tensions supportables par l'instrument. Le résonateur (calebasse recouverte de peau) amplifie le son par résonance acoustique (couplage corde--air), un phénomène d'onde stationnaire dans une cavité (voir chapitre sur les colonnes d'air).
\end{savaistu}

\begin{exoresolu}[Accorder une corde de ngoni]
\textbf{Énoncé.} Un musicien camerounais accorde la corde d'un \emph{ngoni} (instrument à cordes pincées). La corde, de longueur utile $L = \SI{60}{cm}$ et de masse linéique $\mu = \num{1,0e-3}\,\mathrm{kg.m^{-1}}$, doit émettre à vide le son de fréquence fondamentale $f_1 = \SI{220}{Hz}$ (note \emph{la}).

1. Calculer la célérité $v$ des ondes sur cette corde.
2. En déduire la tension $F$ qu'il faut appliquer à la corde.
3. Le musicien pince ensuite la corde de façon à faire apparaître le mode propre de rang $n = 3$. Quelle est la fréquence du son émis ? Combien de fuseaux observe-t-on ?
4. Sans modifier la tension, le musicien bloque la corde au tiers de sa longueur. Quelle est la nouvelle fréquence fondamentale ?

\tcblower
\textbf{Solution détaillée.}

\textbf{1. Célérité.} Pour le mode fondamental, la corde forme un seul fuseau : $L = \dfrac{\lambda_1}{2}$, donc $\lambda_1 = 2L = \SI{1,20}{m}$. Alors :
\[ v = \lambda_1 f_1 = \num{1,20} \times 220 = \SI{264}{m.s^{-1}}. \]

\textbf{2. Tension de la corde.} De $v = \sqrt{F/\mu}$, on tire $F = \mu v^2$ :
\[ F = \num{1,0e-3} \times (264)^2 = \num{1,0e-3} \times \num{69696} \approx \SI{70}{N}. \]
Vérification dimensionnelle : $[\mu v^2] = \mathrm{kg\,m^{-1}} \times \mathrm{m^2\,s^{-2}} = \mathrm{kg\,m\,s^{-2}} = \mathrm{N}$. Correct.

\textbf{3. Mode de rang 3.} Les fréquences propres forment une suite harmonique : $f_n = n f_1$. Donc :
\[ f_3 = 3 \times 220 = \SI{660}{Hz}. \]
Le mode de rang $n = 3$ correspond à $L = 3\,\dfrac{\lambda_3}{2}$ : la corde présente \textbf{3 fuseaux} (4 nœuds, extrémités comprises, et 3 ventres).

\textbf{4. Corde raccourcie au tiers.} La nouvelle longueur vibrante est $L' = \dfrac{L}{3} = \SI{0,20}{m}$. La tension (donc $v$) est inchangée :
\[ f'_1 = \frac{v}{2L'} = \frac{264}{2 \times \num{0,20}} = \SI{660}{Hz}. \]
On retrouve $f'_1 = 3 f_1$ : raccourcir la corde au tiers de sa longueur revient, pour le fondamental, à jouer le 3\textsuperscript{e} harmonique de la corde à vide.

\textbf{Conclusion.} Il faut tendre la corde avec $F \approx \SI{70}{N}$ pour obtenir le \emph{la} de $\SI{220}{Hz}$ ; le mode de rang 3 donne un son de $\SI{660}{Hz}$ avec 3 fuseaux, et bloquer la corde au tiers de sa longueur élève le fondamental à cette même fréquence.
\end{exoresolu}

\courssub{Systèmes technologiques à ondes : échographie, télémétrie, sonar, radar}

\begin{propriete}[Principe de la mesure de distance par écho (Télémétrie active)]
Tous ces systèmes (échographie médicale, sonar, radar, télémètre laser) reposent sur le même principe : \textbf{émission d'une salve d'ondes brève, réflexion sur une interface (changement de milieu), réception de l'écho et mesure de la durée $\Delta t$ de l'aller-retour}.
La distance $d$ séparant l'émetteur-récepteur de l'obstacle est donnée par :
\[ 2d = v\,\Delta t \quad \Rightarrow \quad d = \frac{v\,\Delta t}{2} \]
où $v$ est la célérité de l'onde dans le milieu de propagation traversé.
\begin{itemize}
\item \textbf{Facteur 2} : l'onde parcourt deux fois la distance $d$ (aller + retour). C'est l'erreur n°1 à ne pas commettre.
\item \textbf{Ondes utilisées} :
  \begin{itemize}
  \item \textbf{Ultrasons} ($f > \SI{20}{kHz}$) : échographie médicale ($v \approx \SI{1540}{m.s^{-1}}$ dans les tissus mous), sonar (eau de mer $v \approx \SI{1500}{m.s^{-1}}$), télémètres de parking (air $v \approx \SI{340}{m.s^{-1}}$).
  \item \textbf{Ondes électromagnétiques} (lumière, radio) : radar (avions, météo, $c = \num{3,00e8}\,\mathrm{m.s^{-1}}$), télémétrie laser (lidar, topographie, $c = \num{3,00e8}\,\mathrm{m.s^{-1}}$).
  \end{itemize}
\item \textbf{Résolution spatiale} : de l'ordre de la longueur d'onde $\lambda = v/f$. D'où l'intérêt des hautes fréquences (ultrasons MHz en échographie, ondes centimétriques/millimétriques en radar).
\end{itemize}
\end{propriete}

\begin{methode}[Expliquer le fonctionnement d'un système à écho]
\begin{enumerate}
\item \textbf{Identifier l'onde} : nature (mécanique/électromagnétique), fréquence, milieu, célérité $v$.
\item \textbf{Décrire la chaîne} : Émetteur $\rightarrow$ Propagation (aller) $\rightarrow$ Réflexion sur interface $\rightarrow$ Propagation (retour) $\rightarrow$ Récepteur (souvent confondu avec l'émetteur).
\item \textbf{Écrire la relation} : $d = \dfrac{v\,\Delta t}{2}$. Justifier le facteur 2.
\item \textbf{Calculer} $d$ ou $\Delta t$ selon les données.
\item \textbf{Commenter la résolution} : $\Delta d \sim \lambda = v/f$ (plus $f$ est grand, meilleure est la résolution).
\end{enumerate}
\end{methode}

\begin{exoresolu}[Échographie et sonar : mesures de distances par écho]
\textbf{Énoncé.} \emph{On donne : célérité des ultrasons dans les tissus biologiques $v_1 = \SI{1540}{m.s^{-1}}$ ; dans l'eau de mer $v_2 = \SI{1500}{m.s^{-1}}$.}

\textbf{Partie A.} Lors d'une échographie réalisée dans un hôpital de Yaoundé, la sonde émet une salve d'ultrasons et reçoit l'écho réfléchi par un organe au bout de $\Delta t_1 = \SI{80}{\micro\second}$.

1. Expliquer en quelques lignes le principe de l'échographie.
2. Calculer la profondeur $d_1$ de l'organe réfléchissant.

\textbf{Partie B.} Un navire océanographique au large de Kribi mesure la profondeur du plateau continental grâce à un sonar. L'écho du fond marin revient après $\Delta t_2 = \SI{0,40}{s}$.

3. Calculer la profondeur $d_2$ de l'eau sous le navire.
4. Pourquoi utilise-t-on des ultrasons plutôt que des ondes sonores audibles dans ces deux appareils ?

\tcblower
\textbf{Solution détaillée.}

\textbf{1. Principe de l'échographie.} La sonde joue le rôle d'émetteur-récepteur : elle envoie une salve brève d'ultrasons dans le corps. À chaque interface entre deux tissus de nature différente (impédances acoustiques différentes), une partie de l'onde est \emph{réfléchie} et revient vers la sonde. En mesurant la durée $\Delta t$ de l'aller-retour et connaissant la célérité $v_1$ des ultrasons dans les tissus, on calcule la profondeur de l'interface ; l'ensemble des échos, balayés dans le temps, permet de construire une image (échogramme) des organes.

\textbf{2. Profondeur de l'organe.} Pendant $\Delta t_1$, l'onde effectue un aller-retour, soit le trajet $2d_1$ :
\[ 2d_1 = v_1\,\Delta t_1 \quad \Rightarrow \quad d_1 = \frac{v_1\,\Delta t_1}{2} = \frac{1540 \times \num{80e-6}}{2} = \frac{\num{0,1232}}{2} = \SI{6,16e-2}{m}. \]
L'organe se situe à environ $d_1 = \SI{6,2}{cm}$ sous la sonde. Vérification : $[\mathrm{m\,s^{-1}} \times \mathrm{s}] = \mathrm{m}$. Correct.

\textbf{3. Profondeur sous le navire.} Même raisonnement avec l'aller-retour dans l'eau de mer :
\[ d_2 = \frac{v_2\,\Delta t_2}{2} = \frac{1500 \times \num{0,40}}{2} = \SI{300}{m}. \]

\textbf{4. Pourquoi des ultrasons ?} Les ultrasons ($f > \SI{20}{kHz}$, typiquement \SI{2}{MHz} à \SI{18}{MHz} en échographie, \SI{10}{kHz} à \SI{100}{kHz} en sonar) présentent deux avantages majeurs :
\begin{itemize}
\item Ils sont \emph{inaudibles} (pas de gêne pour le patient ou l'équipage).
\item Leur \emph{longueur d'onde est petite} ($\lambda = v/f$ de l'ordre du millimètre en échographie, de la dizaine de centimètres en sonar), ce qui permet d'obtenir des faisceaux directifs (peu de diffraction) et de détecter de petites structures avec une bonne résolution spatiale : la taille des détails discernables est de l'ordre de $\lambda$.
\end{itemize}

\textbf{Conclusion.} L'échographie localise l'organe à $d_1 \approx \SI{6,2}{cm}$ de profondeur et le sonar mesure $d_2 = \SI{300}{m}$ d'eau sous le navire ; dans les deux cas, la distance s'obtient par $d = v\,\Delta t / 2$, le facteur $1/2$ traduisant l'aller-retour de l'écho.
\end{exoresolu}

\begin{attention}[Les 5 erreurs qui coûtent des points au Bac]
\begin{enumerate}
\item \textbf{Oublier le facteur 2 de l'aller-retour} dans les mesures par écho (échographie, sonar, radar, télémétrie laser) : la distance est $d = \dfrac{v\,\Delta t}{2}$, et non $v\,\Delta t$. Écris explicitement « l'onde parcourt deux fois la distance » pour justifier.
\item \textbf{Confondre les conditions d'interférences} : amplitude maximale pour $\delta = k\lambda$ ($k$ entier) et repos pour $\delta = (2k+1)\dfrac{\lambda}{2}$ ($k$ entier) — à condition que les sources vibrent \emph{en phase}. Vérifie toujours que $\delta/\lambda$ est entier ou demi-entier avant de conclure, et précise l'hypothèse sur la phase des sources.
\item \textbf{Mal compter les interfranges ou les fuseaux} : $n$ franges délimitent $n-1$ interfranges ; une corde de longueur $L$ vibrant en $n$ fuseaux vérifie $L = n\,\dfrac{\lambda}{2}$. Un schéma rapide avec les nœuds placés évite l'erreur.
\item \textbf{Confondre les distances caractéristiques de l'onde stationnaire} : nœud--nœud (ou ventre--ventre) $= \dfrac{\lambda}{2}$, mais nœud--ventre $= \dfrac{\lambda}{4}$. Ne jamais écrire que deux nœuds consécutifs sont distants de $\lambda$.
\item \textbf{Négliger les unités et l'homogénéité} : convertir les cm en m, les $\mu$s en s avant tout calcul ; une fréquence en Hz exige une célérité en $\mathrm{m\,s^{-1}}$ et une longueur en m. Termine chaque calcul par le résultat avec son unité et un nombre cohérent de chiffres significatifs (2 ou 3), et vérifie mentalement l'ordre de grandeur (un son audible : entre \SI{20}{Hz} et \SI{20}{kHz} ; une tension de corde : dizaines de newtons).
\end{enumerate}
\end{attention}

\newpage

\coursec{Exercices}

\courssub{Je vérifie mes connaissances}

\begin{exercice}[QCM : Principe de superposition et cohérence]
Deux sources ponctuelles $S_1$ et $S_2$ émettent des ondes sinusoïdales de même fréquence $f$ et de même amplitude à la surface de l'eau. On considère un point $M$ tel que la différence de marche $\Delta = S_2M - S_1M = \frac{3\lambda}{2}$. Les sources sont en phase.
\begin{enumerate}
    \item Quelle est la phase de l'onde résultante en $M$ par rapport à l'onde émise par $S_1$ ?
    \item Quelle est l'amplitude de l'onde résultante en $M$ ?
    \item Si les sources sont en opposition de phase (déphasage initial $\pi$), quelle devient la nature de l'interférence en $M$ ?
\end{enumerate}
\end{exercice}

\begin{exercice}[Vrai / Faux justifié : Interfrange et figure d'interférences]
Dans une cuve à ondes, deux pointes vibrantes $S_1$ et $S_2$ séparées de $a = 4,0\,\text{cm}$ émettent des ondes de longueur d'onde $\lambda = 1,0\,\text{cm}$. On observe la figure d'interférences sur un écran placé à $D = 50\,\text{cm}$ des sources.
\begin{enumerate}
    \item L'interfrange $i$ vaut $12,5\,\text{cm}$.
    \item Les franges d'amplitude maximale sont des droites perpendiculaires à la droite $(S_1S_2)$.
    \item Si on double la fréquence du vibreur (célérité $c$ constante), l'interfrange est divisée par deux.
    \item En éclairage stroboscopique à la fréquence du vibreur, les franges apparaissent immobiles.
\end{enumerate}
\end{exercice}

\begin{exercice}[Définitions et calculs directs : Ondes stationnaires]
Une corde de longueur $L = 1,20\,\text{m}$, fixée à ses deux extrémités, porte une onde stationnaire à trois ventres ($n=3$). La célérité des ondes progressives sur la corde est $c = 60\,\text{m}\cdot\text{s}^{-1}$.
\begin{enumerate}
    \item Définir un nœud et un ventre d'onde stationnaire.
    \item Calculer la longueur d'onde $\lambda$ de l'onde stationnaire.
    \item Déterminer la fréquence $f$ de cette vibration.
    \item Quelle est la distance entre deux nœuds consécutifs ? Et entre un nœud et le ventre adjacent ?
\end{enumerate}
\end{exercice}

\begin{exercice}[Application numérique : Télémétrie par ultrasons]
Un sonar émet une impulsion ultrasonore de fréquence $f = 50\,\text{kHz}$ dans l'eau de mer ($c = 1500\,\text{m}\cdot\text{s}^{-1}$). L'écho revient au bout de $\Delta t = 0,040\,\text{s}$.
\begin{enumerate}
    \item Calculer la longueur d'onde $\lambda$ de l'ultrason dans l'eau.
    \item Déterminer la distance $d$ séparant le sonar de l'obstacle réfléchissant.
    \item Si l'obstacle se déplace vers le sonar à la vitesse $v = 10\,\text{m}\cdot\text{s}^{-1}$, l'écho revient-il plus tôt ou plus tard ? Justifier sans calcul.
\end{enumerate}
\end{exercice}

\courssub{Je m'entraîne}

\begin{exercice}[Interférences à la surface de l'eau : analyse complète]
Deux pointes $S_1$ et $S_2$ d'un vibreur sont plongées dans une cuve à ondes. Elles sont séparées par $a = 6,0\,\text{cm}$ et vibrent en phase à la fréquence $f = 20\,\text{Hz}$. La célérité des ondes à la surface de l'eau est $c = 40\,\text{cm}\cdot\text{s}^{-1}$.
\begin{enumerate}
    \item Calculer la longueur d'onde $\lambda$ des ondes émises.
    \item On place un écran parallèle à $(S_1S_2)$ à une distance $D = 80\,\text{cm}$. Calculer l'interfrange $i$ observée sur l'écran.
    \item On considère un point $M$ de l'écran tel que $S_1M = 82,0\,\text{cm}$ et $S_2M = 83,5\,\text{cm}$. Déterminer la nature de l'interférence en $M$ (renforcement ou affaiblissement). Justifier par le calcul de la différence de marche.
    \item Tracer schématiquement les lignes de vibration maximale (hyperboles) et indiquer le nombre de franges claires visibles entre les deux sources sur l'axe $(S_1S_2)$.
\end{enumerate}
\end{exercice}

\begin{exercice}[Exploitation d'un cliché stroboscopique]
On photographie une figure d'interférences en éclairage stroboscopique (fréquence du flash = fréquence du vibreur). Sur le cliché, on mesure : distance entre les sources $a = 5,0\,\text{cm}$, distance sources-écran $D = 40\,\text{cm}$, interfrange $i = 1,6\,\text{cm}$.
\begin{enumerate}
    \item Déterminer la longueur d'onde $\lambda$ des ondes à la surface de l'eau.
    \item Si la fréquence du vibreur est $f = 25\,\text{Hz}$, calculer la célérité $c$ des ondes.
    \item Préciser ce que l'on observe sur le cliché si l'on double la fréquence du vibreur (et du flash) en gardant la même géométrie.
    \item Que se passe-t-il si l'on utilise un éclairage normal (continu) au lieu du stroboscopique ?
\end{enumerate}
\end{exercice}

\begin{exercice}[Expérience de Melde : détermination des modes propres]
Une corde de longueur $L = 1,50\,\text{m}$, de masse linéique $\mu = 0,50\,\text{g}\cdot\text{m}^{-1}$, est tendue par une force $F = 80\,\text{N}$ entre un vibreur et une poulie porteuse d'une masse. Le vibreur impose une fréquence $f = 200\,\text{Hz}$.
\begin{enumerate}
    \item Calculer la célérité $c$ des ondes transversales sur la corde.
    \item Déterminer la longueur d'onde $\lambda$ correspondante à $f = 200\,\text{Hz}$.
    \item Combien de fuseaux (ventres) observe-t-on sur la corde ? Justifier.
    \item Quelles sont les fréquences du vibreur pour observer exactement 2, 3 et 4 fuseaux ? (On supposera que la longueur vibrante reste $L$).
\end{enumerate}
\end{exercice}

\begin{exercice}[Corde de guitare et intervalles musicaux]
Une corde de guitare, fixée au chevalet et au sillet, a une longueur vibrante $L_0 = 65\,\text{cm}$. Elle émet la note \emph{La} grave de fréquence fondamentale $f_1 = 110\,\text{Hz}$ (corde à vide).
\begin{enumerate}
    \item Calculer la célérité $c$ des ondes sur cette corde.
    \item Le guitariste appuie sur la corde contre une frette pour raccourcir la longueur vibrante à $L$ et obtenir le \emph{La} une octave au-dessus ($f_2 = 220\,\text{Hz}$). Calculer $L$.
    \item Où doit-il placer son doigt (distance par rapport au sillet) pour obtenir le \emph{Mi} aigu ($f_3 = 330\,\text{Hz}$, tierce juste par rapport au \emph{La} grave) ?
    \item Expliquer pourquoi les fréquences propres d'une corde fixée aux deux bouts forment une suite arithmétique de raison $f_1$.
\end{enumerate}
\end{exercice}

\courssub{Je me prépare au Bac}

\begin{exercice}[Bac Type : Interférences sur l'eau — Analyse d'un dispositif expérimental]
\textbf{Énoncé :} Dans le cadre d'une activité pratique, on réalise l'expérience des interférences à la surface de l'eau. Deux pointes $S_1$ et $S_2$ d'un vibreur sont séparées de $a = 5,0\,\text{cm}$. Elles vibrent en phase à la fréquence $f = 25\,\text{Hz}$. La célérité des ondes à la surface de l'eau est $c = 50\,\text{cm}\cdot\text{s}^{-1}$. On observe la figure d'interférences sur un écran placé à $D = 1,00\,\text{m}$ des sources.

\begin{enumerate}
    \item \textbf{Modélisation :} Rappeler la condition sur la différence de marche $\Delta = S_2M - S_1M$ pour obtenir un renforcement (frange claire) en un point $M$ de l'écran. Exprimer l'interfrange $i$ en fonction de $\lambda$, $D$ et $a$.
    \item \textbf{Calculs :} Calculer $\lambda$ et $i$. Donner les résultats en $\text{cm}$.
    \item \textbf{Analyse de points :} On considère trois points $A$, $B$, $C$ sur l'écran tels que :
    \begin{itemize}
        \item $S_1A = 100,0\,\text{cm}$ ; $S_2A = 101,0\,\text{cm}$
        \item $S_1B = 100,0\,\text{cm}$ ; $S_2B = 102,0\,\text{cm}$
        \item $S_1C = 100,0\,\text{cm}$ ; $S_2C = 103,0\,\text{cm}$
    \end{itemize}
    Pour chaque point, calculer la différence de marche et préciser la nature de l'interférence (renforcement ou affaiblissement).
    \item \textbf{Franges entre les sources :} Sur l'axe $(S_1S_2)$, entre les deux sources, combien observe-t-on de franges d'amplitude maximale (ventres) ? Justifier.
    \item \textbf{Éclairage stroboscopique :} Expliquer l'intérêt de l'éclairage stroboscopique (fréquence du flash = $f$) pour photographier cette figure. Que verrait-on sur une photo prise en éclairage normal avec un temps de pose de $0,1\,\text{s}$ ?
\end{enumerate}
\end{exercice}

\begin{exercice}[Bac Type : Ondes stationnaires — Expérience de Melde et spectre fréquentiel]
\textbf{Énoncé :} On étudie les modes propres de vibration d'une corde tendue entre un vibreur et un point fixe. La corde a une longueur $L = 1,20\,\text{m}$ et une masse linéique $\mu = 4,0 \times 10^{-4}\,\text{kg}\cdot\text{m}^{-1}$. Elle est tendue par une force constante $F = 14,4\,\text{N}$.

\begin{enumerate}
    \item \textbf{Célérité :} Calculer la célérité $c$ des ondes transversales propagées sur cette corde.
    \item \textbf{Fréquences propres :} Établir la relation donnant les fréquences propres $f_n$ en fonction du numéro d'ordre $n$ ($n \in \mathbb{N}^*$). Montrer qu'elles forment une suite arithmétique et donner sa raison.
    \item \textbf{Expérience :} On fait varier la fréquence $f$ du vibreur de $50\,\text{Hz}$ à $500\,\text{Hz}$. On note les fréquences pour lesquelles on observe une onde stationnaire stable (résonance) : $f_1 = 75\,\text{Hz}$, $f_2 = 150\,\text{Hz}$, $f_3 = 225\,\text{Hz}$, $f_4 = 300\,\text{Hz}$, $f_5 = 375\,\text{Hz}$.
    \begin{enumerate}
        \item Vérifier que ces valeurs sont compatibles avec la formule établie.
        \item Pour $f = 300\,\text{Hz}$, déterminer le nombre de ventres et tracer l'allure de la corde à l'instant $t=0$ (position d'allongement maximal). Indiquer les positions des nœuds.
    \end{enumerate}
    \item \textbf{Exploitation spectrale :} On enregistre le son émis par la corde vibrante à $f = 300\,\text{Hz}$ et on en analyse le spectre de fréquences. On observe des raies à $300\,\text{Hz}$, $600\,\text{Hz}$, $900\,\text{Hz}$, $1200\,\text{Hz}$.
    \begin{enumerate}
        \item Identifier la fondamentale et les harmoniques présents.
        \item Pourquoi observe-t-on des fréquences qui ne sont pas des fréquences propres de la corde (ex: $600\,\text{Hz}$ n'est pas une fréquence propre pour $n$ entier avec $L=1,20\,\text{m}$) ? Expliquer l'origine physique de ces composantes spectrales.
    \end{enumerate}
\end{enumerate}
\end{exercice}

\begin{exercice}[Bac Type : Systèmes technologiques — Échographie et Radar]
\textbf{Énoncé :} On compare deux systèmes de télémétrie active par réflexion d'ondes : une sonde d'échographie médicale et un radar aérien.

\textbf{Partie A : Échographie médicale}
Une sonde d'échographie émet une impulsion d'ultrasons de fréquence $f_u = 5,0\,\text{MHz}$ dans les tissus mous (célérité $c_t = 1540\,\text{m}\cdot\text{s}^{-1}$). L'interface entre deux organes réfléchit partiellement l'onde. L'écho est reçu par la sonde après un délai $\Delta t = 65\,\mu\text{s}$.
\begin{enumerate}
    \item Calculer la longueur d'onde $\lambda_u$ des ultrasons dans les tissus.
    \item Déterminer la profondeur $p$ de l'interface réfléchissante.
    \item La résolution axiale (capacité à distinguer deux interfaces proches) est de l'ordre de la longueur d'onde. Quelle est la résolution axiale théorique ici ?
    \item Si l'organe se déplace vers la sonde à la vitesse $v = 0,20\,\text{m}\cdot\text{s}^{-1}$, l'écho subit un effet Doppler. Sans calcul, préciser si la fréquence reçue est supérieure ou inférieure à $f_u$.
\end{enumerate}

\textbf{Partie B : Radar aérien}
Un radar au sol émet une impulsion d'onde électromagnétique (célérité $c = 3,00 \times 10^8\,\text{m}\cdot\text{s}^{-1}$) de fréquence $f_r = 10\,\text{GHz}$. Il détecte un avion : l'écho revient après $\Delta t' = 2,0\,\text{ms}$.
\begin{enumerate}
    \item Calculer la distance $d$ de l'avion.
    \item Comparer les longueurs d'onde $\lambda_u$ (échographie) et $\lambda_r$ (radar). Commenter l'influence de la longueur d'onde sur la résolution et la pénétration dans la matière.
    \item L'avion s'éloigne du radar à $v = 250\,\text{m}\cdot\text{s}^{-1}$. Écrire l'expression approchée de la fréquence $f'$ de l'écho reçu (effet Doppler-Fizeau pour $v \ll c$). Calculer le décalage de fréquence $\Delta f = f' - f_r$.
\end{enumerate}

\textbf{Partie C : Synthèse}
\begin{enumerate}
    \item Rappeler le principe commun à ces deux systèmes (émisson, propagation, réflexion, réception).
    \item Pourquoi utilise-t-on des ultrasons (ondes mécaniques) pour l'échographie et des ondes électromagnétiques pour le radar ? Justifier par la nature du milieu de propagation et l'interaction onde-matière.
\end{enumerate}
\end{exercice}

\newpage

\coursec{Corrigés des exercices}

\begin{corrige}[Exercice 1 — QCM : Principe de superposition et cohérence]
\begin{enumerate}
    \item Le déphasage total en $M$ est la somme du déphasage initial des sources et du déphasage dû à la propagation :
    \[
    \Delta\phi = \phi_0 + \frac{2\pi}{\lambda}\Delta = 0 + \frac{2\pi}{\lambda}\times\frac{3\lambda}{2} = 3\pi \equiv \pi \pmod{2\pi}
    \]
    Les deux ondes arrivent en $M$ \textbf{en opposition de phase}. L'amplitude résultante est :
    \[
    A_R = 2A\left|\cos\frac{\Delta\phi}{2}\right| = 2A\left|\cos\frac{3\pi}{2}\right| = 0
    \]
    L'onde résultante est nulle en permanence : sa phase est \textbf{indéterminée} (il n'y a plus d'oscillation en $M$).
    \item L'amplitude résultante est $\boxed{A_R = 0}$ : interférence \textbf{destructive totale}. Le point $M$ reste immobile (point de repos, frange sombre).
    \item Avec des sources en opposition de phase ($\phi_0 = \pi$) :
    \[
    \Delta\phi = \pi + 3\pi = 4\pi \equiv 0 \pmod{2\pi}
    \]
    L'interférence devient \textbf{constructive (renforcement total)} : $A_R = 2A$. Le point $M$ devient un point de vibration maximale.
\end{enumerate}
\textbf{Point de vigilance :} ne jamais oublier le déphasage initial $\phi_0$ des sources dans le déphasage total ; inverser la phase des sources échange les rôles des franges claires et sombres.
\end{corrige}

\begin{corrige}[Exercice 2 — Vrai / Faux justifié : Interfrange et figure d'interférences]
\begin{enumerate}
    \item \textbf{Vrai.} $i = \dfrac{\lambda D}{a} = \dfrac{1{,}0 \times 50}{4{,}0} = 12{,}5\ \text{cm}$.
    \item \textbf{Faux.} Les lignes d'amplitude maximale sont des \textbf{hyperboles} de foyers $S_1$ et $S_2$. Sur un écran parallèle à $(S_1S_2)$ placé loin ($D \gg a$), leurs traces sont localement des droites \textbf{parallèles à $(S_1S_2)$}, et non perpendiculaires.
    \item \textbf{Vrai.} À célérité constante, doubler $f$ divise $\lambda$ par deux ($\lambda = c/f$). Comme $i = \lambda D/a$, l'interfrange est divisée par deux : les franges se resserrent.
    \item \textbf{Vrai.} Avec $f_{\text{flash}} = f$, la surface de l'eau est éclairée toujours au même instant de sa période : la figure d'interférences, stationnaire dans l'espace, apparaît \textbf{immobile et nette}.
\end{enumerate}
\textbf{Point de vigilance :} citer les conditions d'application de $i = \lambda D/a$ ($D \gg a$, écran parallèle à $(S_1S_2)$) dès qu'on utilise cette formule.
\end{corrige}

\begin{corrige}[Exercice 3 — Définitions et calculs directs : Ondes stationnaires]
\begin{enumerate}
    \item \textbf{Nœud :} point de la corde dont l'amplitude de vibration est nulle à tout instant ($y = 0,\ \forall t$). \textbf{Ventre :} point où l'amplitude de vibration est maximale (le point oscille entre $+A$ et $-A$).
    \item Corde fixée aux deux extrémités : $L = n\dfrac{\lambda}{2}$, d'où
    \[
    \lambda = \frac{2L}{n} = \frac{2 \times 1{,}20}{3} = 0{,}80\ \text{m}
    \]
    \item $f = \dfrac{c}{\lambda} = \dfrac{60}{0{,}80} = 75\ \text{Hz}$.
    \item Distance entre deux nœuds consécutifs : $\dfrac{\lambda}{2} = 0{,}40\ \text{m}$. Distance entre un nœud et le ventre adjacent : $\dfrac{\lambda}{4} = 0{,}20\ \text{m}$.
\end{enumerate}
\textbf{Vérification :} 3 ventres sur $1{,}20\ \text{m}$ donnent des fuseaux de $0{,}40\ \text{m}$, cohérent avec $\lambda/2 = 0{,}40\ \text{m}$.
\end{corrige}

\begin{corrige}[Exercice 4 — Application numérique : Télémétrie par ultrasons (Sonar)]
\begin{enumerate}
    \item $\lambda = \dfrac{c}{f} = \dfrac{1500}{50 \times 10^{3}} = 3{,}0 \times 10^{-2}\ \text{m} = 3{,}0\ \text{cm}$.
    \item L'onde effectue un aller-retour : $2d = c\,\Delta t$, donc
    \[
    d = \frac{c\,\Delta t}{2} = \frac{1500 \times 0{,}040}{2} = 30\ \text{m}
    \]
    \item L'obstacle se rapproche du sonar : pendant le trajet de l'impulsion, la distance à parcourir diminue. Le trajet aller-retour est donc plus court et l'écho revient \textbf{plus tôt}.
\end{enumerate}
\textbf{Point de vigilance :} l'erreur classique est d'oublier le facteur 2 (aller-retour) dans $d = c\Delta t/2$.
\end{corrige}

\begin{corrige}[Exercice 5 — Interférences à la surface de l'eau : analyse complète]
\begin{enumerate}
    \item $\lambda = \dfrac{c}{f} = \dfrac{40}{20} = 2{,}0\ \text{cm}$.
    \item $i = \dfrac{\lambda D}{a} = \dfrac{2{,}0 \times 80}{6{,}0} \approx 26{,}7\ \text{cm}$.
    \item Différence de marche : $\Delta = S_2M - S_1M = 83{,}0 - 82{,}0 = 1{,}0\ \text{cm} = \dfrac{\lambda}{2}$.
    Sources en phase : la condition d'affaiblissement est $\Delta = \left(k + \tfrac{1}{2}\right)\lambda$, satisfaite ici avec $k = 0$. L'interférence en $M$ est \textbf{destructive} : $M$ est un point de repos (frange sombre).
    \item Sur l'axe $(S_1S_2)$, entre les sources : $S_1M + S_2M = a = 6{,}0\ \text{cm}$ et, pour un ventre, $\Delta = S_2M - S_1M = k\lambda = 2k\ \text{cm}$.
    \[
    S_2M = \frac{a + k\lambda}{2} = 3{,}0 + k \quad\text{avec}\quad 0 < S_2M < 6{,}0 \iff -3 < k < 3
    \]
    Les valeurs entières admissibles sont $k \in \{-2, -1, 0, 1, 2\}$ : on observe donc \textbf{5 franges claires} (ventres) entre les sources, dont la frange centrale ($k = 0$, médiatrice).
    \begin{popfigure}
    \begin{tikzpicture}[scale=0.9, >=Stealth]
        \def\a{3}
        % Hyperboles schématiques
        \foreach \k in {-2,-1,1,2} {
            \draw[poporange, thick] plot[smooth, tension=0.9] coordinates {
                (-4.5, \k*1.4) (-3, \k*1.1) (0, \k*0.55) (3, \k*1.1) (4.5, \k*1.4)
            };
        }
        % Frange centrale
        \draw[popgreen, thick] (0,-2.6) -- (0,2.6) node[above] {$\Delta = 0$};
        % Sources
        \fill[popblue] (-\a,0) circle (3pt) node[below left] {$S_1$};
        \fill[popblue] (\a,0) circle (3pt) node[below right] {$S_2$};
        % Ventres sur l'axe : x = k*lambda/2 avec lambda/2 = 1 cm -> echelle 1 cm = 1 unite
        \foreach \k in {-2,-1,0,1,2} {
            \fill[poppink] (\k*1.0, 0) circle (2.5pt);
        }
        \draw[<->] (-\a,-0.9) -- (\a,-0.9) node[midway, below] {$a = 6{,}0\ \text{cm}$};
    \end{tikzpicture}
    \legende{Lignes de vibration maximale (hyperboles, orange) et 5 ventres (points roses) sur l'axe $(S_1S_2)$.}
    \end{popfigure}
\end{enumerate}
\textbf{Point de vigilance :} pour compter les franges sur l'axe, la condition stricte est $|k| < a/\lambda$ (ici $a/\lambda = 3$, donc $k = \pm 3$ exclus car ils correspondraient aux sources elles-mêmes).
\end{corrige}

\begin{corrige}[Exercice 6 — Exploitation d'un cliché stroboscopique]
\begin{enumerate}
    \item De $i = \dfrac{\lambda D}{a}$, on tire :
    \[
    \lambda = \frac{i\,a}{D} = \frac{1{,}6 \times 5{,}0}{40} = 0{,}20\ \text{cm} = 2{,}0\ \text{mm}
    \]
    \item $c = \lambda f = 0{,}20 \times 25 = 5{,}0\ \text{cm}\cdot\text{s}^{-1}$.
    \item Doubler $f$ (à $c$ constante) divise $\lambda$ par deux : $\lambda' = 0{,}10\ \text{cm}$, d'où $i' = \dfrac{\lambda' D}{a} = 0{,}80\ \text{cm}$. On observe \textbf{deux fois plus de franges, deux fois plus serrées}. Le flash restant synchronisé, la figure demeure nette et immobile.
    \item En éclairage continu, le temps de pose ($\gg T$) intègre la luminosité sur de nombreuses périodes : crêtes et creux défilent et se compensent. On observe une \textbf{illumination uniforme et floue, sans franges visibles}.
\end{enumerate}
\end{corrige}

\begin{corrige}[Exercice 7 — Expérience de Melde : détermination des modes propres]
\begin{enumerate}
    \item $\mu = 0{,}50\ \text{g}\cdot\text{m}^{-1} = 5{,}0 \times 10^{-4}\ \text{kg}\cdot\text{m}^{-1}$.
    \[
    c = \sqrt{\frac{F}{\mu}} = \sqrt{\frac{80}{5{,}0 \times 10^{-4}}} = \sqrt{1{,}6 \times 10^{5}} = 400\ \text{m}\cdot\text{s}^{-1}
    \]
    \item $\lambda = \dfrac{c}{f} = \dfrac{400}{200} = 2{,}00\ \text{m}$.
    \item Nombre de fuseaux attendu : $n = \dfrac{2L}{\lambda} = \dfrac{2 \times 1{,}50}{2{,}00} = 1{,}5$.
    $n$ n'est \textbf{pas entier} : la condition de résonance $L = n\lambda/2$ n'est pas satisfaite. On n'observe \textbf{pas d'onde stationnaire stable} : la corde vibre en régime forcé, avec une faible amplitude et sans fuseaux nets.
    \item Fréquences propres : $f_n = n\dfrac{c}{2L} = n \times \dfrac{400}{3{,}00} = n \times 133{,}3\ \text{Hz}$.
    \begin{center}
    \begin{tabular}{|c|c|c|c|}\hline
    \rowcolor{popblueL} Nombre de fuseaux $n$ & 2 & 3 & 4 \\ \hline
    Fréquence $f_n$ (Hz) & 266{,}7 & 400{,}0 & 533{,}3 \\ \hline
    \end{tabular}
    \end{center}
\end{enumerate}
\textbf{Point de vigilance :} une onde stationnaire nette (résonance) n'existe que si $n$ est entier ; c'est la condition aux limites (nœuds aux deux extrémités) qui quantifie les fréquences.
\end{corrige}

\begin{corrige}[Exercice 8 — Corde de guitare et intervalles musicaux]
\begin{enumerate}
    \item Corde à vide, mode fondamental : $f_1 = \dfrac{c}{2L_0}$, d'où
    \[
    c = 2L_0 f_1 = 2 \times 0{,}65 \times 110 = 143\ \text{m}\cdot\text{s}^{-1}
    \]
    \item Octave au-dessus : $f_2 = 220\ \text{Hz} = 2f_1$. Comme $f = \dfrac{c}{2L}$ :
    \[
    L = \frac{c}{2f_2} = \frac{143}{440} = 0{,}325\ \text{m} = 32{,}5\ \text{cm} = \frac{L_0}{2}
    \]
    Le doigt est placé au milieu de la corde ($12^{\text{ième}}$ frette).
    \item $f_3 = 330\ \text{Hz} = 3f_1$ :
    \[
    L = \frac{c}{2f_3} = \frac{143}{660} = 0{,}217\ \text{m} \approx 21{,}7\ \text{cm}
    \]
    Le doigt doit être placé à $\approx 21{,}7\ \text{cm}$ du sillet.
    \item Les extrémités fixes imposent des nœuds : $L = n\dfrac{\lambda_n}{2}$, donc $\lambda_n = \dfrac{2L}{n}$ et
    \[
    f_n = \frac{c}{\lambda_n} = n\,\frac{c}{2L} = n f_1
    \]
    Les fréquences propres forment une suite arithmétique de premier terme $f_1$ et de raison $f_1$ : chaque harmonique est un multiple entier de la fondamentale.
\end{enumerate}
\textbf{Point de vigilance :} la célérité $c$ ne dépend que de la tension et de la masse linéique ; raccourcir la corde ne change pas $c$, seulement $L$.
\end{corrige}

\begin{corrige}[Exercice 9 — Bac Type : Interférences sur l'eau]
\textbf{Barème indicatif (10 pts) :} Q1 : 2 pts ; Q2 : 2 pts ; Q3 : 2,5 pts ; Q4 : 2 pts ; Q5 : 1,5 pt.

\begin{enumerate}
    \item Sources en phase ($\phi_0 = 0$) : renforcement (frange claire) si
    \[
    \Delta = S_2M - S_1M = k\lambda \quad (k \in \mathbb{Z})
    \]
    Sous les conditions $D \gg a$ et écran parallèle à $(S_1S_2)$, l'interfrange est :
    \[
    i = \frac{\lambda D}{a}
    \]
    \item $\lambda = \dfrac{c}{f} = \dfrac{50}{25} = 2{,}0\ \text{cm}$ ; \quad $i = \dfrac{2{,}0 \times 100}{5{,}0} = 40\ \text{cm}$.
    \item 
    \begin{itemize}
        \item Point $A$ : $\Delta_A = 101{,}0 - 100{,}0 = 1{,}0\ \text{cm} = \dfrac{\lambda}{2} = \left(0 + \tfrac{1}{2}\right)\lambda$ $\Rightarrow$ \textbf{affaiblissement} (frange sombre).
        \item Point $B$ : $\Delta_B = 102{,}0 - 100{,}0 = 2{,}0\ \text{cm} = \lambda$ ($k = 1$) $\Rightarrow$ \textbf{renforcement} (frange claire).
        \item Point $C$ : $\Delta_C = 103{,}0 - 100{,}0 = 3{,}0\ \text{cm} = \dfrac{3\lambda}{2} = \left(1 + \tfrac{1}{2}\right)\lambda$ $\Rightarrow$ \textbf{affaiblissement}.
    \end{itemize}
    \item Sur l'axe $(S_1S_2)$ entre les sources : $S_1M + S_2M = a = 5{,}0\ \text{cm}$ et, pour un ventre, $\Delta = k\lambda = 2k\ \text{cm}$.
    \[
    S_2M = \frac{a + k\lambda}{2} = 2{,}5 + k \quad\text{avec}\quad 0 < S_2M < 5{,}0 \iff -2{,}5 < k < 2{,}5
    \]
    $k \in \{-2, -1, 0, 1, 2\}$ : \textbf{5 franges claires} entre les sources (dont la frange centrale $k=0$).
    \item \textbf{Stroboscope ($f_{\text{flash}} = f$) :} l'éclairage a lieu toujours à la même phase du mouvement ; la figure d'interférences, stationnaire dans l'espace, apparaît \textbf{figée, nette et contrastée}, ce qui permet la photographie et les mesures.
    \textbf{Éclairage normal (pose de $0{,}1\ \text{s}$) :} $T = \dfrac{1}{f} = 0{,}04\ \text{s}$, donc la pose couvre $2{,}5$ périodes ; l'intensité est moyennée dans le temps : on obtient une \textbf{image grise uniforme, sans franges}.
\end{enumerate}
\textbf{Points de vigilance (Bac) :} écrire explicitement la condition $\Delta = k\lambda$ avec $k \in \mathbb{Z}$ ; justifier l'usage de $i = \lambda D/a$ par $D \gg a$ ; pour le comptage des franges, exclure les positions des sources elles-mêmes (inégalités strictes).
\end{corrige}

\begin{corrige}[Exercice 10 — Bac Type : Ondes stationnaires — Expérience de Melde et spectre fréquentiel]
\textbf{Barème indicatif (10 pts) :} Q1 : 1 pt ; Q2 : 2 pts ; Q3a : 1 pt ; Q3b : 2,5 pts ; Q4a : 1,5 pt ; Q4b : 2 pts.

\begin{enumerate}
    \item $c = \sqrt{\dfrac{F}{\mu}} = \sqrt{\dfrac{12{,}96}{4{,}0 \times 10^{-4}}} = \sqrt{32\,400} = 180\ \text{m}\cdot\text{s}^{-1}$.
    \item Extrémités fixes $\Rightarrow$ nœuds en $x=0$ et $x=L$ : $L = n\dfrac{\lambda_n}{2}$, soit $\lambda_n = \dfrac{2L}{n}$. Alors :
    \[
    f_n = \frac{c}{\lambda_n} = n\,\frac{c}{2L} = n \times \frac{180}{2{,}40} = 75n\ \text{Hz}
    \]
    C'est une \textbf{suite arithmétique} de premier terme $f_1 = 75\ \text{Hz}$ et de raison $r = 75\ \text{Hz}$.
    \item \begin{enumerate}
        \item Valeurs observées : $75, 150, 225, 300, 375\ \text{Hz} = 1, 2, 3, 4, 5 \times 75\ \text{Hz}$ : parfaitement \textbf{compatibles} avec $f_n = 75n$.
        \item $f = 300\ \text{Hz} = 4 \times 75\ \text{Hz} \Rightarrow n = 4$ : \textbf{4 ventres}.
        Nœuds en $x = 0\,;\ 0{,}30\,;\ 0{,}60\,;\ 0{,}90\,;\ 1{,}20\ \text{m}$ ; ventres en $x = 0{,}15\,;\ 0{,}45\,;\ 0{,}75\,;\ 1{,}05\ \text{m}$.
        \begin{popfigure}
        \begin{tikzpicture}[scale=1.2, >=Stealth]
            \def\L{6} \def\A{0.6}
            \draw[popdark] (0,0) -- (\L,0);
            \draw[popblue, very thick, domain=0:\L, samples=150, variable=\x] plot ({\x}, {\A*sin(720*\x/\L)});
            \draw[popblue!50, dashed, domain=0:\L, samples=150, variable=\x] plot ({\x}, {-\A*sin(720*\x/\L)});
            \foreach \k in {0,1,2,3,4} {
                \fill[popred] (\k*1.5,0) circle (2pt) node[below=3pt, font=\small] {N};
            }
            \foreach \k in {0,1,2,3} {
                \node[popblue, font=\small] at (\k*1.5+0.75, \A+0.25) {V};
            }
            \draw[<->] (0,-1.1) -- (\L,-1.1) node[midway, below] {$L = 1{,}20\ \text{m}$};
        \end{tikzpicture}
        \legende{Mode propre $n = 4$ ($f = 300\ \text{Hz}$) : 4 ventres (V) et 5 nœuds (N), extrémités comprises.}
        \end{popfigure}
    \end{enumerate}
    \item \begin{enumerate}
        \item Fondamentale du son : $300\ \text{Hz}$. Harmoniques : $600\ \text{Hz}$ (rang 2), $900\ \text{Hz}$ (rang 3), $1200\ \text{Hz}$ (rang 4).
        \item \textbf{Oui}, ce sont des fréquences propres de la corde : $600 = 8 \times 75$ ($n=8$), $900 = 12 \times 75$ ($n=12$), $1200 = 16 \times 75$ ($n=16$).
        \textbf{Origine physique :} l'excitation réelle n'est jamais parfaitement sinusoïdale : les \textbf{non-linéarités} (variation de tension aux grandes amplitudes, couplage corde-chevalet, imperfections du vibreur) génèrent des harmoniques de $300\ \text{Hz}$. Ces harmoniques coïncidant avec des modes propres de la corde, ils sont amplifiés par résonance et apparaissent dans le spectre du son émis.
    \end{enumerate}
\end{enumerate}
\textbf{Points de vigilance (Bac) :} la démonstration $f_n = n c/(2L)$ à partir des conditions aux limites est exigible ; sur le schéma, les extrémités doivent être des nœuds ; distinguer clairement le rang de l'harmonique du son ($600\ \text{Hz}$ = $2^{\text{ième}}$ harmonique de $300\ \text{Hz}$) et le numéro du mode propre de la corde ($n = 8$).
\end{corrige}

\begin{corrige}[Exercice 11 — Bac Type : Systèmes technologiques — Échographie et Radar]
\textbf{Barème indicatif (10 pts) :} Partie A : 3 pts ; Partie B : 4 pts ; Partie C : 3 pts.

\textbf{Partie A : Échographie médicale}
\begin{enumerate}
    \item $\lambda_u = \dfrac{c_t}{f_u} = \dfrac{1540}{5{,}0 \times 10^{6}} = 3{,}08 \times 10^{-4}\ \text{m} \approx 0{,}31\ \text{mm}$.
    \item Aller-retour : $p = \dfrac{c_t\,\Delta t}{2} = \dfrac{1540 \times 65 \times 10^{-6}}{2} = 5{,}0 \times 10^{-2}\ \text{m} = 5{,}0\ \text{cm}$.
    \item Résolution axiale théorique : $\delta p \sim \lambda_u \approx 0{,}3\ \text{mm}$.
    \item L'organe se rapproche de la sonde : par effet Doppler, la fréquence reçue est \textbf{supérieure} à $f_u$ (écho plus « aigu »).
\end{enumerate}

\textbf{Partie B : Radar aérien}
\begin{enumerate}
    \item $d = \dfrac{c\,\Delta t'}{2} = \dfrac{3{,}00 \times 10^{8} \times 2{,}0 \times 10^{-3}}{2} = 3{,}0 \times 10^{5}\ \text{m} = 300\ \text{km}$.
    \item $\lambda_r = \dfrac{c}{f_r} = \dfrac{3{,}00 \times 10^{8}}{10 \times 10^{9}} = 3{,}0 \times 10^{-2}\ \text{m} = 3{,}0\ \text{cm}$, soit environ $10^{5}$ fois plus grand que $\lambda_u$.
    \textbf{Résolution :} $\delta d \sim \lambda/2$ : l'échographie distingue des détails millimétriques, le radar des objets de taille métrique ou plus.
    \textbf{Pénétration :} les ultrasons se propagent bien dans les milieux aqueux (tissus) mais sont réfléchis par l'air et les os ; les ondes EM traversent l'air et le vide mais sont absorbées par l'eau et réfléchies par les conducteurs.
    \item Avion en éloignement ($v \ll c$) : $\dfrac{\Delta f}{f_r} \approx -\dfrac{2v}{c}$, soit
    \[
    \Delta f = -\frac{2v}{c}\,f_r = -\frac{2 \times 250}{3{,}00 \times 10^{8}} \times 1{,}0 \times 10^{10} \approx -1{,}7 \times 10^{4}\ \text{Hz} = -17\ \text{kHz}
    \]
    L'écho est reçu à $f' \approx f_r - 17\ \text{kHz}$, fréquence \textbf{inférieure} à $f_r$, conforme à l'éloignement.
\end{enumerate}

\textbf{Partie C : Synthèse}
\begin{enumerate}
    \item \textbf{Principe commun (télémétrie active par écho) :} émission d'une impulsion brève ; propagation dans le milieu à la célérité $c$ ; réflexion partielle sur une interface (changement d'impédance acoustique ou diélectrique) ; réception de l'écho ; mesure du temps de vol $\Delta t$ donnant la distance $d = \dfrac{c\,\Delta t}{2}$, et éventuellement mesure Doppler donnant la vitesse radiale.
    \item \textbf{Choix de l'onde :} le corps humain est constitué à environ $70\,\%$ d'eau : les ultrasons (ondes mécaniques) s'y propagent bien avec une atténuation modérée, alors que les ondes EM y sont fortement absorbées. Inversement, l'air et le vide sont transparents aux ondes radio, dont la grande célérité ($3 \times 10^{8}\ \text{m/s}$) autorise des portées de plusieurs centaines de kilomètres, tandis que le son s'atténue rapidement dans l'air. On adapte donc la nature de l'onde au milieu pour optimiser pénétration, portée et résolution.
\end{enumerate}
\textbf{Points de vigilance (Bac) :} ne pas oublier le facteur $1/2$ (aller-retour) dans le calcul des distances ; convertir correctement les unités ($\mu$s, ms, MHz, GHz) ; pour le Doppler, justifier le signe du décalage par le sens du mouvement de la cible.
\end{corrige}

\newpage

\coursec{Fiche bilan}

\begin{popfigure}
\begin{center}
\begin{tikzpicture}[mindmap, grow cyclic,
  every node/.style={concept, font=\small},
  root concept/.append style={concept color=popblue, font=\small\bfseries, text=white},
  level 1/.append style={level distance=3.4cm, sibling angle=51.5},
  level 2/.append style={level distance=2.2cm, font=\scriptsize},
  level 1 concept/.append style={font=\scriptsize\bfseries, text=white}]
\node[root concept]{Superposition des ondes mécaniques}
  child[concept color=poporange]{ node {Principe de superposition}
    child { node[concept color=poporange!70] {$\vec{s}=\vec{s}_1+\vec{s}_2$} }
    child { node[concept color=poporange!70] {Sources cohérentes : même $f$, déphasage constant} }
  }
  child[concept color=popgreen]{ node {Interférences sur l'eau}
    child { node[concept color=popgreen!70] {Max : $\delta=k\lambda$} }
    child { node[concept color=popgreen!70] {Min : $\delta=(2k+1)\dfrac{\lambda}{2}$} }
    child { node[concept color=popgreen!70] {Franges = hyperboles} }
  }
  child[concept color=poppurple]{ node {Interfrange}
    child { node[concept color=poppurple!70] {$i=\dfrac{\lambda D}{a}$} }
    child { node[concept color=poppurple!70] {Stroboscope : fige la figure} }
  }
  child[concept color=popgold]{ node {Onde stationnaire (Melde)}
    child { node[concept color=popgold!70] {Nœuds : $A=0$} }
    child { node[concept color=popgold!70] {Ventres : $A=2a$} }
    child { node[concept color=popgold!70] {Nœud--nœud : $\lambda/2$} }
  }
  child[concept color=popblue!80]{ node {Modes propres}
    child { node[concept color=popblue!50] {$L=n\dfrac{\lambda}{2}$} }
    child { node[concept color=popblue!50] {$f_n=\dfrac{n}{2L}\sqrt{\dfrac{F}{\mu}}$} }
  }
  child[concept color=poppink]{ node {Applications}
    child { node[concept color=poppink!70] {Échographie, sonar} }
    child { node[concept color=poppink!70] {Radar, télémétrie : $d=\dfrac{c\,\Delta t}{2}$} }
  };
\end{tikzpicture}
\end{center}
\legende{Carte mentale de la leçon : les six idées forces à maîtriser.}
\end{popfigure}

\begin{bilan}
\textbf{1. Définitions clés}
\begin{itemize}
  \item \cle{Principe de superposition} : lorsque deux ondes se croisent, l'élongation résultante en un point est la somme algébrique des élongations : $s(M,t)=s_1(M,t)+s_2(M,t)$. Après le croisement, chaque onde poursuit sa propagation sans modification.
  \item \cle{Sources cohérentes} : deux sources de même fréquence $f$ dont le déphasage reste constant au cours du temps. Seules des sources cohérentes donnent une figure d'interférences stable.
  \item \cle{Différence de marche} : $\delta = d_2 - d_1 = S_2M - S_1M$, différence des distances du point $M$ aux deux sources.
  \item \cle{Onde stationnaire} : onde sans propagation apparente, obtenue par superposition de l'onde incidente et de l'onde réfléchie ; elle présente des \cle{nœuds} (amplitude nulle) et des \cle{ventres} (amplitude maximale $2a$).
\end{itemize}

\textbf{2. Formules à connaître par cœur}
\begin{center}
\begin{tabularx}{\linewidth}{|l|X|X|}
\hline
\rowcolor{popblueL} \textbf{Grandeur} & \textbf{Formule} & \textbf{Condition / remarque} \\
\hline
Vibration maximale & $\delta = k\lambda$, $k\in\mathbb{Z}$ & franges d'amplitude $2a$ (hyperboles) \\
\hline
Vibration nulle (repos) & $\delta = (2k+1)\dfrac{\lambda}{2}$ & franges de repos (hyperboles) \\
\hline
Interfrange & $i = \dfrac{\lambda D}{a}$ & $a$ : écart des sources, $D$ : distance sources--écran \\
\hline
Distance nœud--nœud & $d = \dfrac{\lambda}{2}$ & nœud--ventre : $\dfrac{\lambda}{4}$ \\
\hline
Modes propres d'une corde & $L = n\dfrac{\lambda_n}{2}$ \;;\; $f_n = n\,f_1$ & 2 extrémités fixes, $n\in\mathbb{N}^*$ \\
\hline
Fréquence fondamentale & $f_1 = \dfrac{1}{2L}\sqrt{\dfrac{F}{\mu}}$ & $F$ : tension (N), $\mu$ : masse linéique (kg/m) \\
\hline
Télémétrie (aller-retour) & $d = \dfrac{v\,\Delta t}{2}$ & $\Delta t$ : durée de l'aller-retour de l'écho \\
\hline
\end{tabularx}
\end{center}

\textbf{3. Méthodes express}
\begin{itemize}
  \item \emph{Nature d'un point $M$} : calculer $\delta = |S_2M - S_1M|$, puis le quotient $\delta/\lambda$. Entier $\Rightarrow$ maximum ; demi-entier $\Rightarrow$ repos.
  \item \emph{Compter les franges} : $k$ varie entre $-\dfrac{a}{\lambda}$ et $+\dfrac{a}{\lambda}$ car $|\delta|\le S_1S_2=a$.
  \item \emph{Exploiter une figure d'interférences} : mesurer $n$ interfranges sur une longueur $\ell$ : $i=\ell/n$, puis $\lambda = \dfrac{ia}{D}$.
  \item \emph{Onde stationnaire} : compter le nombre $n$ de fuseaux sur la longueur $L$ : $L=n\dfrac{\lambda}{2}$, d'où $\lambda$ puis $v=\lambda f$.
  \item \emph{Échographie / sonar} : ne jamais oublier le facteur $1/2$ (trajet aller-retour) : $d=\dfrac{v\,\Delta t}{2}$.
\end{itemize}
\end{bilan}

\begin{aretenir}[Checklist avant le Bac]
\begin{itemize}
  \item $\square$ Je sais énoncer le principe de superposition et définir des sources cohérentes.
  \item $\square$ Je sais décrire et schématiser l'expérience des interférences à la surface de l'eau (cuve à ondes, deux pointes solidaires du même vibreur).
  \item $\square$ Je connais les conditions sur $\delta$ pour un point de vibration maximale ($\delta=k\lambda$) et un point au repos ($\delta=(2k+1)\lambda/2$).
  \item $\square$ Je sais que les franges d'interférences sont des hyperboles de foyers $S_1$ et $S_2$.
  \item $\square$ Je sais exploiter $i=\dfrac{\lambda D}{a}$ avec les unités SI et vérifier l'homogénéité.
  \item $\square$ Je sais expliquer le rôle du stroboscope (figer ou ralentir l'observation de la figure).
  \item $\square$ Je sais décrire l'expérience de Melde (corde, vibreur, poulie, masses marquées) et reconnaître nœuds et ventres.
  \item $\square$ Je retiens : distance nœud--nœud $=\lambda/2$, nœud--ventre $=\lambda/4$, amplitude au ventre $=2a$.
  \item $\square$ Je sais exploiter la condition aux limites $L=n\lambda/2$ et la relation $f_n=\dfrac{n}{2L}\sqrt{F/\mu}$.
  \item $\square$ Je sais relier les modes propres aux instruments à cordes (guitare, mvet) : corde plus tendue ou plus courte $\Rightarrow$ son plus aigu.
  \item $\square$ Je sais expliquer le principe de l'échographie, du sonar et du radar avec $d=\dfrac{v\,\Delta t}{2}$ (facteur 2 !).
  \item $\square$ Je vérifie toujours mes unités (m, s, Hz) et l'analyse dimensionnelle de chaque formule utilisée.
\end{itemize}
\end{aretenir}

\findecours

\end{document}
